%% -------------------------------------------------------------------------
%% GENEROVANO: make_all.py -- needitovat rucne.
%%
%% Popisky, tabulky a poznamky bez preambule. Do clanku staci jeden radek:
%%
%%     %% -------------------------------------------------------------------------
%% GENEROVANO: make_all.py -- needitovat rucne.
%%
%% Popisky, tabulky a poznamky bez preambule. Do clanku staci jeden radek:
%%
%%     %% -------------------------------------------------------------------------
%% GENEROVANO: make_all.py -- needitovat rucne.
%%
%% Popisky, tabulky a poznamky bez preambule. Do clanku staci jeden radek:
%%
%%     %% -------------------------------------------------------------------------
%% GENEROVANO: make_all.py -- needitovat rucne.
%%
%% Popisky, tabulky a poznamky bez preambule. Do clanku staci jeden radek:
%%
%%     \input{tables_tom/output/supplement_body}
%%
%% Tabulky se vejdou i do uzke sazby `elsarticle` (390 pt), ale zdrobni.
%% Aby vypadaly jako v samostatnem supplementu, obal se vlozeni geometrii:
%%
%%     \usepackage{geometry}          % v preambuli, za elsarticle
%%     ...
%%     \clearpage
%%     \newgeometry{margin=15mm}
%%     \input{tables_tom/output/supplement_body}
%%     \restoregeometry
%%
%% Cislovani S1--S7 misto 1--7 zaridi pred vlozenim:
%%     \setcounter{table}{0}\renewcommand{\thetable}{S\arabic{table}}
%%
%% Preambule potrebuje: booktabs, array, amsmath, amssymb.
%% Cesta k tabulkam se bere z \tablesdir; vychozi hodnota plati pro preklad
%% z korene projektu.
%% -------------------------------------------------------------------------
\providecommand{\tablesdir}{tables_tom/output}

%% Zaloha, kdyby `manuscriptnotes` nebyl nactveny: vysadi se jen kotva,
%% poznamka se zahodi -- tedy tentyz vysledek jako v rezimu odesilani.
\providecommand{\NOTE}[2][]{#1}
\providecommand{\NOTEWARN}[2][]{#1}
\providecommand{\NOTECHECK}[2][]{#1}
\providecommand{\NOTETODO}[2][]{#1}
\providecommand{\NOTEQUERY}[2][]{#1}

{\large\bfseries Supplementary tables\par}
\bigskip
\NOTEWARN{the captions and notes in this supplement were drafted with AI assistance and are provisional; every one of them has to be rewritten and checked by an author before submission}

\begin{table}[htbp]
\centering
\caption{Patient selection in the two cohorts. MR CLEAN NO-IV was used for model development and cross-validation, the MR CLEAN Registry as the independent external cohort. Entries are numbers of patients; an em dash marks an exclusion step that does not apply to that cohort. eTICI, expanded Thrombolysis in Cerebral Infarction score; the outcome modeled throughout this article is near-complete or complete reperfusion, defined as a final eTICI score of 2C or 3. Slice spacing refers to the through-plane spacing of the non-contrast CT acquisition.}
\label{tab:cohorts}
\smallskip
\input{\tablesdir/table1_cohorts}
\end{table}

\smallskip
{\footnotesize
\sloppy
\noindent
\textit{Notes to Table~\ref{tab:cohorts}.}
The manual review of segmentation quality was carried out on the MR CLEAN NO-IV cohort only.
No segmentation was reviewed or rejected in the MR CLEAN Registry cohort, so the two cohorts differ in segmentation quality control as well as in population.
This is a property of the design and is the single largest asymmetry between the development and the external cohort.
The review was blinded to the eTICI score: it judged the quality and the anatomical placement of the segmentation only, so the 155 exclusions it produced cannot have been informed by the outcome.
Of the 141 patients that remain, 82 (58.2\%) reached eTICI 2C--3; the proportion varies across the five outer folds (0.448, 0.714, 0.536, 0.679, 0.536), so in the first fold the majority class is the absence rather than the presence of near-complete or complete reperfusion.
\NOTETODO{breakdown of how many of the 155 reviewed cases were rejected by which criterion}
\NOTETODO{number and proportion of patients with eTICI 2C--3 reperfusion in the Registry cohort, which the table leaves as an em dash}
\NOTEWARN{size of the Registry cohort: the exclusion cascade gives 1906, the stored identifier list holds 1839, and the denominators reconstructed from the decimal expansion of the reported external accuracies are 1649 (8 models), 1760 (6 models), 1761 (2 models). Neither 1906 nor 1839 is compatible with any of the sixteen transferred models, and the reconstructed values are mutually incompatible, so the external cohort is not the same across feature sets}
\par}

\clearpage

\begin{table}[htbp]
\centering
\caption{Hyperparameter search spaces of the four pipelines. LASSO, $\ell_1$-penalized linear regression on the binary outcome used as a feature selector; RFI, ranking by the gain-based feature importance of a random-forest gradient-boosting model; XGBRF, \texttt{XGBRFClassifier}; LinearSVC, linear support vector classifier with an $\ell_1$ penalty, which acts as selector and classifier simultaneously; NuSVC, $\nu$-support vector classifier with a radial basis function kernel. In the Component column, (S) marks the component acting as feature selector, (C) as classifier, and (S+C) a single object that performs both roles. Scale indicates how Optuna sampled the parameter: \emph{log}, log-uniform over the interval; \emph{int}, uniform over the integers of the interval; \emph{float}, uniform over the interval.}
\label{tab:hyperparameters}
\smallskip
\input{\tablesdir/table2_hyperparameters}
\end{table}

\smallskip
{\footnotesize
\sloppy
\noindent
\textit{Notes to Table~\ref{tab:hyperparameters}.}
$n$ and \texttt{X.shape[1]} denote the number of rows and the number of columns of the training matrix passed to the search.
Rows are patients for the global-based and distribution-based datasets and voxels for the voxel-based datasets, so the same expression yields different intervals across datasets; the number of columns reaches 20660 for the distribution-based datasets.
The minority fraction is the proportion of the smaller outcome class among the training labels of the given fold.
In the RFI--RF pipeline a single parameter dictionary configures both the model whose gain importances rank the features and the classifier fitted on the retained ones, which is why its rows are marked (S+C) rather than (C).
All search spaces were explored with Optuna using the Gaussian-process sampler (\texttt{GPSampler}) with 20 trials per dataset-by-pipeline combination.
The objective maximized was the \emph{mean accuracy over the five inner cross-validation folds}, not AUC; the inner folds were generated by \texttt{StratifiedGroupKFold} grouped by \texttt{patient\_id}.
Settings that were fixed and not searched, one per line so that the long identifiers have somewhere to break:\\
\texttt{Lasso(fit\_intercept=True, max\_iter=5000)}\\
\texttt{LinearSVC(penalty='l1', max\_iter=5000, class\_weight='balanced')}\\
\texttt{NuSVC(kernel='rbf', probability=True, class\_weight='balanced')}\\
\texttt{XGBRFClassifier(random\_state=42)}\\
Class weighting was therefore applied in the two support vector branches only.
The models were seeded with \texttt{random\_state=42}, but \texttt{GPSampler()} was instantiated without a seed, so the sequence of 20 trials is not reproducible.
\par}

\clearpage

\begin{table}[htbp]
\centering
\caption{Cross-validated performance of all 24 dataset-by-pipeline combinations in the MR CLEAN NO-IV cohort ($n = 141$), over five outer folds. Feature set: \emph{global}, eight summary statistics of each radiomic feature over the thrombus, one row per patient; \emph{distribution}, a ten-bin histogram of each radiomic feature over the thrombus, one row per patient; \emph{voxel}, the radiomic feature values themselves, one row per voxel. ROI: \emph{whole}, the dilated thrombus mask; \emph{centerline}, the dilated skeleton of the thrombus mask. Selector (Sel.): LASSO, $\ell_1$-penalized regression; L1, embedded selection by the non-zero coefficients of the classifier itself; RFI, gain-based feature importance of an \texttt{XGBRFClassifier}. Classifier (Class.): XGBRF, \texttt{XGBRFClassifier}; LinSVC, linear support vector classifier with an $\ell_1$ penalty; NuSVC, $\nu$-support vector classifier with a radial basis function kernel. In the body of the table the distribution-based feature set is abbreviated \emph{distrib.}\ and the centerline ROI \emph{cent.} AUC, area under the receiver operating characteristic curve; ACC, accuracy; IQR, interquartile range across the five folds; min--max, the smallest and largest value across the five folds; Med.\ $n$, median number of features retained by the selector; Null AUC, mean AUC of the permutation baseline; $\kappa$, the $\kappa$-type index of the observed AUC relative to that baseline; $p$, permutation $p$-value; $q$, the same $p$-value after Benjamini--Hochberg correction across all 24 combinations.}
\label{tab:cv-performance}
\smallskip
\input{\tablesdir/table3_cv_performance}
\end{table}

\smallskip
{\footnotesize
\sloppy
\noindent
\textit{Notes to Table~\ref{tab:cv-performance}.}
Accuracy must be read against the majority-class accuracy of the development cohort, $82/141 = 0.582$; only 4 of the 24 combinations exceed it by median accuracy, and 86 of the 120 individual fold results (71.7\%) fall below it.
The class ratio itself varies across folds (0.448, 0.714, 0.536, 0.679, 0.536), so in the first fold the majority class is the absence rather than the presence of near-complete or complete reperfusion.
The largest single fold result of the whole study is 0.810, obtained by distribution/whole/RFI/NuSVC in fold 4.
Averaged over the eight whole-thrombus patient-level combinations, AUC per fold is 0.510, 0.661, 0.606, 0.587, 0.623, a spread within a single fold of up to 0.531; between-fold variability therefore exceeds between-method variability.
The permutation null was obtained from 100 label permutations at the patient level per outer fold; with $m = 100$ permutations the smallest attainable $p$-value is $p = (0+1)/(100+1) = 0.0099$, which bounds the resolution of every permutation result in this table.
The Null AUC, $\kappa$, $p$ and $q$ columns were computed from the \emph{mean} AUC across folds, whereas the AUC column reports the \emph{median}; the two summaries are not interchangeable and are shown side by side deliberately.
10 of the 24 combinations reach $p < 0.05$ uncorrected and 6 survive Benjamini--Hochberg correction\NOTETODO{cite Benjamini and Hochberg 1995; the key is missing from BibFile.bib} for AUC, while none survives it for accuracy.
No centerline combination survives correction; across the eight of them the mean AUC ranges from 0.396 to 0.495 against a null of 0.489--0.510, $\kappa$ from $-0.206$ to $-0.014$ and $p$ from 0.5347 to 1.0000.\\
\textsuperscript{a} In these 2 combinations the $\ell_1$ penalty shrank every coefficient to zero in every fold and the code substituted the index of the first column; the retained feature is a fallback, not one the selector identified as informative.
Combinations with a median of one retained feature but no mark reached that median with genuinely non-zero coefficients.\\
\textsuperscript{b} The unit of analysis of the voxel-based datasets is the voxel, not the patient, and their metrics are therefore \emph{not} comparable with the remaining rows.
Across the five folds the voxel/whole datasets contribute 396085 voxel-level predictions (91500, 79327, 75748, 82369, 67141) and the voxel/centerline datasets 8557 (1690, 1571, 1811, 1806, 1679), against 141 patient-level predictions (29, 28, 28, 28, 28) for every other dataset.
The consequence is visible in the permutation null, whose standard deviation is 0.040--0.052 for the voxel-based combinations against 0.063--0.120 for all others, so that the voxel-based $p$-values are computed against a spuriously narrow null and are not comparable with the rest of the table.
3 voxel-based combinations are formally significant for AUC after correction; they must not be read as evidence for a signal.
\NOTETODO{patient-level majority-vote aggregation of the voxel-based predictions, which would make these rows comparable}
\par}

\clearpage

\begin{table}[htbp]
\centering
\caption{External validation in the MR CLEAN Registry cohort. Each pipeline was refitted on all 141 development patients with its tuned hyperparameters and applied once to the Registry cohort. CV, cross-validated median over the five outer folds of the same combination in Table~\ref{tab:cv-performance}; Reg., Registry; ACC, accuracy; $\Delta$AUC, Registry AUC minus cross-validated median AUC, a negative value being a loss of discrimination on transfer; other abbreviations as in Table~\ref{tab:cv-performance}. The voxel-based datasets were not transferred externally and therefore have no row here, which is why the table has 16 rows rather than 24.}
\label{tab:external-validation}
\smallskip
\input{\tablesdir/table4_external_validation}
\end{table}

\smallskip
{\footnotesize
\sloppy
\noindent
\textit{Notes to Table~\ref{tab:external-validation}.}
Across the 16 models the Registry AUC ranged from 0.471 to 0.521 with a median of 0.496, and 5 of them exceeded 0.50.
Registry accuracy ranged from 0.440 to 0.550 with a median of 0.481, and no model reached the majority-class accuracy of the development cohort, 0.582.
The sign of $\Delta$ follows the cross-validated value and nothing else: the 8 combinations whose cross-validated median AUC was above 0.500 all moved downwards on transfer, by $-0.034$ to $-0.190$, and the 8 that were at or below it all moved upwards, by $0.002$ to $0.128$.
The upward movement is therefore a regression to chance of models that had performed below chance internally, not a gain in generalization, and it must not be presented as one.
The only combinations for which the comparison is informative are the 6 that discriminated in cross-validation, that is reached a cross-validated median AUC of at least 0.600; all of them lose between 0.150 and 0.190 of AUC.
\NOTEWARN{$n$ Registry is not stated in the results files and is not the same for every model. Reconstructing the smallest denominator compatible with the decimal expansion of each reported accuracy leaves exactly three cohort sizes: 1649 for eight models, 1760 for six and 1761 for two. Neither the 1906 of the exclusion cascade nor the 1839 identifiers stored with the cohort is compatible with any of the sixteen. Until this is resolved the AUC values in this table are not strictly comparable with one another}
\par}

\clearpage

\begin{table}[htbp]
\centering
\caption{Comparison with the reference pipeline of van Voorst et al. \cite{vanVoorst2022} The first row summarizes the present study; rows two to four were computed on the MR CLEAN NO-IV development cohort with the same five outer folds as in Table~\ref{tab:cv-performance}, so that the comparison is made within one dataset rather than across publications; the last row is the value reported in the original publication, on a different cohort and for a different endpoint. ICC, intraclass correlation coefficient: of the 107 features the reference pipeline extracts, the 52 with ICC $\geq$ 0.75 are retained and the rest discarded. AUC$_{\text{prob}}$ denotes the AUC computed from the predicted class probabilities, and is the quantity reported for the reference pipeline. Values are medians with the interquartile range in brackets, or means $\pm$ standard deviation across folds, as indicated in the note below. An em dash marks a quantity the row does not have: the pooled row spans combinations with different accuracies, and the published row reports neither a range nor an accuracy.}
\label{tab:reference-comparison}
\smallskip
\input{\tablesdir/table5_reference_comparison}
\end{table}

\smallskip
{\footnotesize
\sloppy
\noindent
\textit{Notes to Table~\ref{tab:reference-comparison}.}
The first row pools 25 fold results, 5 combinations $\times$ 5 outer folds, and its bracketed values are the interquartile range; the second row is a mean $\pm$ standard deviation over the 5 splits; rows three and four are medians over the 5 folds, with no further feature selection applied to the 52 ICC-filtered features.
The last row is quoted from the publication: thrombus radiomics for eTICI $\geq$ 2B, a more permissive endpoint than the eTICI 2C--3 outcome of this article, as a mean over 100 bootstrap samples of 100 patients.
The plain \texttt{AUC} column of the reference pipeline degenerates to exactly 0.50 in 4 of the five splits, because the fitted model assigns the majority class to every patient in those splits and its hard predictions therefore carry no ranking information.
AUC$_{\text{prob}}$, computed from the predicted probabilities, is the only quantity of that pipeline that is comparable with the AUC values reported elsewhere in this article, and it is the one used here.
By median accuracy the reference pipeline and both of our classifiers restricted to its 52 ICC-filtered features stay below the majority-class accuracy of 0.582, although the reference pipeline exceeds it in one individual split.
Rows two to four were recomputed against the eTICI 2C--3 outcome of this article --- the stored label column is named \texttt{TICI\_2B-3}, but its per-fold positive counts are exactly the 82/141 of this study --- and are directly comparable with the first row, whereas the published 0.52 of the last row was obtained for eTICI $\geq$ 2B in a different cohort and differs from them in endpoint and cohort as much as in pipeline.
\NOTEQUERY{the pool of the first row omits one whole-thrombus combination that also discriminated in cross-validation, and the reason is not documented in the source material}
\NOTEQUERY{decide which of the two mutually incompatible sets of numbers in the source data applies before this table is finalized: set A, 0.481 $\pm$ 0.145 on 141 patients against set B, 0.467 $\pm$ 0.099 on 170}
\par}

\clearpage

\begin{table}[htbp]
\centering
\caption{Ensemble analysis: the 15 configurations of discrete, that is binary, AdaBoost \cite{Freund1997} applied to the stored prediction vectors of the 16 patient-level weak classifiers (four patient-level datasets $\times$ four pipelines). Each configuration was fitted on the five outer folds, giving 75 ensemble models. ACC $\varnothing$, mean accuracy of the configuration across the five folds; best weak classifier, the single pipeline in the configuration's pool with the highest mean accuracy, and its value; $\Delta$, the mean accuracy of the ensemble minus that of the best weak classifier. A negative $\Delta$ means the ensemble is worse than the best model it combines. The configuration name states which of the four patient-level datasets form its pool, each contributing all four pipelines: \texttt{type1} pools all four datasets (16 weak classifiers), \texttt{type2} a whole-thrombus dataset with its own centerline counterpart (8), \texttt{type3} a single dataset (4), \texttt{type4} two datasets that are not such a pair (8) and \texttt{type5} all datasets but one (12). \texttt{distrib.}\ abbreviates the distribution-based dataset and \emph{cent.} the centerline region of interest, as in Table~\ref{tab:cv-performance}. The two voxel-based datasets were excluded from every pool because their predictions are defined per voxel.}
\label{tab:ensemble}
\smallskip
\input{\tablesdir/table6_ensemble}
\end{table}

\smallskip
{\footnotesize
\sloppy
\noindent
\textit{Notes to Table~\ref{tab:ensemble}.}
$\Delta$ is negative for all 15 configurations, between $-0.022$ and $-0.121$; no ensemble outperformed the best single pipeline in its own pool.
The worst result belongs to \texttt{type1\_\allowbreak all\_\allowbreak featuresets}, the configuration that pools all feature sets.
Across all 75 fitted ensemble models accuracy averages $0.543 \pm 0.083$ with a range of 0.414--0.714.
The mechanism is documented by the fitted weights: in the all-feature-set configuration, 55 of the 80 boosting rounds (16 rounds $\times$ five folds) received a negative weight $\alpha$, meaning that the weak classifier admitted in that round performed worse than chance on the reweighted data.
The individual pipelines therefore do not make complementary errors, and the boosting assumption of weak classifiers reliably better than chance is not met.
Labels were coded as $\{-1,+1\}$, the weight of a round was $\alpha = \tfrac{1}{2}\ln\frac{1-\varepsilon}{\varepsilon}$ with $\varepsilon$ the error of the admitted classifier on the currently reweighted data, and each weak classifier could be admitted at most once, so a pool of $k$ members yields at most $k$ rounds.
The weak classifiers were not retrained: the procedure is applied to models that had already been fitted, which makes it methodologically a form of stacking rather than boosting proper.
The best weak classifier of \texttt{type3\_\allowbreak global} and of \texttt{type1\_\allowbreak all\_\allowbreak featuresets} is the same pipeline, because the pool of the latter contains the pool of the former.
\par}

\clearpage

\begin{table}[htbp]
\centering
\caption{Pairwise comparison of the pipelines by the corrected resampled $t$-test of Nadeau and Bengio \cite{Nadeau2003}. The upper block compares the four pipelines of this study with one another on the whole-thrombus global-based dataset; the lower block compares them, and our two classifiers restricted to the reference pipeline's own 52 ICC-filtered features, against the reference pipeline of van Voorst et al.\ recomputed on the same cohort. Every model in the table was evaluated on the same 5 outer folds of the same 141 patients, which is what makes the paired test applicable. Because the training sets of $k$-fold cross-validation overlap, the per-fold differences are not independent and the ordinary paired $t$-test underestimates their variance; the correction multiplies the variance by $1/k + n_{\text{test}}/n_{\text{train}}$, here $1/5 + 28.2/112.8 = 0.4500$, which inflates the standard error 1.50-fold at $\mathrm{df} = 4$. Fold sizes are taken from the data, not assumed. Mean $\Delta$AUC is the mean of the 5 per-fold differences in AUC between pipeline A and pipeline B; a positive value favours A. $t$ (uncorrected) is the ordinary paired statistic, $t$ (corrected) the same quantity with the inflated variance. The two-sided critical value at $\alpha = 0.05$ and $\mathrm{df} = 4$ is 2.776. Pipeline names combine selector and classifier as in Table~\ref{tab:cv-performance}.}
\label{tab:nadeau-bengio}
\smallskip
\input{\tablesdir/table7_nadeau_bengio}
\end{table}

\smallskip
{\footnotesize
\sloppy
\noindent
\textit{Notes to Table~\ref{tab:nadeau-bengio}.}
Not one of the 12 comparisons reaches significance: the largest corrected statistic in absolute value is 1.54, against a critical value of 2.776.
The four pipelines of this study are therefore statistically indistinguishable from one another, and none of them is distinguishable from the reference pipeline either, in either direction.
For the reference pipeline the compared quantity is AUC$_{\text{prob}}$, because its plain \texttt{AUC} column degenerates to exactly 0.50 in four of the five splits; see Table~\ref{tab:reference-comparison}.
The pooled row of Table~\ref{tab:reference-comparison} is deliberately absent here: those 5 combinations were picked by their results on these same folds, so testing them against the reference pipeline would be biased in favour of the present study. All four pipelines of the dataset are tested instead, whatever their outcome.
With 5 folds the test has $\mathrm{df} = 4$ and very little power, so these results are an absence of evidence for a difference, not evidence that the pipelines perform identically. The comparison that the data would support best, a DeLong test on the external ROC curves, was not possible because per-patient predicted scores on the Registry were not retained.
\par}

\clearpage


%%
%% Tabulky se vejdou i do uzke sazby `elsarticle` (390 pt), ale zdrobni.
%% Aby vypadaly jako v samostatnem supplementu, obal se vlozeni geometrii:
%%
%%     \usepackage{geometry}          % v preambuli, za elsarticle
%%     ...
%%     \clearpage
%%     \newgeometry{margin=15mm}
%%     %% -------------------------------------------------------------------------
%% GENEROVANO: make_all.py -- needitovat rucne.
%%
%% Popisky, tabulky a poznamky bez preambule. Do clanku staci jeden radek:
%%
%%     \input{tables_tom/output/supplement_body}
%%
%% Tabulky se vejdou i do uzke sazby `elsarticle` (390 pt), ale zdrobni.
%% Aby vypadaly jako v samostatnem supplementu, obal se vlozeni geometrii:
%%
%%     \usepackage{geometry}          % v preambuli, za elsarticle
%%     ...
%%     \clearpage
%%     \newgeometry{margin=15mm}
%%     \input{tables_tom/output/supplement_body}
%%     \restoregeometry
%%
%% Cislovani S1--S7 misto 1--7 zaridi pred vlozenim:
%%     \setcounter{table}{0}\renewcommand{\thetable}{S\arabic{table}}
%%
%% Preambule potrebuje: booktabs, array, amsmath, amssymb.
%% Cesta k tabulkam se bere z \tablesdir; vychozi hodnota plati pro preklad
%% z korene projektu.
%% -------------------------------------------------------------------------
\providecommand{\tablesdir}{tables_tom/output}

%% Zaloha, kdyby `manuscriptnotes` nebyl nactveny: vysadi se jen kotva,
%% poznamka se zahodi -- tedy tentyz vysledek jako v rezimu odesilani.
\providecommand{\NOTE}[2][]{#1}
\providecommand{\NOTEWARN}[2][]{#1}
\providecommand{\NOTECHECK}[2][]{#1}
\providecommand{\NOTETODO}[2][]{#1}
\providecommand{\NOTEQUERY}[2][]{#1}

{\large\bfseries Supplementary tables\par}
\bigskip
\NOTEWARN{the captions and notes in this supplement were drafted with AI assistance and are provisional; every one of them has to be rewritten and checked by an author before submission}

\begin{table}[htbp]
\centering
\caption{Patient selection in the two cohorts. MR CLEAN NO-IV was used for model development and cross-validation, the MR CLEAN Registry as the independent external cohort. Entries are numbers of patients; an em dash marks an exclusion step that does not apply to that cohort. eTICI, expanded Thrombolysis in Cerebral Infarction score; the outcome modeled throughout this article is near-complete or complete reperfusion, defined as a final eTICI score of 2C or 3. Slice spacing refers to the through-plane spacing of the non-contrast CT acquisition.}
\label{tab:cohorts}
\smallskip
\input{\tablesdir/table1_cohorts}
\end{table}

\smallskip
{\footnotesize
\sloppy
\noindent
\textit{Notes to Table~\ref{tab:cohorts}.}
The manual review of segmentation quality was carried out on the MR CLEAN NO-IV cohort only.
No segmentation was reviewed or rejected in the MR CLEAN Registry cohort, so the two cohorts differ in segmentation quality control as well as in population.
This is a property of the design and is the single largest asymmetry between the development and the external cohort.
The review was blinded to the eTICI score: it judged the quality and the anatomical placement of the segmentation only, so the 155 exclusions it produced cannot have been informed by the outcome.
Of the 141 patients that remain, 82 (58.2\%) reached eTICI 2C--3; the proportion varies across the five outer folds (0.448, 0.714, 0.536, 0.679, 0.536), so in the first fold the majority class is the absence rather than the presence of near-complete or complete reperfusion.
\NOTETODO{breakdown of how many of the 155 reviewed cases were rejected by which criterion}
\NOTETODO{number and proportion of patients with eTICI 2C--3 reperfusion in the Registry cohort, which the table leaves as an em dash}
\NOTEWARN{size of the Registry cohort: the exclusion cascade gives 1906, the stored identifier list holds 1839, and the denominators reconstructed from the decimal expansion of the reported external accuracies are 1649 (8 models), 1760 (6 models), 1761 (2 models). Neither 1906 nor 1839 is compatible with any of the sixteen transferred models, and the reconstructed values are mutually incompatible, so the external cohort is not the same across feature sets}
\par}

\clearpage

\begin{table}[htbp]
\centering
\caption{Hyperparameter search spaces of the four pipelines. LASSO, $\ell_1$-penalized linear regression on the binary outcome used as a feature selector; RFI, ranking by the gain-based feature importance of a random-forest gradient-boosting model; XGBRF, \texttt{XGBRFClassifier}; LinearSVC, linear support vector classifier with an $\ell_1$ penalty, which acts as selector and classifier simultaneously; NuSVC, $\nu$-support vector classifier with a radial basis function kernel. In the Component column, (S) marks the component acting as feature selector, (C) as classifier, and (S+C) a single object that performs both roles. Scale indicates how Optuna sampled the parameter: \emph{log}, log-uniform over the interval; \emph{int}, uniform over the integers of the interval; \emph{float}, uniform over the interval.}
\label{tab:hyperparameters}
\smallskip
\input{\tablesdir/table2_hyperparameters}
\end{table}

\smallskip
{\footnotesize
\sloppy
\noindent
\textit{Notes to Table~\ref{tab:hyperparameters}.}
$n$ and \texttt{X.shape[1]} denote the number of rows and the number of columns of the training matrix passed to the search.
Rows are patients for the global-based and distribution-based datasets and voxels for the voxel-based datasets, so the same expression yields different intervals across datasets; the number of columns reaches 20660 for the distribution-based datasets.
The minority fraction is the proportion of the smaller outcome class among the training labels of the given fold.
In the RFI--RF pipeline a single parameter dictionary configures both the model whose gain importances rank the features and the classifier fitted on the retained ones, which is why its rows are marked (S+C) rather than (C).
All search spaces were explored with Optuna using the Gaussian-process sampler (\texttt{GPSampler}) with 20 trials per dataset-by-pipeline combination.
The objective maximized was the \emph{mean accuracy over the five inner cross-validation folds}, not AUC; the inner folds were generated by \texttt{StratifiedGroupKFold} grouped by \texttt{patient\_id}.
Settings that were fixed and not searched, one per line so that the long identifiers have somewhere to break:\\
\texttt{Lasso(fit\_intercept=True, max\_iter=5000)}\\
\texttt{LinearSVC(penalty='l1', max\_iter=5000, class\_weight='balanced')}\\
\texttt{NuSVC(kernel='rbf', probability=True, class\_weight='balanced')}\\
\texttt{XGBRFClassifier(random\_state=42)}\\
Class weighting was therefore applied in the two support vector branches only.
The models were seeded with \texttt{random\_state=42}, but \texttt{GPSampler()} was instantiated without a seed, so the sequence of 20 trials is not reproducible.
\par}

\clearpage

\begin{table}[htbp]
\centering
\caption{Cross-validated performance of all 24 dataset-by-pipeline combinations in the MR CLEAN NO-IV cohort ($n = 141$), over five outer folds. Feature set: \emph{global}, eight summary statistics of each radiomic feature over the thrombus, one row per patient; \emph{distribution}, a ten-bin histogram of each radiomic feature over the thrombus, one row per patient; \emph{voxel}, the radiomic feature values themselves, one row per voxel. ROI: \emph{whole}, the dilated thrombus mask; \emph{centerline}, the dilated skeleton of the thrombus mask. Selector (Sel.): LASSO, $\ell_1$-penalized regression; L1, embedded selection by the non-zero coefficients of the classifier itself; RFI, gain-based feature importance of an \texttt{XGBRFClassifier}. Classifier (Class.): XGBRF, \texttt{XGBRFClassifier}; LinSVC, linear support vector classifier with an $\ell_1$ penalty; NuSVC, $\nu$-support vector classifier with a radial basis function kernel. In the body of the table the distribution-based feature set is abbreviated \emph{distrib.}\ and the centerline ROI \emph{cent.} AUC, area under the receiver operating characteristic curve; ACC, accuracy; IQR, interquartile range across the five folds; min--max, the smallest and largest value across the five folds; Med.\ $n$, median number of features retained by the selector; Null AUC, mean AUC of the permutation baseline; $\kappa$, the $\kappa$-type index of the observed AUC relative to that baseline; $p$, permutation $p$-value; $q$, the same $p$-value after Benjamini--Hochberg correction across all 24 combinations.}
\label{tab:cv-performance}
\smallskip
\input{\tablesdir/table3_cv_performance}
\end{table}

\smallskip
{\footnotesize
\sloppy
\noindent
\textit{Notes to Table~\ref{tab:cv-performance}.}
Accuracy must be read against the majority-class accuracy of the development cohort, $82/141 = 0.582$; only 4 of the 24 combinations exceed it by median accuracy, and 86 of the 120 individual fold results (71.7\%) fall below it.
The class ratio itself varies across folds (0.448, 0.714, 0.536, 0.679, 0.536), so in the first fold the majority class is the absence rather than the presence of near-complete or complete reperfusion.
The largest single fold result of the whole study is 0.810, obtained by distribution/whole/RFI/NuSVC in fold 4.
Averaged over the eight whole-thrombus patient-level combinations, AUC per fold is 0.510, 0.661, 0.606, 0.587, 0.623, a spread within a single fold of up to 0.531; between-fold variability therefore exceeds between-method variability.
The permutation null was obtained from 100 label permutations at the patient level per outer fold; with $m = 100$ permutations the smallest attainable $p$-value is $p = (0+1)/(100+1) = 0.0099$, which bounds the resolution of every permutation result in this table.
The Null AUC, $\kappa$, $p$ and $q$ columns were computed from the \emph{mean} AUC across folds, whereas the AUC column reports the \emph{median}; the two summaries are not interchangeable and are shown side by side deliberately.
10 of the 24 combinations reach $p < 0.05$ uncorrected and 6 survive Benjamini--Hochberg correction\NOTETODO{cite Benjamini and Hochberg 1995; the key is missing from BibFile.bib} for AUC, while none survives it for accuracy.
No centerline combination survives correction; across the eight of them the mean AUC ranges from 0.396 to 0.495 against a null of 0.489--0.510, $\kappa$ from $-0.206$ to $-0.014$ and $p$ from 0.5347 to 1.0000.\\
\textsuperscript{a} In these 2 combinations the $\ell_1$ penalty shrank every coefficient to zero in every fold and the code substituted the index of the first column; the retained feature is a fallback, not one the selector identified as informative.
Combinations with a median of one retained feature but no mark reached that median with genuinely non-zero coefficients.\\
\textsuperscript{b} The unit of analysis of the voxel-based datasets is the voxel, not the patient, and their metrics are therefore \emph{not} comparable with the remaining rows.
Across the five folds the voxel/whole datasets contribute 396085 voxel-level predictions (91500, 79327, 75748, 82369, 67141) and the voxel/centerline datasets 8557 (1690, 1571, 1811, 1806, 1679), against 141 patient-level predictions (29, 28, 28, 28, 28) for every other dataset.
The consequence is visible in the permutation null, whose standard deviation is 0.040--0.052 for the voxel-based combinations against 0.063--0.120 for all others, so that the voxel-based $p$-values are computed against a spuriously narrow null and are not comparable with the rest of the table.
3 voxel-based combinations are formally significant for AUC after correction; they must not be read as evidence for a signal.
\NOTETODO{patient-level majority-vote aggregation of the voxel-based predictions, which would make these rows comparable}
\par}

\clearpage

\begin{table}[htbp]
\centering
\caption{External validation in the MR CLEAN Registry cohort. Each pipeline was refitted on all 141 development patients with its tuned hyperparameters and applied once to the Registry cohort. CV, cross-validated median over the five outer folds of the same combination in Table~\ref{tab:cv-performance}; Reg., Registry; ACC, accuracy; $\Delta$AUC, Registry AUC minus cross-validated median AUC, a negative value being a loss of discrimination on transfer; other abbreviations as in Table~\ref{tab:cv-performance}. The voxel-based datasets were not transferred externally and therefore have no row here, which is why the table has 16 rows rather than 24.}
\label{tab:external-validation}
\smallskip
\input{\tablesdir/table4_external_validation}
\end{table}

\smallskip
{\footnotesize
\sloppy
\noindent
\textit{Notes to Table~\ref{tab:external-validation}.}
Across the 16 models the Registry AUC ranged from 0.471 to 0.521 with a median of 0.496, and 5 of them exceeded 0.50.
Registry accuracy ranged from 0.440 to 0.550 with a median of 0.481, and no model reached the majority-class accuracy of the development cohort, 0.582.
The sign of $\Delta$ follows the cross-validated value and nothing else: the 8 combinations whose cross-validated median AUC was above 0.500 all moved downwards on transfer, by $-0.034$ to $-0.190$, and the 8 that were at or below it all moved upwards, by $0.002$ to $0.128$.
The upward movement is therefore a regression to chance of models that had performed below chance internally, not a gain in generalization, and it must not be presented as one.
The only combinations for which the comparison is informative are the 6 that discriminated in cross-validation, that is reached a cross-validated median AUC of at least 0.600; all of them lose between 0.150 and 0.190 of AUC.
\NOTEWARN{$n$ Registry is not stated in the results files and is not the same for every model. Reconstructing the smallest denominator compatible with the decimal expansion of each reported accuracy leaves exactly three cohort sizes: 1649 for eight models, 1760 for six and 1761 for two. Neither the 1906 of the exclusion cascade nor the 1839 identifiers stored with the cohort is compatible with any of the sixteen. Until this is resolved the AUC values in this table are not strictly comparable with one another}
\par}

\clearpage

\begin{table}[htbp]
\centering
\caption{Comparison with the reference pipeline of van Voorst et al. \cite{vanVoorst2022} The first row summarizes the present study; rows two to four were computed on the MR CLEAN NO-IV development cohort with the same five outer folds as in Table~\ref{tab:cv-performance}, so that the comparison is made within one dataset rather than across publications; the last row is the value reported in the original publication, on a different cohort and for a different endpoint. ICC, intraclass correlation coefficient: of the 107 features the reference pipeline extracts, the 52 with ICC $\geq$ 0.75 are retained and the rest discarded. AUC$_{\text{prob}}$ denotes the AUC computed from the predicted class probabilities, and is the quantity reported for the reference pipeline. Values are medians with the interquartile range in brackets, or means $\pm$ standard deviation across folds, as indicated in the note below. An em dash marks a quantity the row does not have: the pooled row spans combinations with different accuracies, and the published row reports neither a range nor an accuracy.}
\label{tab:reference-comparison}
\smallskip
\input{\tablesdir/table5_reference_comparison}
\end{table}

\smallskip
{\footnotesize
\sloppy
\noindent
\textit{Notes to Table~\ref{tab:reference-comparison}.}
The first row pools 25 fold results, 5 combinations $\times$ 5 outer folds, and its bracketed values are the interquartile range; the second row is a mean $\pm$ standard deviation over the 5 splits; rows three and four are medians over the 5 folds, with no further feature selection applied to the 52 ICC-filtered features.
The last row is quoted from the publication: thrombus radiomics for eTICI $\geq$ 2B, a more permissive endpoint than the eTICI 2C--3 outcome of this article, as a mean over 100 bootstrap samples of 100 patients.
The plain \texttt{AUC} column of the reference pipeline degenerates to exactly 0.50 in 4 of the five splits, because the fitted model assigns the majority class to every patient in those splits and its hard predictions therefore carry no ranking information.
AUC$_{\text{prob}}$, computed from the predicted probabilities, is the only quantity of that pipeline that is comparable with the AUC values reported elsewhere in this article, and it is the one used here.
By median accuracy the reference pipeline and both of our classifiers restricted to its 52 ICC-filtered features stay below the majority-class accuracy of 0.582, although the reference pipeline exceeds it in one individual split.
Rows two to four were recomputed against the eTICI 2C--3 outcome of this article --- the stored label column is named \texttt{TICI\_2B-3}, but its per-fold positive counts are exactly the 82/141 of this study --- and are directly comparable with the first row, whereas the published 0.52 of the last row was obtained for eTICI $\geq$ 2B in a different cohort and differs from them in endpoint and cohort as much as in pipeline.
\NOTEQUERY{the pool of the first row omits one whole-thrombus combination that also discriminated in cross-validation, and the reason is not documented in the source material}
\NOTEQUERY{decide which of the two mutually incompatible sets of numbers in the source data applies before this table is finalized: set A, 0.481 $\pm$ 0.145 on 141 patients against set B, 0.467 $\pm$ 0.099 on 170}
\par}

\clearpage

\begin{table}[htbp]
\centering
\caption{Ensemble analysis: the 15 configurations of discrete, that is binary, AdaBoost \cite{Freund1997} applied to the stored prediction vectors of the 16 patient-level weak classifiers (four patient-level datasets $\times$ four pipelines). Each configuration was fitted on the five outer folds, giving 75 ensemble models. ACC $\varnothing$, mean accuracy of the configuration across the five folds; best weak classifier, the single pipeline in the configuration's pool with the highest mean accuracy, and its value; $\Delta$, the mean accuracy of the ensemble minus that of the best weak classifier. A negative $\Delta$ means the ensemble is worse than the best model it combines. The configuration name states which of the four patient-level datasets form its pool, each contributing all four pipelines: \texttt{type1} pools all four datasets (16 weak classifiers), \texttt{type2} a whole-thrombus dataset with its own centerline counterpart (8), \texttt{type3} a single dataset (4), \texttt{type4} two datasets that are not such a pair (8) and \texttt{type5} all datasets but one (12). \texttt{distrib.}\ abbreviates the distribution-based dataset and \emph{cent.} the centerline region of interest, as in Table~\ref{tab:cv-performance}. The two voxel-based datasets were excluded from every pool because their predictions are defined per voxel.}
\label{tab:ensemble}
\smallskip
\input{\tablesdir/table6_ensemble}
\end{table}

\smallskip
{\footnotesize
\sloppy
\noindent
\textit{Notes to Table~\ref{tab:ensemble}.}
$\Delta$ is negative for all 15 configurations, between $-0.022$ and $-0.121$; no ensemble outperformed the best single pipeline in its own pool.
The worst result belongs to \texttt{type1\_\allowbreak all\_\allowbreak featuresets}, the configuration that pools all feature sets.
Across all 75 fitted ensemble models accuracy averages $0.543 \pm 0.083$ with a range of 0.414--0.714.
The mechanism is documented by the fitted weights: in the all-feature-set configuration, 55 of the 80 boosting rounds (16 rounds $\times$ five folds) received a negative weight $\alpha$, meaning that the weak classifier admitted in that round performed worse than chance on the reweighted data.
The individual pipelines therefore do not make complementary errors, and the boosting assumption of weak classifiers reliably better than chance is not met.
Labels were coded as $\{-1,+1\}$, the weight of a round was $\alpha = \tfrac{1}{2}\ln\frac{1-\varepsilon}{\varepsilon}$ with $\varepsilon$ the error of the admitted classifier on the currently reweighted data, and each weak classifier could be admitted at most once, so a pool of $k$ members yields at most $k$ rounds.
The weak classifiers were not retrained: the procedure is applied to models that had already been fitted, which makes it methodologically a form of stacking rather than boosting proper.
The best weak classifier of \texttt{type3\_\allowbreak global} and of \texttt{type1\_\allowbreak all\_\allowbreak featuresets} is the same pipeline, because the pool of the latter contains the pool of the former.
\par}

\clearpage

\begin{table}[htbp]
\centering
\caption{Pairwise comparison of the pipelines by the corrected resampled $t$-test of Nadeau and Bengio \cite{Nadeau2003}. The upper block compares the four pipelines of this study with one another on the whole-thrombus global-based dataset; the lower block compares them, and our two classifiers restricted to the reference pipeline's own 52 ICC-filtered features, against the reference pipeline of van Voorst et al.\ recomputed on the same cohort. Every model in the table was evaluated on the same 5 outer folds of the same 141 patients, which is what makes the paired test applicable. Because the training sets of $k$-fold cross-validation overlap, the per-fold differences are not independent and the ordinary paired $t$-test underestimates their variance; the correction multiplies the variance by $1/k + n_{\text{test}}/n_{\text{train}}$, here $1/5 + 28.2/112.8 = 0.4500$, which inflates the standard error 1.50-fold at $\mathrm{df} = 4$. Fold sizes are taken from the data, not assumed. Mean $\Delta$AUC is the mean of the 5 per-fold differences in AUC between pipeline A and pipeline B; a positive value favours A. $t$ (uncorrected) is the ordinary paired statistic, $t$ (corrected) the same quantity with the inflated variance. The two-sided critical value at $\alpha = 0.05$ and $\mathrm{df} = 4$ is 2.776. Pipeline names combine selector and classifier as in Table~\ref{tab:cv-performance}.}
\label{tab:nadeau-bengio}
\smallskip
\input{\tablesdir/table7_nadeau_bengio}
\end{table}

\smallskip
{\footnotesize
\sloppy
\noindent
\textit{Notes to Table~\ref{tab:nadeau-bengio}.}
Not one of the 12 comparisons reaches significance: the largest corrected statistic in absolute value is 1.54, against a critical value of 2.776.
The four pipelines of this study are therefore statistically indistinguishable from one another, and none of them is distinguishable from the reference pipeline either, in either direction.
For the reference pipeline the compared quantity is AUC$_{\text{prob}}$, because its plain \texttt{AUC} column degenerates to exactly 0.50 in four of the five splits; see Table~\ref{tab:reference-comparison}.
The pooled row of Table~\ref{tab:reference-comparison} is deliberately absent here: those 5 combinations were picked by their results on these same folds, so testing them against the reference pipeline would be biased in favour of the present study. All four pipelines of the dataset are tested instead, whatever their outcome.
With 5 folds the test has $\mathrm{df} = 4$ and very little power, so these results are an absence of evidence for a difference, not evidence that the pipelines perform identically. The comparison that the data would support best, a DeLong test on the external ROC curves, was not possible because per-patient predicted scores on the Registry were not retained.
\par}

\clearpage


%%     \restoregeometry
%%
%% Cislovani S1--S7 misto 1--7 zaridi pred vlozenim:
%%     \setcounter{table}{0}\renewcommand{\thetable}{S\arabic{table}}
%%
%% Preambule potrebuje: booktabs, array, amsmath, amssymb.
%% Cesta k tabulkam se bere z \tablesdir; vychozi hodnota plati pro preklad
%% z korene projektu.
%% -------------------------------------------------------------------------
\providecommand{\tablesdir}{tables_tom/output}

%% Zaloha, kdyby `manuscriptnotes` nebyl nactveny: vysadi se jen kotva,
%% poznamka se zahodi -- tedy tentyz vysledek jako v rezimu odesilani.
\providecommand{\NOTE}[2][]{#1}
\providecommand{\NOTEWARN}[2][]{#1}
\providecommand{\NOTECHECK}[2][]{#1}
\providecommand{\NOTETODO}[2][]{#1}
\providecommand{\NOTEQUERY}[2][]{#1}

{\large\bfseries Supplementary tables\par}
\bigskip
\NOTEWARN{the captions and notes in this supplement were drafted with AI assistance and are provisional; every one of them has to be rewritten and checked by an author before submission}

\begin{table}[htbp]
\centering
\caption{Patient selection in the two cohorts. MR CLEAN NO-IV was used for model development and cross-validation, the MR CLEAN Registry as the independent external cohort. Entries are numbers of patients; an em dash marks an exclusion step that does not apply to that cohort. eTICI, expanded Thrombolysis in Cerebral Infarction score; the outcome modeled throughout this article is near-complete or complete reperfusion, defined as a final eTICI score of 2C or 3. Slice spacing refers to the through-plane spacing of the non-contrast CT acquisition.}
\label{tab:cohorts}
\smallskip
\input{\tablesdir/table1_cohorts}
\end{table}

\smallskip
{\footnotesize
\sloppy
\noindent
\textit{Notes to Table~\ref{tab:cohorts}.}
The manual review of segmentation quality was carried out on the MR CLEAN NO-IV cohort only.
No segmentation was reviewed or rejected in the MR CLEAN Registry cohort, so the two cohorts differ in segmentation quality control as well as in population.
This is a property of the design and is the single largest asymmetry between the development and the external cohort.
The review was blinded to the eTICI score: it judged the quality and the anatomical placement of the segmentation only, so the 155 exclusions it produced cannot have been informed by the outcome.
Of the 141 patients that remain, 82 (58.2\%) reached eTICI 2C--3; the proportion varies across the five outer folds (0.448, 0.714, 0.536, 0.679, 0.536), so in the first fold the majority class is the absence rather than the presence of near-complete or complete reperfusion.
\NOTETODO{breakdown of how many of the 155 reviewed cases were rejected by which criterion}
\NOTETODO{number and proportion of patients with eTICI 2C--3 reperfusion in the Registry cohort, which the table leaves as an em dash}
\NOTEWARN{size of the Registry cohort: the exclusion cascade gives 1906, the stored identifier list holds 1839, and the denominators reconstructed from the decimal expansion of the reported external accuracies are 1649 (8 models), 1760 (6 models), 1761 (2 models). Neither 1906 nor 1839 is compatible with any of the sixteen transferred models, and the reconstructed values are mutually incompatible, so the external cohort is not the same across feature sets}
\par}

\clearpage

\begin{table}[htbp]
\centering
\caption{Hyperparameter search spaces of the four pipelines. LASSO, $\ell_1$-penalized linear regression on the binary outcome used as a feature selector; RFI, ranking by the gain-based feature importance of a random-forest gradient-boosting model; XGBRF, \texttt{XGBRFClassifier}; LinearSVC, linear support vector classifier with an $\ell_1$ penalty, which acts as selector and classifier simultaneously; NuSVC, $\nu$-support vector classifier with a radial basis function kernel. In the Component column, (S) marks the component acting as feature selector, (C) as classifier, and (S+C) a single object that performs both roles. Scale indicates how Optuna sampled the parameter: \emph{log}, log-uniform over the interval; \emph{int}, uniform over the integers of the interval; \emph{float}, uniform over the interval.}
\label{tab:hyperparameters}
\smallskip
\input{\tablesdir/table2_hyperparameters}
\end{table}

\smallskip
{\footnotesize
\sloppy
\noindent
\textit{Notes to Table~\ref{tab:hyperparameters}.}
$n$ and \texttt{X.shape[1]} denote the number of rows and the number of columns of the training matrix passed to the search.
Rows are patients for the global-based and distribution-based datasets and voxels for the voxel-based datasets, so the same expression yields different intervals across datasets; the number of columns reaches 20660 for the distribution-based datasets.
The minority fraction is the proportion of the smaller outcome class among the training labels of the given fold.
In the RFI--RF pipeline a single parameter dictionary configures both the model whose gain importances rank the features and the classifier fitted on the retained ones, which is why its rows are marked (S+C) rather than (C).
All search spaces were explored with Optuna using the Gaussian-process sampler (\texttt{GPSampler}) with 20 trials per dataset-by-pipeline combination.
The objective maximized was the \emph{mean accuracy over the five inner cross-validation folds}, not AUC; the inner folds were generated by \texttt{StratifiedGroupKFold} grouped by \texttt{patient\_id}.
Settings that were fixed and not searched, one per line so that the long identifiers have somewhere to break:\\
\texttt{Lasso(fit\_intercept=True, max\_iter=5000)}\\
\texttt{LinearSVC(penalty='l1', max\_iter=5000, class\_weight='balanced')}\\
\texttt{NuSVC(kernel='rbf', probability=True, class\_weight='balanced')}\\
\texttt{XGBRFClassifier(random\_state=42)}\\
Class weighting was therefore applied in the two support vector branches only.
The models were seeded with \texttt{random\_state=42}, but \texttt{GPSampler()} was instantiated without a seed, so the sequence of 20 trials is not reproducible.
\par}

\clearpage

\begin{table}[htbp]
\centering
\caption{Cross-validated performance of all 24 dataset-by-pipeline combinations in the MR CLEAN NO-IV cohort ($n = 141$), over five outer folds. Feature set: \emph{global}, eight summary statistics of each radiomic feature over the thrombus, one row per patient; \emph{distribution}, a ten-bin histogram of each radiomic feature over the thrombus, one row per patient; \emph{voxel}, the radiomic feature values themselves, one row per voxel. ROI: \emph{whole}, the dilated thrombus mask; \emph{centerline}, the dilated skeleton of the thrombus mask. Selector (Sel.): LASSO, $\ell_1$-penalized regression; L1, embedded selection by the non-zero coefficients of the classifier itself; RFI, gain-based feature importance of an \texttt{XGBRFClassifier}. Classifier (Class.): XGBRF, \texttt{XGBRFClassifier}; LinSVC, linear support vector classifier with an $\ell_1$ penalty; NuSVC, $\nu$-support vector classifier with a radial basis function kernel. In the body of the table the distribution-based feature set is abbreviated \emph{distrib.}\ and the centerline ROI \emph{cent.} AUC, area under the receiver operating characteristic curve; ACC, accuracy; IQR, interquartile range across the five folds; min--max, the smallest and largest value across the five folds; Med.\ $n$, median number of features retained by the selector; Null AUC, mean AUC of the permutation baseline; $\kappa$, the $\kappa$-type index of the observed AUC relative to that baseline; $p$, permutation $p$-value; $q$, the same $p$-value after Benjamini--Hochberg correction across all 24 combinations.}
\label{tab:cv-performance}
\smallskip
\input{\tablesdir/table3_cv_performance}
\end{table}

\smallskip
{\footnotesize
\sloppy
\noindent
\textit{Notes to Table~\ref{tab:cv-performance}.}
Accuracy must be read against the majority-class accuracy of the development cohort, $82/141 = 0.582$; only 4 of the 24 combinations exceed it by median accuracy, and 86 of the 120 individual fold results (71.7\%) fall below it.
The class ratio itself varies across folds (0.448, 0.714, 0.536, 0.679, 0.536), so in the first fold the majority class is the absence rather than the presence of near-complete or complete reperfusion.
The largest single fold result of the whole study is 0.810, obtained by distribution/whole/RFI/NuSVC in fold 4.
Averaged over the eight whole-thrombus patient-level combinations, AUC per fold is 0.510, 0.661, 0.606, 0.587, 0.623, a spread within a single fold of up to 0.531; between-fold variability therefore exceeds between-method variability.
The permutation null was obtained from 100 label permutations at the patient level per outer fold; with $m = 100$ permutations the smallest attainable $p$-value is $p = (0+1)/(100+1) = 0.0099$, which bounds the resolution of every permutation result in this table.
The Null AUC, $\kappa$, $p$ and $q$ columns were computed from the \emph{mean} AUC across folds, whereas the AUC column reports the \emph{median}; the two summaries are not interchangeable and are shown side by side deliberately.
10 of the 24 combinations reach $p < 0.05$ uncorrected and 6 survive Benjamini--Hochberg correction\NOTETODO{cite Benjamini and Hochberg 1995; the key is missing from BibFile.bib} for AUC, while none survives it for accuracy.
No centerline combination survives correction; across the eight of them the mean AUC ranges from 0.396 to 0.495 against a null of 0.489--0.510, $\kappa$ from $-0.206$ to $-0.014$ and $p$ from 0.5347 to 1.0000.\\
\textsuperscript{a} In these 2 combinations the $\ell_1$ penalty shrank every coefficient to zero in every fold and the code substituted the index of the first column; the retained feature is a fallback, not one the selector identified as informative.
Combinations with a median of one retained feature but no mark reached that median with genuinely non-zero coefficients.\\
\textsuperscript{b} The unit of analysis of the voxel-based datasets is the voxel, not the patient, and their metrics are therefore \emph{not} comparable with the remaining rows.
Across the five folds the voxel/whole datasets contribute 396085 voxel-level predictions (91500, 79327, 75748, 82369, 67141) and the voxel/centerline datasets 8557 (1690, 1571, 1811, 1806, 1679), against 141 patient-level predictions (29, 28, 28, 28, 28) for every other dataset.
The consequence is visible in the permutation null, whose standard deviation is 0.040--0.052 for the voxel-based combinations against 0.063--0.120 for all others, so that the voxel-based $p$-values are computed against a spuriously narrow null and are not comparable with the rest of the table.
3 voxel-based combinations are formally significant for AUC after correction; they must not be read as evidence for a signal.
\NOTETODO{patient-level majority-vote aggregation of the voxel-based predictions, which would make these rows comparable}
\par}

\clearpage

\begin{table}[htbp]
\centering
\caption{External validation in the MR CLEAN Registry cohort. Each pipeline was refitted on all 141 development patients with its tuned hyperparameters and applied once to the Registry cohort. CV, cross-validated median over the five outer folds of the same combination in Table~\ref{tab:cv-performance}; Reg., Registry; ACC, accuracy; $\Delta$AUC, Registry AUC minus cross-validated median AUC, a negative value being a loss of discrimination on transfer; other abbreviations as in Table~\ref{tab:cv-performance}. The voxel-based datasets were not transferred externally and therefore have no row here, which is why the table has 16 rows rather than 24.}
\label{tab:external-validation}
\smallskip
\input{\tablesdir/table4_external_validation}
\end{table}

\smallskip
{\footnotesize
\sloppy
\noindent
\textit{Notes to Table~\ref{tab:external-validation}.}
Across the 16 models the Registry AUC ranged from 0.471 to 0.521 with a median of 0.496, and 5 of them exceeded 0.50.
Registry accuracy ranged from 0.440 to 0.550 with a median of 0.481, and no model reached the majority-class accuracy of the development cohort, 0.582.
The sign of $\Delta$ follows the cross-validated value and nothing else: the 8 combinations whose cross-validated median AUC was above 0.500 all moved downwards on transfer, by $-0.034$ to $-0.190$, and the 8 that were at or below it all moved upwards, by $0.002$ to $0.128$.
The upward movement is therefore a regression to chance of models that had performed below chance internally, not a gain in generalization, and it must not be presented as one.
The only combinations for which the comparison is informative are the 6 that discriminated in cross-validation, that is reached a cross-validated median AUC of at least 0.600; all of them lose between 0.150 and 0.190 of AUC.
\NOTEWARN{$n$ Registry is not stated in the results files and is not the same for every model. Reconstructing the smallest denominator compatible with the decimal expansion of each reported accuracy leaves exactly three cohort sizes: 1649 for eight models, 1760 for six and 1761 for two. Neither the 1906 of the exclusion cascade nor the 1839 identifiers stored with the cohort is compatible with any of the sixteen. Until this is resolved the AUC values in this table are not strictly comparable with one another}
\par}

\clearpage

\begin{table}[htbp]
\centering
\caption{Comparison with the reference pipeline of van Voorst et al. \cite{vanVoorst2022} The first row summarizes the present study; rows two to four were computed on the MR CLEAN NO-IV development cohort with the same five outer folds as in Table~\ref{tab:cv-performance}, so that the comparison is made within one dataset rather than across publications; the last row is the value reported in the original publication, on a different cohort and for a different endpoint. ICC, intraclass correlation coefficient: of the 107 features the reference pipeline extracts, the 52 with ICC $\geq$ 0.75 are retained and the rest discarded. AUC$_{\text{prob}}$ denotes the AUC computed from the predicted class probabilities, and is the quantity reported for the reference pipeline. Values are medians with the interquartile range in brackets, or means $\pm$ standard deviation across folds, as indicated in the note below. An em dash marks a quantity the row does not have: the pooled row spans combinations with different accuracies, and the published row reports neither a range nor an accuracy.}
\label{tab:reference-comparison}
\smallskip
\input{\tablesdir/table5_reference_comparison}
\end{table}

\smallskip
{\footnotesize
\sloppy
\noindent
\textit{Notes to Table~\ref{tab:reference-comparison}.}
The first row pools 25 fold results, 5 combinations $\times$ 5 outer folds, and its bracketed values are the interquartile range; the second row is a mean $\pm$ standard deviation over the 5 splits; rows three and four are medians over the 5 folds, with no further feature selection applied to the 52 ICC-filtered features.
The last row is quoted from the publication: thrombus radiomics for eTICI $\geq$ 2B, a more permissive endpoint than the eTICI 2C--3 outcome of this article, as a mean over 100 bootstrap samples of 100 patients.
The plain \texttt{AUC} column of the reference pipeline degenerates to exactly 0.50 in 4 of the five splits, because the fitted model assigns the majority class to every patient in those splits and its hard predictions therefore carry no ranking information.
AUC$_{\text{prob}}$, computed from the predicted probabilities, is the only quantity of that pipeline that is comparable with the AUC values reported elsewhere in this article, and it is the one used here.
By median accuracy the reference pipeline and both of our classifiers restricted to its 52 ICC-filtered features stay below the majority-class accuracy of 0.582, although the reference pipeline exceeds it in one individual split.
Rows two to four were recomputed against the eTICI 2C--3 outcome of this article --- the stored label column is named \texttt{TICI\_2B-3}, but its per-fold positive counts are exactly the 82/141 of this study --- and are directly comparable with the first row, whereas the published 0.52 of the last row was obtained for eTICI $\geq$ 2B in a different cohort and differs from them in endpoint and cohort as much as in pipeline.
\NOTEQUERY{the pool of the first row omits one whole-thrombus combination that also discriminated in cross-validation, and the reason is not documented in the source material}
\NOTEQUERY{decide which of the two mutually incompatible sets of numbers in the source data applies before this table is finalized: set A, 0.481 $\pm$ 0.145 on 141 patients against set B, 0.467 $\pm$ 0.099 on 170}
\par}

\clearpage

\begin{table}[htbp]
\centering
\caption{Ensemble analysis: the 15 configurations of discrete, that is binary, AdaBoost \cite{Freund1997} applied to the stored prediction vectors of the 16 patient-level weak classifiers (four patient-level datasets $\times$ four pipelines). Each configuration was fitted on the five outer folds, giving 75 ensemble models. ACC $\varnothing$, mean accuracy of the configuration across the five folds; best weak classifier, the single pipeline in the configuration's pool with the highest mean accuracy, and its value; $\Delta$, the mean accuracy of the ensemble minus that of the best weak classifier. A negative $\Delta$ means the ensemble is worse than the best model it combines. The configuration name states which of the four patient-level datasets form its pool, each contributing all four pipelines: \texttt{type1} pools all four datasets (16 weak classifiers), \texttt{type2} a whole-thrombus dataset with its own centerline counterpart (8), \texttt{type3} a single dataset (4), \texttt{type4} two datasets that are not such a pair (8) and \texttt{type5} all datasets but one (12). \texttt{distrib.}\ abbreviates the distribution-based dataset and \emph{cent.} the centerline region of interest, as in Table~\ref{tab:cv-performance}. The two voxel-based datasets were excluded from every pool because their predictions are defined per voxel.}
\label{tab:ensemble}
\smallskip
\input{\tablesdir/table6_ensemble}
\end{table}

\smallskip
{\footnotesize
\sloppy
\noindent
\textit{Notes to Table~\ref{tab:ensemble}.}
$\Delta$ is negative for all 15 configurations, between $-0.022$ and $-0.121$; no ensemble outperformed the best single pipeline in its own pool.
The worst result belongs to \texttt{type1\_\allowbreak all\_\allowbreak featuresets}, the configuration that pools all feature sets.
Across all 75 fitted ensemble models accuracy averages $0.543 \pm 0.083$ with a range of 0.414--0.714.
The mechanism is documented by the fitted weights: in the all-feature-set configuration, 55 of the 80 boosting rounds (16 rounds $\times$ five folds) received a negative weight $\alpha$, meaning that the weak classifier admitted in that round performed worse than chance on the reweighted data.
The individual pipelines therefore do not make complementary errors, and the boosting assumption of weak classifiers reliably better than chance is not met.
Labels were coded as $\{-1,+1\}$, the weight of a round was $\alpha = \tfrac{1}{2}\ln\frac{1-\varepsilon}{\varepsilon}$ with $\varepsilon$ the error of the admitted classifier on the currently reweighted data, and each weak classifier could be admitted at most once, so a pool of $k$ members yields at most $k$ rounds.
The weak classifiers were not retrained: the procedure is applied to models that had already been fitted, which makes it methodologically a form of stacking rather than boosting proper.
The best weak classifier of \texttt{type3\_\allowbreak global} and of \texttt{type1\_\allowbreak all\_\allowbreak featuresets} is the same pipeline, because the pool of the latter contains the pool of the former.
\par}

\clearpage

\begin{table}[htbp]
\centering
\caption{Pairwise comparison of the pipelines by the corrected resampled $t$-test of Nadeau and Bengio \cite{Nadeau2003}. The upper block compares the four pipelines of this study with one another on the whole-thrombus global-based dataset; the lower block compares them, and our two classifiers restricted to the reference pipeline's own 52 ICC-filtered features, against the reference pipeline of van Voorst et al.\ recomputed on the same cohort. Every model in the table was evaluated on the same 5 outer folds of the same 141 patients, which is what makes the paired test applicable. Because the training sets of $k$-fold cross-validation overlap, the per-fold differences are not independent and the ordinary paired $t$-test underestimates their variance; the correction multiplies the variance by $1/k + n_{\text{test}}/n_{\text{train}}$, here $1/5 + 28.2/112.8 = 0.4500$, which inflates the standard error 1.50-fold at $\mathrm{df} = 4$. Fold sizes are taken from the data, not assumed. Mean $\Delta$AUC is the mean of the 5 per-fold differences in AUC between pipeline A and pipeline B; a positive value favours A. $t$ (uncorrected) is the ordinary paired statistic, $t$ (corrected) the same quantity with the inflated variance. The two-sided critical value at $\alpha = 0.05$ and $\mathrm{df} = 4$ is 2.776. Pipeline names combine selector and classifier as in Table~\ref{tab:cv-performance}.}
\label{tab:nadeau-bengio}
\smallskip
\input{\tablesdir/table7_nadeau_bengio}
\end{table}

\smallskip
{\footnotesize
\sloppy
\noindent
\textit{Notes to Table~\ref{tab:nadeau-bengio}.}
Not one of the 12 comparisons reaches significance: the largest corrected statistic in absolute value is 1.54, against a critical value of 2.776.
The four pipelines of this study are therefore statistically indistinguishable from one another, and none of them is distinguishable from the reference pipeline either, in either direction.
For the reference pipeline the compared quantity is AUC$_{\text{prob}}$, because its plain \texttt{AUC} column degenerates to exactly 0.50 in four of the five splits; see Table~\ref{tab:reference-comparison}.
The pooled row of Table~\ref{tab:reference-comparison} is deliberately absent here: those 5 combinations were picked by their results on these same folds, so testing them against the reference pipeline would be biased in favour of the present study. All four pipelines of the dataset are tested instead, whatever their outcome.
With 5 folds the test has $\mathrm{df} = 4$ and very little power, so these results are an absence of evidence for a difference, not evidence that the pipelines perform identically. The comparison that the data would support best, a DeLong test on the external ROC curves, was not possible because per-patient predicted scores on the Registry were not retained.
\par}

\clearpage


%%
%% Tabulky se vejdou i do uzke sazby `elsarticle` (390 pt), ale zdrobni.
%% Aby vypadaly jako v samostatnem supplementu, obal se vlozeni geometrii:
%%
%%     \usepackage{geometry}          % v preambuli, za elsarticle
%%     ...
%%     \clearpage
%%     \newgeometry{margin=15mm}
%%     %% -------------------------------------------------------------------------
%% GENEROVANO: make_all.py -- needitovat rucne.
%%
%% Popisky, tabulky a poznamky bez preambule. Do clanku staci jeden radek:
%%
%%     %% -------------------------------------------------------------------------
%% GENEROVANO: make_all.py -- needitovat rucne.
%%
%% Popisky, tabulky a poznamky bez preambule. Do clanku staci jeden radek:
%%
%%     \input{tables_tom/output/supplement_body}
%%
%% Tabulky se vejdou i do uzke sazby `elsarticle` (390 pt), ale zdrobni.
%% Aby vypadaly jako v samostatnem supplementu, obal se vlozeni geometrii:
%%
%%     \usepackage{geometry}          % v preambuli, za elsarticle
%%     ...
%%     \clearpage
%%     \newgeometry{margin=15mm}
%%     \input{tables_tom/output/supplement_body}
%%     \restoregeometry
%%
%% Cislovani S1--S7 misto 1--7 zaridi pred vlozenim:
%%     \setcounter{table}{0}\renewcommand{\thetable}{S\arabic{table}}
%%
%% Preambule potrebuje: booktabs, array, amsmath, amssymb.
%% Cesta k tabulkam se bere z \tablesdir; vychozi hodnota plati pro preklad
%% z korene projektu.
%% -------------------------------------------------------------------------
\providecommand{\tablesdir}{tables_tom/output}

%% Zaloha, kdyby `manuscriptnotes` nebyl nactveny: vysadi se jen kotva,
%% poznamka se zahodi -- tedy tentyz vysledek jako v rezimu odesilani.
\providecommand{\NOTE}[2][]{#1}
\providecommand{\NOTEWARN}[2][]{#1}
\providecommand{\NOTECHECK}[2][]{#1}
\providecommand{\NOTETODO}[2][]{#1}
\providecommand{\NOTEQUERY}[2][]{#1}

{\large\bfseries Supplementary tables\par}
\bigskip
\NOTEWARN{the captions and notes in this supplement were drafted with AI assistance and are provisional; every one of them has to be rewritten and checked by an author before submission}

\begin{table}[htbp]
\centering
\caption{Patient selection in the two cohorts. MR CLEAN NO-IV was used for model development and cross-validation, the MR CLEAN Registry as the independent external cohort. Entries are numbers of patients; an em dash marks an exclusion step that does not apply to that cohort. eTICI, expanded Thrombolysis in Cerebral Infarction score; the outcome modeled throughout this article is near-complete or complete reperfusion, defined as a final eTICI score of 2C or 3. Slice spacing refers to the through-plane spacing of the non-contrast CT acquisition.}
\label{tab:cohorts}
\smallskip
\input{\tablesdir/table1_cohorts}
\end{table}

\smallskip
{\footnotesize
\sloppy
\noindent
\textit{Notes to Table~\ref{tab:cohorts}.}
The manual review of segmentation quality was carried out on the MR CLEAN NO-IV cohort only.
No segmentation was reviewed or rejected in the MR CLEAN Registry cohort, so the two cohorts differ in segmentation quality control as well as in population.
This is a property of the design and is the single largest asymmetry between the development and the external cohort.
The review was blinded to the eTICI score: it judged the quality and the anatomical placement of the segmentation only, so the 155 exclusions it produced cannot have been informed by the outcome.
Of the 141 patients that remain, 82 (58.2\%) reached eTICI 2C--3; the proportion varies across the five outer folds (0.448, 0.714, 0.536, 0.679, 0.536), so in the first fold the majority class is the absence rather than the presence of near-complete or complete reperfusion.
\NOTETODO{breakdown of how many of the 155 reviewed cases were rejected by which criterion}
\NOTETODO{number and proportion of patients with eTICI 2C--3 reperfusion in the Registry cohort, which the table leaves as an em dash}
\NOTEWARN{size of the Registry cohort: the exclusion cascade gives 1906, the stored identifier list holds 1839, and the denominators reconstructed from the decimal expansion of the reported external accuracies are 1649 (8 models), 1760 (6 models), 1761 (2 models). Neither 1906 nor 1839 is compatible with any of the sixteen transferred models, and the reconstructed values are mutually incompatible, so the external cohort is not the same across feature sets}
\par}

\clearpage

\begin{table}[htbp]
\centering
\caption{Hyperparameter search spaces of the four pipelines. LASSO, $\ell_1$-penalized linear regression on the binary outcome used as a feature selector; RFI, ranking by the gain-based feature importance of a random-forest gradient-boosting model; XGBRF, \texttt{XGBRFClassifier}; LinearSVC, linear support vector classifier with an $\ell_1$ penalty, which acts as selector and classifier simultaneously; NuSVC, $\nu$-support vector classifier with a radial basis function kernel. In the Component column, (S) marks the component acting as feature selector, (C) as classifier, and (S+C) a single object that performs both roles. Scale indicates how Optuna sampled the parameter: \emph{log}, log-uniform over the interval; \emph{int}, uniform over the integers of the interval; \emph{float}, uniform over the interval.}
\label{tab:hyperparameters}
\smallskip
\input{\tablesdir/table2_hyperparameters}
\end{table}

\smallskip
{\footnotesize
\sloppy
\noindent
\textit{Notes to Table~\ref{tab:hyperparameters}.}
$n$ and \texttt{X.shape[1]} denote the number of rows and the number of columns of the training matrix passed to the search.
Rows are patients for the global-based and distribution-based datasets and voxels for the voxel-based datasets, so the same expression yields different intervals across datasets; the number of columns reaches 20660 for the distribution-based datasets.
The minority fraction is the proportion of the smaller outcome class among the training labels of the given fold.
In the RFI--RF pipeline a single parameter dictionary configures both the model whose gain importances rank the features and the classifier fitted on the retained ones, which is why its rows are marked (S+C) rather than (C).
All search spaces were explored with Optuna using the Gaussian-process sampler (\texttt{GPSampler}) with 20 trials per dataset-by-pipeline combination.
The objective maximized was the \emph{mean accuracy over the five inner cross-validation folds}, not AUC; the inner folds were generated by \texttt{StratifiedGroupKFold} grouped by \texttt{patient\_id}.
Settings that were fixed and not searched, one per line so that the long identifiers have somewhere to break:\\
\texttt{Lasso(fit\_intercept=True, max\_iter=5000)}\\
\texttt{LinearSVC(penalty='l1', max\_iter=5000, class\_weight='balanced')}\\
\texttt{NuSVC(kernel='rbf', probability=True, class\_weight='balanced')}\\
\texttt{XGBRFClassifier(random\_state=42)}\\
Class weighting was therefore applied in the two support vector branches only.
The models were seeded with \texttt{random\_state=42}, but \texttt{GPSampler()} was instantiated without a seed, so the sequence of 20 trials is not reproducible.
\par}

\clearpage

\begin{table}[htbp]
\centering
\caption{Cross-validated performance of all 24 dataset-by-pipeline combinations in the MR CLEAN NO-IV cohort ($n = 141$), over five outer folds. Feature set: \emph{global}, eight summary statistics of each radiomic feature over the thrombus, one row per patient; \emph{distribution}, a ten-bin histogram of each radiomic feature over the thrombus, one row per patient; \emph{voxel}, the radiomic feature values themselves, one row per voxel. ROI: \emph{whole}, the dilated thrombus mask; \emph{centerline}, the dilated skeleton of the thrombus mask. Selector (Sel.): LASSO, $\ell_1$-penalized regression; L1, embedded selection by the non-zero coefficients of the classifier itself; RFI, gain-based feature importance of an \texttt{XGBRFClassifier}. Classifier (Class.): XGBRF, \texttt{XGBRFClassifier}; LinSVC, linear support vector classifier with an $\ell_1$ penalty; NuSVC, $\nu$-support vector classifier with a radial basis function kernel. In the body of the table the distribution-based feature set is abbreviated \emph{distrib.}\ and the centerline ROI \emph{cent.} AUC, area under the receiver operating characteristic curve; ACC, accuracy; IQR, interquartile range across the five folds; min--max, the smallest and largest value across the five folds; Med.\ $n$, median number of features retained by the selector; Null AUC, mean AUC of the permutation baseline; $\kappa$, the $\kappa$-type index of the observed AUC relative to that baseline; $p$, permutation $p$-value; $q$, the same $p$-value after Benjamini--Hochberg correction across all 24 combinations.}
\label{tab:cv-performance}
\smallskip
\input{\tablesdir/table3_cv_performance}
\end{table}

\smallskip
{\footnotesize
\sloppy
\noindent
\textit{Notes to Table~\ref{tab:cv-performance}.}
Accuracy must be read against the majority-class accuracy of the development cohort, $82/141 = 0.582$; only 4 of the 24 combinations exceed it by median accuracy, and 86 of the 120 individual fold results (71.7\%) fall below it.
The class ratio itself varies across folds (0.448, 0.714, 0.536, 0.679, 0.536), so in the first fold the majority class is the absence rather than the presence of near-complete or complete reperfusion.
The largest single fold result of the whole study is 0.810, obtained by distribution/whole/RFI/NuSVC in fold 4.
Averaged over the eight whole-thrombus patient-level combinations, AUC per fold is 0.510, 0.661, 0.606, 0.587, 0.623, a spread within a single fold of up to 0.531; between-fold variability therefore exceeds between-method variability.
The permutation null was obtained from 100 label permutations at the patient level per outer fold; with $m = 100$ permutations the smallest attainable $p$-value is $p = (0+1)/(100+1) = 0.0099$, which bounds the resolution of every permutation result in this table.
The Null AUC, $\kappa$, $p$ and $q$ columns were computed from the \emph{mean} AUC across folds, whereas the AUC column reports the \emph{median}; the two summaries are not interchangeable and are shown side by side deliberately.
10 of the 24 combinations reach $p < 0.05$ uncorrected and 6 survive Benjamini--Hochberg correction\NOTETODO{cite Benjamini and Hochberg 1995; the key is missing from BibFile.bib} for AUC, while none survives it for accuracy.
No centerline combination survives correction; across the eight of them the mean AUC ranges from 0.396 to 0.495 against a null of 0.489--0.510, $\kappa$ from $-0.206$ to $-0.014$ and $p$ from 0.5347 to 1.0000.\\
\textsuperscript{a} In these 2 combinations the $\ell_1$ penalty shrank every coefficient to zero in every fold and the code substituted the index of the first column; the retained feature is a fallback, not one the selector identified as informative.
Combinations with a median of one retained feature but no mark reached that median with genuinely non-zero coefficients.\\
\textsuperscript{b} The unit of analysis of the voxel-based datasets is the voxel, not the patient, and their metrics are therefore \emph{not} comparable with the remaining rows.
Across the five folds the voxel/whole datasets contribute 396085 voxel-level predictions (91500, 79327, 75748, 82369, 67141) and the voxel/centerline datasets 8557 (1690, 1571, 1811, 1806, 1679), against 141 patient-level predictions (29, 28, 28, 28, 28) for every other dataset.
The consequence is visible in the permutation null, whose standard deviation is 0.040--0.052 for the voxel-based combinations against 0.063--0.120 for all others, so that the voxel-based $p$-values are computed against a spuriously narrow null and are not comparable with the rest of the table.
3 voxel-based combinations are formally significant for AUC after correction; they must not be read as evidence for a signal.
\NOTETODO{patient-level majority-vote aggregation of the voxel-based predictions, which would make these rows comparable}
\par}

\clearpage

\begin{table}[htbp]
\centering
\caption{External validation in the MR CLEAN Registry cohort. Each pipeline was refitted on all 141 development patients with its tuned hyperparameters and applied once to the Registry cohort. CV, cross-validated median over the five outer folds of the same combination in Table~\ref{tab:cv-performance}; Reg., Registry; ACC, accuracy; $\Delta$AUC, Registry AUC minus cross-validated median AUC, a negative value being a loss of discrimination on transfer; other abbreviations as in Table~\ref{tab:cv-performance}. The voxel-based datasets were not transferred externally and therefore have no row here, which is why the table has 16 rows rather than 24.}
\label{tab:external-validation}
\smallskip
\input{\tablesdir/table4_external_validation}
\end{table}

\smallskip
{\footnotesize
\sloppy
\noindent
\textit{Notes to Table~\ref{tab:external-validation}.}
Across the 16 models the Registry AUC ranged from 0.471 to 0.521 with a median of 0.496, and 5 of them exceeded 0.50.
Registry accuracy ranged from 0.440 to 0.550 with a median of 0.481, and no model reached the majority-class accuracy of the development cohort, 0.582.
The sign of $\Delta$ follows the cross-validated value and nothing else: the 8 combinations whose cross-validated median AUC was above 0.500 all moved downwards on transfer, by $-0.034$ to $-0.190$, and the 8 that were at or below it all moved upwards, by $0.002$ to $0.128$.
The upward movement is therefore a regression to chance of models that had performed below chance internally, not a gain in generalization, and it must not be presented as one.
The only combinations for which the comparison is informative are the 6 that discriminated in cross-validation, that is reached a cross-validated median AUC of at least 0.600; all of them lose between 0.150 and 0.190 of AUC.
\NOTEWARN{$n$ Registry is not stated in the results files and is not the same for every model. Reconstructing the smallest denominator compatible with the decimal expansion of each reported accuracy leaves exactly three cohort sizes: 1649 for eight models, 1760 for six and 1761 for two. Neither the 1906 of the exclusion cascade nor the 1839 identifiers stored with the cohort is compatible with any of the sixteen. Until this is resolved the AUC values in this table are not strictly comparable with one another}
\par}

\clearpage

\begin{table}[htbp]
\centering
\caption{Comparison with the reference pipeline of van Voorst et al. \cite{vanVoorst2022} The first row summarizes the present study; rows two to four were computed on the MR CLEAN NO-IV development cohort with the same five outer folds as in Table~\ref{tab:cv-performance}, so that the comparison is made within one dataset rather than across publications; the last row is the value reported in the original publication, on a different cohort and for a different endpoint. ICC, intraclass correlation coefficient: of the 107 features the reference pipeline extracts, the 52 with ICC $\geq$ 0.75 are retained and the rest discarded. AUC$_{\text{prob}}$ denotes the AUC computed from the predicted class probabilities, and is the quantity reported for the reference pipeline. Values are medians with the interquartile range in brackets, or means $\pm$ standard deviation across folds, as indicated in the note below. An em dash marks a quantity the row does not have: the pooled row spans combinations with different accuracies, and the published row reports neither a range nor an accuracy.}
\label{tab:reference-comparison}
\smallskip
\input{\tablesdir/table5_reference_comparison}
\end{table}

\smallskip
{\footnotesize
\sloppy
\noindent
\textit{Notes to Table~\ref{tab:reference-comparison}.}
The first row pools 25 fold results, 5 combinations $\times$ 5 outer folds, and its bracketed values are the interquartile range; the second row is a mean $\pm$ standard deviation over the 5 splits; rows three and four are medians over the 5 folds, with no further feature selection applied to the 52 ICC-filtered features.
The last row is quoted from the publication: thrombus radiomics for eTICI $\geq$ 2B, a more permissive endpoint than the eTICI 2C--3 outcome of this article, as a mean over 100 bootstrap samples of 100 patients.
The plain \texttt{AUC} column of the reference pipeline degenerates to exactly 0.50 in 4 of the five splits, because the fitted model assigns the majority class to every patient in those splits and its hard predictions therefore carry no ranking information.
AUC$_{\text{prob}}$, computed from the predicted probabilities, is the only quantity of that pipeline that is comparable with the AUC values reported elsewhere in this article, and it is the one used here.
By median accuracy the reference pipeline and both of our classifiers restricted to its 52 ICC-filtered features stay below the majority-class accuracy of 0.582, although the reference pipeline exceeds it in one individual split.
Rows two to four were recomputed against the eTICI 2C--3 outcome of this article --- the stored label column is named \texttt{TICI\_2B-3}, but its per-fold positive counts are exactly the 82/141 of this study --- and are directly comparable with the first row, whereas the published 0.52 of the last row was obtained for eTICI $\geq$ 2B in a different cohort and differs from them in endpoint and cohort as much as in pipeline.
\NOTEQUERY{the pool of the first row omits one whole-thrombus combination that also discriminated in cross-validation, and the reason is not documented in the source material}
\NOTEQUERY{decide which of the two mutually incompatible sets of numbers in the source data applies before this table is finalized: set A, 0.481 $\pm$ 0.145 on 141 patients against set B, 0.467 $\pm$ 0.099 on 170}
\par}

\clearpage

\begin{table}[htbp]
\centering
\caption{Ensemble analysis: the 15 configurations of discrete, that is binary, AdaBoost \cite{Freund1997} applied to the stored prediction vectors of the 16 patient-level weak classifiers (four patient-level datasets $\times$ four pipelines). Each configuration was fitted on the five outer folds, giving 75 ensemble models. ACC $\varnothing$, mean accuracy of the configuration across the five folds; best weak classifier, the single pipeline in the configuration's pool with the highest mean accuracy, and its value; $\Delta$, the mean accuracy of the ensemble minus that of the best weak classifier. A negative $\Delta$ means the ensemble is worse than the best model it combines. The configuration name states which of the four patient-level datasets form its pool, each contributing all four pipelines: \texttt{type1} pools all four datasets (16 weak classifiers), \texttt{type2} a whole-thrombus dataset with its own centerline counterpart (8), \texttt{type3} a single dataset (4), \texttt{type4} two datasets that are not such a pair (8) and \texttt{type5} all datasets but one (12). \texttt{distrib.}\ abbreviates the distribution-based dataset and \emph{cent.} the centerline region of interest, as in Table~\ref{tab:cv-performance}. The two voxel-based datasets were excluded from every pool because their predictions are defined per voxel.}
\label{tab:ensemble}
\smallskip
\input{\tablesdir/table6_ensemble}
\end{table}

\smallskip
{\footnotesize
\sloppy
\noindent
\textit{Notes to Table~\ref{tab:ensemble}.}
$\Delta$ is negative for all 15 configurations, between $-0.022$ and $-0.121$; no ensemble outperformed the best single pipeline in its own pool.
The worst result belongs to \texttt{type1\_\allowbreak all\_\allowbreak featuresets}, the configuration that pools all feature sets.
Across all 75 fitted ensemble models accuracy averages $0.543 \pm 0.083$ with a range of 0.414--0.714.
The mechanism is documented by the fitted weights: in the all-feature-set configuration, 55 of the 80 boosting rounds (16 rounds $\times$ five folds) received a negative weight $\alpha$, meaning that the weak classifier admitted in that round performed worse than chance on the reweighted data.
The individual pipelines therefore do not make complementary errors, and the boosting assumption of weak classifiers reliably better than chance is not met.
Labels were coded as $\{-1,+1\}$, the weight of a round was $\alpha = \tfrac{1}{2}\ln\frac{1-\varepsilon}{\varepsilon}$ with $\varepsilon$ the error of the admitted classifier on the currently reweighted data, and each weak classifier could be admitted at most once, so a pool of $k$ members yields at most $k$ rounds.
The weak classifiers were not retrained: the procedure is applied to models that had already been fitted, which makes it methodologically a form of stacking rather than boosting proper.
The best weak classifier of \texttt{type3\_\allowbreak global} and of \texttt{type1\_\allowbreak all\_\allowbreak featuresets} is the same pipeline, because the pool of the latter contains the pool of the former.
\par}

\clearpage

\begin{table}[htbp]
\centering
\caption{Pairwise comparison of the pipelines by the corrected resampled $t$-test of Nadeau and Bengio \cite{Nadeau2003}. The upper block compares the four pipelines of this study with one another on the whole-thrombus global-based dataset; the lower block compares them, and our two classifiers restricted to the reference pipeline's own 52 ICC-filtered features, against the reference pipeline of van Voorst et al.\ recomputed on the same cohort. Every model in the table was evaluated on the same 5 outer folds of the same 141 patients, which is what makes the paired test applicable. Because the training sets of $k$-fold cross-validation overlap, the per-fold differences are not independent and the ordinary paired $t$-test underestimates their variance; the correction multiplies the variance by $1/k + n_{\text{test}}/n_{\text{train}}$, here $1/5 + 28.2/112.8 = 0.4500$, which inflates the standard error 1.50-fold at $\mathrm{df} = 4$. Fold sizes are taken from the data, not assumed. Mean $\Delta$AUC is the mean of the 5 per-fold differences in AUC between pipeline A and pipeline B; a positive value favours A. $t$ (uncorrected) is the ordinary paired statistic, $t$ (corrected) the same quantity with the inflated variance. The two-sided critical value at $\alpha = 0.05$ and $\mathrm{df} = 4$ is 2.776. Pipeline names combine selector and classifier as in Table~\ref{tab:cv-performance}.}
\label{tab:nadeau-bengio}
\smallskip
\input{\tablesdir/table7_nadeau_bengio}
\end{table}

\smallskip
{\footnotesize
\sloppy
\noindent
\textit{Notes to Table~\ref{tab:nadeau-bengio}.}
Not one of the 12 comparisons reaches significance: the largest corrected statistic in absolute value is 1.54, against a critical value of 2.776.
The four pipelines of this study are therefore statistically indistinguishable from one another, and none of them is distinguishable from the reference pipeline either, in either direction.
For the reference pipeline the compared quantity is AUC$_{\text{prob}}$, because its plain \texttt{AUC} column degenerates to exactly 0.50 in four of the five splits; see Table~\ref{tab:reference-comparison}.
The pooled row of Table~\ref{tab:reference-comparison} is deliberately absent here: those 5 combinations were picked by their results on these same folds, so testing them against the reference pipeline would be biased in favour of the present study. All four pipelines of the dataset are tested instead, whatever their outcome.
With 5 folds the test has $\mathrm{df} = 4$ and very little power, so these results are an absence of evidence for a difference, not evidence that the pipelines perform identically. The comparison that the data would support best, a DeLong test on the external ROC curves, was not possible because per-patient predicted scores on the Registry were not retained.
\par}

\clearpage


%%
%% Tabulky se vejdou i do uzke sazby `elsarticle` (390 pt), ale zdrobni.
%% Aby vypadaly jako v samostatnem supplementu, obal se vlozeni geometrii:
%%
%%     \usepackage{geometry}          % v preambuli, za elsarticle
%%     ...
%%     \clearpage
%%     \newgeometry{margin=15mm}
%%     %% -------------------------------------------------------------------------
%% GENEROVANO: make_all.py -- needitovat rucne.
%%
%% Popisky, tabulky a poznamky bez preambule. Do clanku staci jeden radek:
%%
%%     \input{tables_tom/output/supplement_body}
%%
%% Tabulky se vejdou i do uzke sazby `elsarticle` (390 pt), ale zdrobni.
%% Aby vypadaly jako v samostatnem supplementu, obal se vlozeni geometrii:
%%
%%     \usepackage{geometry}          % v preambuli, za elsarticle
%%     ...
%%     \clearpage
%%     \newgeometry{margin=15mm}
%%     \input{tables_tom/output/supplement_body}
%%     \restoregeometry
%%
%% Cislovani S1--S7 misto 1--7 zaridi pred vlozenim:
%%     \setcounter{table}{0}\renewcommand{\thetable}{S\arabic{table}}
%%
%% Preambule potrebuje: booktabs, array, amsmath, amssymb.
%% Cesta k tabulkam se bere z \tablesdir; vychozi hodnota plati pro preklad
%% z korene projektu.
%% -------------------------------------------------------------------------
\providecommand{\tablesdir}{tables_tom/output}

%% Zaloha, kdyby `manuscriptnotes` nebyl nactveny: vysadi se jen kotva,
%% poznamka se zahodi -- tedy tentyz vysledek jako v rezimu odesilani.
\providecommand{\NOTE}[2][]{#1}
\providecommand{\NOTEWARN}[2][]{#1}
\providecommand{\NOTECHECK}[2][]{#1}
\providecommand{\NOTETODO}[2][]{#1}
\providecommand{\NOTEQUERY}[2][]{#1}

{\large\bfseries Supplementary tables\par}
\bigskip
\NOTEWARN{the captions and notes in this supplement were drafted with AI assistance and are provisional; every one of them has to be rewritten and checked by an author before submission}

\begin{table}[htbp]
\centering
\caption{Patient selection in the two cohorts. MR CLEAN NO-IV was used for model development and cross-validation, the MR CLEAN Registry as the independent external cohort. Entries are numbers of patients; an em dash marks an exclusion step that does not apply to that cohort. eTICI, expanded Thrombolysis in Cerebral Infarction score; the outcome modeled throughout this article is near-complete or complete reperfusion, defined as a final eTICI score of 2C or 3. Slice spacing refers to the through-plane spacing of the non-contrast CT acquisition.}
\label{tab:cohorts}
\smallskip
\input{\tablesdir/table1_cohorts}
\end{table}

\smallskip
{\footnotesize
\sloppy
\noindent
\textit{Notes to Table~\ref{tab:cohorts}.}
The manual review of segmentation quality was carried out on the MR CLEAN NO-IV cohort only.
No segmentation was reviewed or rejected in the MR CLEAN Registry cohort, so the two cohorts differ in segmentation quality control as well as in population.
This is a property of the design and is the single largest asymmetry between the development and the external cohort.
The review was blinded to the eTICI score: it judged the quality and the anatomical placement of the segmentation only, so the 155 exclusions it produced cannot have been informed by the outcome.
Of the 141 patients that remain, 82 (58.2\%) reached eTICI 2C--3; the proportion varies across the five outer folds (0.448, 0.714, 0.536, 0.679, 0.536), so in the first fold the majority class is the absence rather than the presence of near-complete or complete reperfusion.
\NOTETODO{breakdown of how many of the 155 reviewed cases were rejected by which criterion}
\NOTETODO{number and proportion of patients with eTICI 2C--3 reperfusion in the Registry cohort, which the table leaves as an em dash}
\NOTEWARN{size of the Registry cohort: the exclusion cascade gives 1906, the stored identifier list holds 1839, and the denominators reconstructed from the decimal expansion of the reported external accuracies are 1649 (8 models), 1760 (6 models), 1761 (2 models). Neither 1906 nor 1839 is compatible with any of the sixteen transferred models, and the reconstructed values are mutually incompatible, so the external cohort is not the same across feature sets}
\par}

\clearpage

\begin{table}[htbp]
\centering
\caption{Hyperparameter search spaces of the four pipelines. LASSO, $\ell_1$-penalized linear regression on the binary outcome used as a feature selector; RFI, ranking by the gain-based feature importance of a random-forest gradient-boosting model; XGBRF, \texttt{XGBRFClassifier}; LinearSVC, linear support vector classifier with an $\ell_1$ penalty, which acts as selector and classifier simultaneously; NuSVC, $\nu$-support vector classifier with a radial basis function kernel. In the Component column, (S) marks the component acting as feature selector, (C) as classifier, and (S+C) a single object that performs both roles. Scale indicates how Optuna sampled the parameter: \emph{log}, log-uniform over the interval; \emph{int}, uniform over the integers of the interval; \emph{float}, uniform over the interval.}
\label{tab:hyperparameters}
\smallskip
\input{\tablesdir/table2_hyperparameters}
\end{table}

\smallskip
{\footnotesize
\sloppy
\noindent
\textit{Notes to Table~\ref{tab:hyperparameters}.}
$n$ and \texttt{X.shape[1]} denote the number of rows and the number of columns of the training matrix passed to the search.
Rows are patients for the global-based and distribution-based datasets and voxels for the voxel-based datasets, so the same expression yields different intervals across datasets; the number of columns reaches 20660 for the distribution-based datasets.
The minority fraction is the proportion of the smaller outcome class among the training labels of the given fold.
In the RFI--RF pipeline a single parameter dictionary configures both the model whose gain importances rank the features and the classifier fitted on the retained ones, which is why its rows are marked (S+C) rather than (C).
All search spaces were explored with Optuna using the Gaussian-process sampler (\texttt{GPSampler}) with 20 trials per dataset-by-pipeline combination.
The objective maximized was the \emph{mean accuracy over the five inner cross-validation folds}, not AUC; the inner folds were generated by \texttt{StratifiedGroupKFold} grouped by \texttt{patient\_id}.
Settings that were fixed and not searched, one per line so that the long identifiers have somewhere to break:\\
\texttt{Lasso(fit\_intercept=True, max\_iter=5000)}\\
\texttt{LinearSVC(penalty='l1', max\_iter=5000, class\_weight='balanced')}\\
\texttt{NuSVC(kernel='rbf', probability=True, class\_weight='balanced')}\\
\texttt{XGBRFClassifier(random\_state=42)}\\
Class weighting was therefore applied in the two support vector branches only.
The models were seeded with \texttt{random\_state=42}, but \texttt{GPSampler()} was instantiated without a seed, so the sequence of 20 trials is not reproducible.
\par}

\clearpage

\begin{table}[htbp]
\centering
\caption{Cross-validated performance of all 24 dataset-by-pipeline combinations in the MR CLEAN NO-IV cohort ($n = 141$), over five outer folds. Feature set: \emph{global}, eight summary statistics of each radiomic feature over the thrombus, one row per patient; \emph{distribution}, a ten-bin histogram of each radiomic feature over the thrombus, one row per patient; \emph{voxel}, the radiomic feature values themselves, one row per voxel. ROI: \emph{whole}, the dilated thrombus mask; \emph{centerline}, the dilated skeleton of the thrombus mask. Selector (Sel.): LASSO, $\ell_1$-penalized regression; L1, embedded selection by the non-zero coefficients of the classifier itself; RFI, gain-based feature importance of an \texttt{XGBRFClassifier}. Classifier (Class.): XGBRF, \texttt{XGBRFClassifier}; LinSVC, linear support vector classifier with an $\ell_1$ penalty; NuSVC, $\nu$-support vector classifier with a radial basis function kernel. In the body of the table the distribution-based feature set is abbreviated \emph{distrib.}\ and the centerline ROI \emph{cent.} AUC, area under the receiver operating characteristic curve; ACC, accuracy; IQR, interquartile range across the five folds; min--max, the smallest and largest value across the five folds; Med.\ $n$, median number of features retained by the selector; Null AUC, mean AUC of the permutation baseline; $\kappa$, the $\kappa$-type index of the observed AUC relative to that baseline; $p$, permutation $p$-value; $q$, the same $p$-value after Benjamini--Hochberg correction across all 24 combinations.}
\label{tab:cv-performance}
\smallskip
\input{\tablesdir/table3_cv_performance}
\end{table}

\smallskip
{\footnotesize
\sloppy
\noindent
\textit{Notes to Table~\ref{tab:cv-performance}.}
Accuracy must be read against the majority-class accuracy of the development cohort, $82/141 = 0.582$; only 4 of the 24 combinations exceed it by median accuracy, and 86 of the 120 individual fold results (71.7\%) fall below it.
The class ratio itself varies across folds (0.448, 0.714, 0.536, 0.679, 0.536), so in the first fold the majority class is the absence rather than the presence of near-complete or complete reperfusion.
The largest single fold result of the whole study is 0.810, obtained by distribution/whole/RFI/NuSVC in fold 4.
Averaged over the eight whole-thrombus patient-level combinations, AUC per fold is 0.510, 0.661, 0.606, 0.587, 0.623, a spread within a single fold of up to 0.531; between-fold variability therefore exceeds between-method variability.
The permutation null was obtained from 100 label permutations at the patient level per outer fold; with $m = 100$ permutations the smallest attainable $p$-value is $p = (0+1)/(100+1) = 0.0099$, which bounds the resolution of every permutation result in this table.
The Null AUC, $\kappa$, $p$ and $q$ columns were computed from the \emph{mean} AUC across folds, whereas the AUC column reports the \emph{median}; the two summaries are not interchangeable and are shown side by side deliberately.
10 of the 24 combinations reach $p < 0.05$ uncorrected and 6 survive Benjamini--Hochberg correction\NOTETODO{cite Benjamini and Hochberg 1995; the key is missing from BibFile.bib} for AUC, while none survives it for accuracy.
No centerline combination survives correction; across the eight of them the mean AUC ranges from 0.396 to 0.495 against a null of 0.489--0.510, $\kappa$ from $-0.206$ to $-0.014$ and $p$ from 0.5347 to 1.0000.\\
\textsuperscript{a} In these 2 combinations the $\ell_1$ penalty shrank every coefficient to zero in every fold and the code substituted the index of the first column; the retained feature is a fallback, not one the selector identified as informative.
Combinations with a median of one retained feature but no mark reached that median with genuinely non-zero coefficients.\\
\textsuperscript{b} The unit of analysis of the voxel-based datasets is the voxel, not the patient, and their metrics are therefore \emph{not} comparable with the remaining rows.
Across the five folds the voxel/whole datasets contribute 396085 voxel-level predictions (91500, 79327, 75748, 82369, 67141) and the voxel/centerline datasets 8557 (1690, 1571, 1811, 1806, 1679), against 141 patient-level predictions (29, 28, 28, 28, 28) for every other dataset.
The consequence is visible in the permutation null, whose standard deviation is 0.040--0.052 for the voxel-based combinations against 0.063--0.120 for all others, so that the voxel-based $p$-values are computed against a spuriously narrow null and are not comparable with the rest of the table.
3 voxel-based combinations are formally significant for AUC after correction; they must not be read as evidence for a signal.
\NOTETODO{patient-level majority-vote aggregation of the voxel-based predictions, which would make these rows comparable}
\par}

\clearpage

\begin{table}[htbp]
\centering
\caption{External validation in the MR CLEAN Registry cohort. Each pipeline was refitted on all 141 development patients with its tuned hyperparameters and applied once to the Registry cohort. CV, cross-validated median over the five outer folds of the same combination in Table~\ref{tab:cv-performance}; Reg., Registry; ACC, accuracy; $\Delta$AUC, Registry AUC minus cross-validated median AUC, a negative value being a loss of discrimination on transfer; other abbreviations as in Table~\ref{tab:cv-performance}. The voxel-based datasets were not transferred externally and therefore have no row here, which is why the table has 16 rows rather than 24.}
\label{tab:external-validation}
\smallskip
\input{\tablesdir/table4_external_validation}
\end{table}

\smallskip
{\footnotesize
\sloppy
\noindent
\textit{Notes to Table~\ref{tab:external-validation}.}
Across the 16 models the Registry AUC ranged from 0.471 to 0.521 with a median of 0.496, and 5 of them exceeded 0.50.
Registry accuracy ranged from 0.440 to 0.550 with a median of 0.481, and no model reached the majority-class accuracy of the development cohort, 0.582.
The sign of $\Delta$ follows the cross-validated value and nothing else: the 8 combinations whose cross-validated median AUC was above 0.500 all moved downwards on transfer, by $-0.034$ to $-0.190$, and the 8 that were at or below it all moved upwards, by $0.002$ to $0.128$.
The upward movement is therefore a regression to chance of models that had performed below chance internally, not a gain in generalization, and it must not be presented as one.
The only combinations for which the comparison is informative are the 6 that discriminated in cross-validation, that is reached a cross-validated median AUC of at least 0.600; all of them lose between 0.150 and 0.190 of AUC.
\NOTEWARN{$n$ Registry is not stated in the results files and is not the same for every model. Reconstructing the smallest denominator compatible with the decimal expansion of each reported accuracy leaves exactly three cohort sizes: 1649 for eight models, 1760 for six and 1761 for two. Neither the 1906 of the exclusion cascade nor the 1839 identifiers stored with the cohort is compatible with any of the sixteen. Until this is resolved the AUC values in this table are not strictly comparable with one another}
\par}

\clearpage

\begin{table}[htbp]
\centering
\caption{Comparison with the reference pipeline of van Voorst et al. \cite{vanVoorst2022} The first row summarizes the present study; rows two to four were computed on the MR CLEAN NO-IV development cohort with the same five outer folds as in Table~\ref{tab:cv-performance}, so that the comparison is made within one dataset rather than across publications; the last row is the value reported in the original publication, on a different cohort and for a different endpoint. ICC, intraclass correlation coefficient: of the 107 features the reference pipeline extracts, the 52 with ICC $\geq$ 0.75 are retained and the rest discarded. AUC$_{\text{prob}}$ denotes the AUC computed from the predicted class probabilities, and is the quantity reported for the reference pipeline. Values are medians with the interquartile range in brackets, or means $\pm$ standard deviation across folds, as indicated in the note below. An em dash marks a quantity the row does not have: the pooled row spans combinations with different accuracies, and the published row reports neither a range nor an accuracy.}
\label{tab:reference-comparison}
\smallskip
\input{\tablesdir/table5_reference_comparison}
\end{table}

\smallskip
{\footnotesize
\sloppy
\noindent
\textit{Notes to Table~\ref{tab:reference-comparison}.}
The first row pools 25 fold results, 5 combinations $\times$ 5 outer folds, and its bracketed values are the interquartile range; the second row is a mean $\pm$ standard deviation over the 5 splits; rows three and four are medians over the 5 folds, with no further feature selection applied to the 52 ICC-filtered features.
The last row is quoted from the publication: thrombus radiomics for eTICI $\geq$ 2B, a more permissive endpoint than the eTICI 2C--3 outcome of this article, as a mean over 100 bootstrap samples of 100 patients.
The plain \texttt{AUC} column of the reference pipeline degenerates to exactly 0.50 in 4 of the five splits, because the fitted model assigns the majority class to every patient in those splits and its hard predictions therefore carry no ranking information.
AUC$_{\text{prob}}$, computed from the predicted probabilities, is the only quantity of that pipeline that is comparable with the AUC values reported elsewhere in this article, and it is the one used here.
By median accuracy the reference pipeline and both of our classifiers restricted to its 52 ICC-filtered features stay below the majority-class accuracy of 0.582, although the reference pipeline exceeds it in one individual split.
Rows two to four were recomputed against the eTICI 2C--3 outcome of this article --- the stored label column is named \texttt{TICI\_2B-3}, but its per-fold positive counts are exactly the 82/141 of this study --- and are directly comparable with the first row, whereas the published 0.52 of the last row was obtained for eTICI $\geq$ 2B in a different cohort and differs from them in endpoint and cohort as much as in pipeline.
\NOTEQUERY{the pool of the first row omits one whole-thrombus combination that also discriminated in cross-validation, and the reason is not documented in the source material}
\NOTEQUERY{decide which of the two mutually incompatible sets of numbers in the source data applies before this table is finalized: set A, 0.481 $\pm$ 0.145 on 141 patients against set B, 0.467 $\pm$ 0.099 on 170}
\par}

\clearpage

\begin{table}[htbp]
\centering
\caption{Ensemble analysis: the 15 configurations of discrete, that is binary, AdaBoost \cite{Freund1997} applied to the stored prediction vectors of the 16 patient-level weak classifiers (four patient-level datasets $\times$ four pipelines). Each configuration was fitted on the five outer folds, giving 75 ensemble models. ACC $\varnothing$, mean accuracy of the configuration across the five folds; best weak classifier, the single pipeline in the configuration's pool with the highest mean accuracy, and its value; $\Delta$, the mean accuracy of the ensemble minus that of the best weak classifier. A negative $\Delta$ means the ensemble is worse than the best model it combines. The configuration name states which of the four patient-level datasets form its pool, each contributing all four pipelines: \texttt{type1} pools all four datasets (16 weak classifiers), \texttt{type2} a whole-thrombus dataset with its own centerline counterpart (8), \texttt{type3} a single dataset (4), \texttt{type4} two datasets that are not such a pair (8) and \texttt{type5} all datasets but one (12). \texttt{distrib.}\ abbreviates the distribution-based dataset and \emph{cent.} the centerline region of interest, as in Table~\ref{tab:cv-performance}. The two voxel-based datasets were excluded from every pool because their predictions are defined per voxel.}
\label{tab:ensemble}
\smallskip
\input{\tablesdir/table6_ensemble}
\end{table}

\smallskip
{\footnotesize
\sloppy
\noindent
\textit{Notes to Table~\ref{tab:ensemble}.}
$\Delta$ is negative for all 15 configurations, between $-0.022$ and $-0.121$; no ensemble outperformed the best single pipeline in its own pool.
The worst result belongs to \texttt{type1\_\allowbreak all\_\allowbreak featuresets}, the configuration that pools all feature sets.
Across all 75 fitted ensemble models accuracy averages $0.543 \pm 0.083$ with a range of 0.414--0.714.
The mechanism is documented by the fitted weights: in the all-feature-set configuration, 55 of the 80 boosting rounds (16 rounds $\times$ five folds) received a negative weight $\alpha$, meaning that the weak classifier admitted in that round performed worse than chance on the reweighted data.
The individual pipelines therefore do not make complementary errors, and the boosting assumption of weak classifiers reliably better than chance is not met.
Labels were coded as $\{-1,+1\}$, the weight of a round was $\alpha = \tfrac{1}{2}\ln\frac{1-\varepsilon}{\varepsilon}$ with $\varepsilon$ the error of the admitted classifier on the currently reweighted data, and each weak classifier could be admitted at most once, so a pool of $k$ members yields at most $k$ rounds.
The weak classifiers were not retrained: the procedure is applied to models that had already been fitted, which makes it methodologically a form of stacking rather than boosting proper.
The best weak classifier of \texttt{type3\_\allowbreak global} and of \texttt{type1\_\allowbreak all\_\allowbreak featuresets} is the same pipeline, because the pool of the latter contains the pool of the former.
\par}

\clearpage

\begin{table}[htbp]
\centering
\caption{Pairwise comparison of the pipelines by the corrected resampled $t$-test of Nadeau and Bengio \cite{Nadeau2003}. The upper block compares the four pipelines of this study with one another on the whole-thrombus global-based dataset; the lower block compares them, and our two classifiers restricted to the reference pipeline's own 52 ICC-filtered features, against the reference pipeline of van Voorst et al.\ recomputed on the same cohort. Every model in the table was evaluated on the same 5 outer folds of the same 141 patients, which is what makes the paired test applicable. Because the training sets of $k$-fold cross-validation overlap, the per-fold differences are not independent and the ordinary paired $t$-test underestimates their variance; the correction multiplies the variance by $1/k + n_{\text{test}}/n_{\text{train}}$, here $1/5 + 28.2/112.8 = 0.4500$, which inflates the standard error 1.50-fold at $\mathrm{df} = 4$. Fold sizes are taken from the data, not assumed. Mean $\Delta$AUC is the mean of the 5 per-fold differences in AUC between pipeline A and pipeline B; a positive value favours A. $t$ (uncorrected) is the ordinary paired statistic, $t$ (corrected) the same quantity with the inflated variance. The two-sided critical value at $\alpha = 0.05$ and $\mathrm{df} = 4$ is 2.776. Pipeline names combine selector and classifier as in Table~\ref{tab:cv-performance}.}
\label{tab:nadeau-bengio}
\smallskip
\input{\tablesdir/table7_nadeau_bengio}
\end{table}

\smallskip
{\footnotesize
\sloppy
\noindent
\textit{Notes to Table~\ref{tab:nadeau-bengio}.}
Not one of the 12 comparisons reaches significance: the largest corrected statistic in absolute value is 1.54, against a critical value of 2.776.
The four pipelines of this study are therefore statistically indistinguishable from one another, and none of them is distinguishable from the reference pipeline either, in either direction.
For the reference pipeline the compared quantity is AUC$_{\text{prob}}$, because its plain \texttt{AUC} column degenerates to exactly 0.50 in four of the five splits; see Table~\ref{tab:reference-comparison}.
The pooled row of Table~\ref{tab:reference-comparison} is deliberately absent here: those 5 combinations were picked by their results on these same folds, so testing them against the reference pipeline would be biased in favour of the present study. All four pipelines of the dataset are tested instead, whatever their outcome.
With 5 folds the test has $\mathrm{df} = 4$ and very little power, so these results are an absence of evidence for a difference, not evidence that the pipelines perform identically. The comparison that the data would support best, a DeLong test on the external ROC curves, was not possible because per-patient predicted scores on the Registry were not retained.
\par}

\clearpage


%%     \restoregeometry
%%
%% Cislovani S1--S7 misto 1--7 zaridi pred vlozenim:
%%     \setcounter{table}{0}\renewcommand{\thetable}{S\arabic{table}}
%%
%% Preambule potrebuje: booktabs, array, amsmath, amssymb.
%% Cesta k tabulkam se bere z \tablesdir; vychozi hodnota plati pro preklad
%% z korene projektu.
%% -------------------------------------------------------------------------
\providecommand{\tablesdir}{tables_tom/output}

%% Zaloha, kdyby `manuscriptnotes` nebyl nactveny: vysadi se jen kotva,
%% poznamka se zahodi -- tedy tentyz vysledek jako v rezimu odesilani.
\providecommand{\NOTE}[2][]{#1}
\providecommand{\NOTEWARN}[2][]{#1}
\providecommand{\NOTECHECK}[2][]{#1}
\providecommand{\NOTETODO}[2][]{#1}
\providecommand{\NOTEQUERY}[2][]{#1}

{\large\bfseries Supplementary tables\par}
\bigskip
\NOTEWARN{the captions and notes in this supplement were drafted with AI assistance and are provisional; every one of them has to be rewritten and checked by an author before submission}

\begin{table}[htbp]
\centering
\caption{Patient selection in the two cohorts. MR CLEAN NO-IV was used for model development and cross-validation, the MR CLEAN Registry as the independent external cohort. Entries are numbers of patients; an em dash marks an exclusion step that does not apply to that cohort. eTICI, expanded Thrombolysis in Cerebral Infarction score; the outcome modeled throughout this article is near-complete or complete reperfusion, defined as a final eTICI score of 2C or 3. Slice spacing refers to the through-plane spacing of the non-contrast CT acquisition.}
\label{tab:cohorts}
\smallskip
\input{\tablesdir/table1_cohorts}
\end{table}

\smallskip
{\footnotesize
\sloppy
\noindent
\textit{Notes to Table~\ref{tab:cohorts}.}
The manual review of segmentation quality was carried out on the MR CLEAN NO-IV cohort only.
No segmentation was reviewed or rejected in the MR CLEAN Registry cohort, so the two cohorts differ in segmentation quality control as well as in population.
This is a property of the design and is the single largest asymmetry between the development and the external cohort.
The review was blinded to the eTICI score: it judged the quality and the anatomical placement of the segmentation only, so the 155 exclusions it produced cannot have been informed by the outcome.
Of the 141 patients that remain, 82 (58.2\%) reached eTICI 2C--3; the proportion varies across the five outer folds (0.448, 0.714, 0.536, 0.679, 0.536), so in the first fold the majority class is the absence rather than the presence of near-complete or complete reperfusion.
\NOTETODO{breakdown of how many of the 155 reviewed cases were rejected by which criterion}
\NOTETODO{number and proportion of patients with eTICI 2C--3 reperfusion in the Registry cohort, which the table leaves as an em dash}
\NOTEWARN{size of the Registry cohort: the exclusion cascade gives 1906, the stored identifier list holds 1839, and the denominators reconstructed from the decimal expansion of the reported external accuracies are 1649 (8 models), 1760 (6 models), 1761 (2 models). Neither 1906 nor 1839 is compatible with any of the sixteen transferred models, and the reconstructed values are mutually incompatible, so the external cohort is not the same across feature sets}
\par}

\clearpage

\begin{table}[htbp]
\centering
\caption{Hyperparameter search spaces of the four pipelines. LASSO, $\ell_1$-penalized linear regression on the binary outcome used as a feature selector; RFI, ranking by the gain-based feature importance of a random-forest gradient-boosting model; XGBRF, \texttt{XGBRFClassifier}; LinearSVC, linear support vector classifier with an $\ell_1$ penalty, which acts as selector and classifier simultaneously; NuSVC, $\nu$-support vector classifier with a radial basis function kernel. In the Component column, (S) marks the component acting as feature selector, (C) as classifier, and (S+C) a single object that performs both roles. Scale indicates how Optuna sampled the parameter: \emph{log}, log-uniform over the interval; \emph{int}, uniform over the integers of the interval; \emph{float}, uniform over the interval.}
\label{tab:hyperparameters}
\smallskip
\input{\tablesdir/table2_hyperparameters}
\end{table}

\smallskip
{\footnotesize
\sloppy
\noindent
\textit{Notes to Table~\ref{tab:hyperparameters}.}
$n$ and \texttt{X.shape[1]} denote the number of rows and the number of columns of the training matrix passed to the search.
Rows are patients for the global-based and distribution-based datasets and voxels for the voxel-based datasets, so the same expression yields different intervals across datasets; the number of columns reaches 20660 for the distribution-based datasets.
The minority fraction is the proportion of the smaller outcome class among the training labels of the given fold.
In the RFI--RF pipeline a single parameter dictionary configures both the model whose gain importances rank the features and the classifier fitted on the retained ones, which is why its rows are marked (S+C) rather than (C).
All search spaces were explored with Optuna using the Gaussian-process sampler (\texttt{GPSampler}) with 20 trials per dataset-by-pipeline combination.
The objective maximized was the \emph{mean accuracy over the five inner cross-validation folds}, not AUC; the inner folds were generated by \texttt{StratifiedGroupKFold} grouped by \texttt{patient\_id}.
Settings that were fixed and not searched, one per line so that the long identifiers have somewhere to break:\\
\texttt{Lasso(fit\_intercept=True, max\_iter=5000)}\\
\texttt{LinearSVC(penalty='l1', max\_iter=5000, class\_weight='balanced')}\\
\texttt{NuSVC(kernel='rbf', probability=True, class\_weight='balanced')}\\
\texttt{XGBRFClassifier(random\_state=42)}\\
Class weighting was therefore applied in the two support vector branches only.
The models were seeded with \texttt{random\_state=42}, but \texttt{GPSampler()} was instantiated without a seed, so the sequence of 20 trials is not reproducible.
\par}

\clearpage

\begin{table}[htbp]
\centering
\caption{Cross-validated performance of all 24 dataset-by-pipeline combinations in the MR CLEAN NO-IV cohort ($n = 141$), over five outer folds. Feature set: \emph{global}, eight summary statistics of each radiomic feature over the thrombus, one row per patient; \emph{distribution}, a ten-bin histogram of each radiomic feature over the thrombus, one row per patient; \emph{voxel}, the radiomic feature values themselves, one row per voxel. ROI: \emph{whole}, the dilated thrombus mask; \emph{centerline}, the dilated skeleton of the thrombus mask. Selector (Sel.): LASSO, $\ell_1$-penalized regression; L1, embedded selection by the non-zero coefficients of the classifier itself; RFI, gain-based feature importance of an \texttt{XGBRFClassifier}. Classifier (Class.): XGBRF, \texttt{XGBRFClassifier}; LinSVC, linear support vector classifier with an $\ell_1$ penalty; NuSVC, $\nu$-support vector classifier with a radial basis function kernel. In the body of the table the distribution-based feature set is abbreviated \emph{distrib.}\ and the centerline ROI \emph{cent.} AUC, area under the receiver operating characteristic curve; ACC, accuracy; IQR, interquartile range across the five folds; min--max, the smallest and largest value across the five folds; Med.\ $n$, median number of features retained by the selector; Null AUC, mean AUC of the permutation baseline; $\kappa$, the $\kappa$-type index of the observed AUC relative to that baseline; $p$, permutation $p$-value; $q$, the same $p$-value after Benjamini--Hochberg correction across all 24 combinations.}
\label{tab:cv-performance}
\smallskip
\input{\tablesdir/table3_cv_performance}
\end{table}

\smallskip
{\footnotesize
\sloppy
\noindent
\textit{Notes to Table~\ref{tab:cv-performance}.}
Accuracy must be read against the majority-class accuracy of the development cohort, $82/141 = 0.582$; only 4 of the 24 combinations exceed it by median accuracy, and 86 of the 120 individual fold results (71.7\%) fall below it.
The class ratio itself varies across folds (0.448, 0.714, 0.536, 0.679, 0.536), so in the first fold the majority class is the absence rather than the presence of near-complete or complete reperfusion.
The largest single fold result of the whole study is 0.810, obtained by distribution/whole/RFI/NuSVC in fold 4.
Averaged over the eight whole-thrombus patient-level combinations, AUC per fold is 0.510, 0.661, 0.606, 0.587, 0.623, a spread within a single fold of up to 0.531; between-fold variability therefore exceeds between-method variability.
The permutation null was obtained from 100 label permutations at the patient level per outer fold; with $m = 100$ permutations the smallest attainable $p$-value is $p = (0+1)/(100+1) = 0.0099$, which bounds the resolution of every permutation result in this table.
The Null AUC, $\kappa$, $p$ and $q$ columns were computed from the \emph{mean} AUC across folds, whereas the AUC column reports the \emph{median}; the two summaries are not interchangeable and are shown side by side deliberately.
10 of the 24 combinations reach $p < 0.05$ uncorrected and 6 survive Benjamini--Hochberg correction\NOTETODO{cite Benjamini and Hochberg 1995; the key is missing from BibFile.bib} for AUC, while none survives it for accuracy.
No centerline combination survives correction; across the eight of them the mean AUC ranges from 0.396 to 0.495 against a null of 0.489--0.510, $\kappa$ from $-0.206$ to $-0.014$ and $p$ from 0.5347 to 1.0000.\\
\textsuperscript{a} In these 2 combinations the $\ell_1$ penalty shrank every coefficient to zero in every fold and the code substituted the index of the first column; the retained feature is a fallback, not one the selector identified as informative.
Combinations with a median of one retained feature but no mark reached that median with genuinely non-zero coefficients.\\
\textsuperscript{b} The unit of analysis of the voxel-based datasets is the voxel, not the patient, and their metrics are therefore \emph{not} comparable with the remaining rows.
Across the five folds the voxel/whole datasets contribute 396085 voxel-level predictions (91500, 79327, 75748, 82369, 67141) and the voxel/centerline datasets 8557 (1690, 1571, 1811, 1806, 1679), against 141 patient-level predictions (29, 28, 28, 28, 28) for every other dataset.
The consequence is visible in the permutation null, whose standard deviation is 0.040--0.052 for the voxel-based combinations against 0.063--0.120 for all others, so that the voxel-based $p$-values are computed against a spuriously narrow null and are not comparable with the rest of the table.
3 voxel-based combinations are formally significant for AUC after correction; they must not be read as evidence for a signal.
\NOTETODO{patient-level majority-vote aggregation of the voxel-based predictions, which would make these rows comparable}
\par}

\clearpage

\begin{table}[htbp]
\centering
\caption{External validation in the MR CLEAN Registry cohort. Each pipeline was refitted on all 141 development patients with its tuned hyperparameters and applied once to the Registry cohort. CV, cross-validated median over the five outer folds of the same combination in Table~\ref{tab:cv-performance}; Reg., Registry; ACC, accuracy; $\Delta$AUC, Registry AUC minus cross-validated median AUC, a negative value being a loss of discrimination on transfer; other abbreviations as in Table~\ref{tab:cv-performance}. The voxel-based datasets were not transferred externally and therefore have no row here, which is why the table has 16 rows rather than 24.}
\label{tab:external-validation}
\smallskip
\input{\tablesdir/table4_external_validation}
\end{table}

\smallskip
{\footnotesize
\sloppy
\noindent
\textit{Notes to Table~\ref{tab:external-validation}.}
Across the 16 models the Registry AUC ranged from 0.471 to 0.521 with a median of 0.496, and 5 of them exceeded 0.50.
Registry accuracy ranged from 0.440 to 0.550 with a median of 0.481, and no model reached the majority-class accuracy of the development cohort, 0.582.
The sign of $\Delta$ follows the cross-validated value and nothing else: the 8 combinations whose cross-validated median AUC was above 0.500 all moved downwards on transfer, by $-0.034$ to $-0.190$, and the 8 that were at or below it all moved upwards, by $0.002$ to $0.128$.
The upward movement is therefore a regression to chance of models that had performed below chance internally, not a gain in generalization, and it must not be presented as one.
The only combinations for which the comparison is informative are the 6 that discriminated in cross-validation, that is reached a cross-validated median AUC of at least 0.600; all of them lose between 0.150 and 0.190 of AUC.
\NOTEWARN{$n$ Registry is not stated in the results files and is not the same for every model. Reconstructing the smallest denominator compatible with the decimal expansion of each reported accuracy leaves exactly three cohort sizes: 1649 for eight models, 1760 for six and 1761 for two. Neither the 1906 of the exclusion cascade nor the 1839 identifiers stored with the cohort is compatible with any of the sixteen. Until this is resolved the AUC values in this table are not strictly comparable with one another}
\par}

\clearpage

\begin{table}[htbp]
\centering
\caption{Comparison with the reference pipeline of van Voorst et al. \cite{vanVoorst2022} The first row summarizes the present study; rows two to four were computed on the MR CLEAN NO-IV development cohort with the same five outer folds as in Table~\ref{tab:cv-performance}, so that the comparison is made within one dataset rather than across publications; the last row is the value reported in the original publication, on a different cohort and for a different endpoint. ICC, intraclass correlation coefficient: of the 107 features the reference pipeline extracts, the 52 with ICC $\geq$ 0.75 are retained and the rest discarded. AUC$_{\text{prob}}$ denotes the AUC computed from the predicted class probabilities, and is the quantity reported for the reference pipeline. Values are medians with the interquartile range in brackets, or means $\pm$ standard deviation across folds, as indicated in the note below. An em dash marks a quantity the row does not have: the pooled row spans combinations with different accuracies, and the published row reports neither a range nor an accuracy.}
\label{tab:reference-comparison}
\smallskip
\input{\tablesdir/table5_reference_comparison}
\end{table}

\smallskip
{\footnotesize
\sloppy
\noindent
\textit{Notes to Table~\ref{tab:reference-comparison}.}
The first row pools 25 fold results, 5 combinations $\times$ 5 outer folds, and its bracketed values are the interquartile range; the second row is a mean $\pm$ standard deviation over the 5 splits; rows three and four are medians over the 5 folds, with no further feature selection applied to the 52 ICC-filtered features.
The last row is quoted from the publication: thrombus radiomics for eTICI $\geq$ 2B, a more permissive endpoint than the eTICI 2C--3 outcome of this article, as a mean over 100 bootstrap samples of 100 patients.
The plain \texttt{AUC} column of the reference pipeline degenerates to exactly 0.50 in 4 of the five splits, because the fitted model assigns the majority class to every patient in those splits and its hard predictions therefore carry no ranking information.
AUC$_{\text{prob}}$, computed from the predicted probabilities, is the only quantity of that pipeline that is comparable with the AUC values reported elsewhere in this article, and it is the one used here.
By median accuracy the reference pipeline and both of our classifiers restricted to its 52 ICC-filtered features stay below the majority-class accuracy of 0.582, although the reference pipeline exceeds it in one individual split.
Rows two to four were recomputed against the eTICI 2C--3 outcome of this article --- the stored label column is named \texttt{TICI\_2B-3}, but its per-fold positive counts are exactly the 82/141 of this study --- and are directly comparable with the first row, whereas the published 0.52 of the last row was obtained for eTICI $\geq$ 2B in a different cohort and differs from them in endpoint and cohort as much as in pipeline.
\NOTEQUERY{the pool of the first row omits one whole-thrombus combination that also discriminated in cross-validation, and the reason is not documented in the source material}
\NOTEQUERY{decide which of the two mutually incompatible sets of numbers in the source data applies before this table is finalized: set A, 0.481 $\pm$ 0.145 on 141 patients against set B, 0.467 $\pm$ 0.099 on 170}
\par}

\clearpage

\begin{table}[htbp]
\centering
\caption{Ensemble analysis: the 15 configurations of discrete, that is binary, AdaBoost \cite{Freund1997} applied to the stored prediction vectors of the 16 patient-level weak classifiers (four patient-level datasets $\times$ four pipelines). Each configuration was fitted on the five outer folds, giving 75 ensemble models. ACC $\varnothing$, mean accuracy of the configuration across the five folds; best weak classifier, the single pipeline in the configuration's pool with the highest mean accuracy, and its value; $\Delta$, the mean accuracy of the ensemble minus that of the best weak classifier. A negative $\Delta$ means the ensemble is worse than the best model it combines. The configuration name states which of the four patient-level datasets form its pool, each contributing all four pipelines: \texttt{type1} pools all four datasets (16 weak classifiers), \texttt{type2} a whole-thrombus dataset with its own centerline counterpart (8), \texttt{type3} a single dataset (4), \texttt{type4} two datasets that are not such a pair (8) and \texttt{type5} all datasets but one (12). \texttt{distrib.}\ abbreviates the distribution-based dataset and \emph{cent.} the centerline region of interest, as in Table~\ref{tab:cv-performance}. The two voxel-based datasets were excluded from every pool because their predictions are defined per voxel.}
\label{tab:ensemble}
\smallskip
\input{\tablesdir/table6_ensemble}
\end{table}

\smallskip
{\footnotesize
\sloppy
\noindent
\textit{Notes to Table~\ref{tab:ensemble}.}
$\Delta$ is negative for all 15 configurations, between $-0.022$ and $-0.121$; no ensemble outperformed the best single pipeline in its own pool.
The worst result belongs to \texttt{type1\_\allowbreak all\_\allowbreak featuresets}, the configuration that pools all feature sets.
Across all 75 fitted ensemble models accuracy averages $0.543 \pm 0.083$ with a range of 0.414--0.714.
The mechanism is documented by the fitted weights: in the all-feature-set configuration, 55 of the 80 boosting rounds (16 rounds $\times$ five folds) received a negative weight $\alpha$, meaning that the weak classifier admitted in that round performed worse than chance on the reweighted data.
The individual pipelines therefore do not make complementary errors, and the boosting assumption of weak classifiers reliably better than chance is not met.
Labels were coded as $\{-1,+1\}$, the weight of a round was $\alpha = \tfrac{1}{2}\ln\frac{1-\varepsilon}{\varepsilon}$ with $\varepsilon$ the error of the admitted classifier on the currently reweighted data, and each weak classifier could be admitted at most once, so a pool of $k$ members yields at most $k$ rounds.
The weak classifiers were not retrained: the procedure is applied to models that had already been fitted, which makes it methodologically a form of stacking rather than boosting proper.
The best weak classifier of \texttt{type3\_\allowbreak global} and of \texttt{type1\_\allowbreak all\_\allowbreak featuresets} is the same pipeline, because the pool of the latter contains the pool of the former.
\par}

\clearpage

\begin{table}[htbp]
\centering
\caption{Pairwise comparison of the pipelines by the corrected resampled $t$-test of Nadeau and Bengio \cite{Nadeau2003}. The upper block compares the four pipelines of this study with one another on the whole-thrombus global-based dataset; the lower block compares them, and our two classifiers restricted to the reference pipeline's own 52 ICC-filtered features, against the reference pipeline of van Voorst et al.\ recomputed on the same cohort. Every model in the table was evaluated on the same 5 outer folds of the same 141 patients, which is what makes the paired test applicable. Because the training sets of $k$-fold cross-validation overlap, the per-fold differences are not independent and the ordinary paired $t$-test underestimates their variance; the correction multiplies the variance by $1/k + n_{\text{test}}/n_{\text{train}}$, here $1/5 + 28.2/112.8 = 0.4500$, which inflates the standard error 1.50-fold at $\mathrm{df} = 4$. Fold sizes are taken from the data, not assumed. Mean $\Delta$AUC is the mean of the 5 per-fold differences in AUC between pipeline A and pipeline B; a positive value favours A. $t$ (uncorrected) is the ordinary paired statistic, $t$ (corrected) the same quantity with the inflated variance. The two-sided critical value at $\alpha = 0.05$ and $\mathrm{df} = 4$ is 2.776. Pipeline names combine selector and classifier as in Table~\ref{tab:cv-performance}.}
\label{tab:nadeau-bengio}
\smallskip
\input{\tablesdir/table7_nadeau_bengio}
\end{table}

\smallskip
{\footnotesize
\sloppy
\noindent
\textit{Notes to Table~\ref{tab:nadeau-bengio}.}
Not one of the 12 comparisons reaches significance: the largest corrected statistic in absolute value is 1.54, against a critical value of 2.776.
The four pipelines of this study are therefore statistically indistinguishable from one another, and none of them is distinguishable from the reference pipeline either, in either direction.
For the reference pipeline the compared quantity is AUC$_{\text{prob}}$, because its plain \texttt{AUC} column degenerates to exactly 0.50 in four of the five splits; see Table~\ref{tab:reference-comparison}.
The pooled row of Table~\ref{tab:reference-comparison} is deliberately absent here: those 5 combinations were picked by their results on these same folds, so testing them against the reference pipeline would be biased in favour of the present study. All four pipelines of the dataset are tested instead, whatever their outcome.
With 5 folds the test has $\mathrm{df} = 4$ and very little power, so these results are an absence of evidence for a difference, not evidence that the pipelines perform identically. The comparison that the data would support best, a DeLong test on the external ROC curves, was not possible because per-patient predicted scores on the Registry were not retained.
\par}

\clearpage


%%     \restoregeometry
%%
%% Cislovani S1--S7 misto 1--7 zaridi pred vlozenim:
%%     \setcounter{table}{0}\renewcommand{\thetable}{S\arabic{table}}
%%
%% Preambule potrebuje: booktabs, array, amsmath, amssymb.
%% Cesta k tabulkam se bere z \tablesdir; vychozi hodnota plati pro preklad
%% z korene projektu.
%% -------------------------------------------------------------------------
\providecommand{\tablesdir}{tables_tom/output}

%% Zaloha, kdyby `manuscriptnotes` nebyl nactveny: vysadi se jen kotva,
%% poznamka se zahodi -- tedy tentyz vysledek jako v rezimu odesilani.
\providecommand{\NOTE}[2][]{#1}
\providecommand{\NOTEWARN}[2][]{#1}
\providecommand{\NOTECHECK}[2][]{#1}
\providecommand{\NOTETODO}[2][]{#1}
\providecommand{\NOTEQUERY}[2][]{#1}

{\large\bfseries Supplementary tables\par}
\bigskip
\NOTEWARN{the captions and notes in this supplement were drafted with AI assistance and are provisional; every one of them has to be rewritten and checked by an author before submission}

\begin{table}[htbp]
\centering
\caption{Patient selection in the two cohorts. MR CLEAN NO-IV was used for model development and cross-validation, the MR CLEAN Registry as the independent external cohort. Entries are numbers of patients; an em dash marks an exclusion step that does not apply to that cohort. eTICI, expanded Thrombolysis in Cerebral Infarction score; the outcome modeled throughout this article is near-complete or complete reperfusion, defined as a final eTICI score of 2C or 3. Slice spacing refers to the through-plane spacing of the non-contrast CT acquisition.}
\label{tab:cohorts}
\smallskip
\input{\tablesdir/table1_cohorts}
\end{table}

\smallskip
{\footnotesize
\sloppy
\noindent
\textit{Notes to Table~\ref{tab:cohorts}.}
The manual review of segmentation quality was carried out on the MR CLEAN NO-IV cohort only.
No segmentation was reviewed or rejected in the MR CLEAN Registry cohort, so the two cohorts differ in segmentation quality control as well as in population.
This is a property of the design and is the single largest asymmetry between the development and the external cohort.
The review was blinded to the eTICI score: it judged the quality and the anatomical placement of the segmentation only, so the 155 exclusions it produced cannot have been informed by the outcome.
Of the 141 patients that remain, 82 (58.2\%) reached eTICI 2C--3; the proportion varies across the five outer folds (0.448, 0.714, 0.536, 0.679, 0.536), so in the first fold the majority class is the absence rather than the presence of near-complete or complete reperfusion.
\NOTETODO{breakdown of how many of the 155 reviewed cases were rejected by which criterion}
\NOTETODO{number and proportion of patients with eTICI 2C--3 reperfusion in the Registry cohort, which the table leaves as an em dash}
\NOTEWARN{size of the Registry cohort: the exclusion cascade gives 1906, the stored identifier list holds 1839, and the denominators reconstructed from the decimal expansion of the reported external accuracies are 1649 (8 models), 1760 (6 models), 1761 (2 models). Neither 1906 nor 1839 is compatible with any of the sixteen transferred models, and the reconstructed values are mutually incompatible, so the external cohort is not the same across feature sets}
\par}

\clearpage

\begin{table}[htbp]
\centering
\caption{Hyperparameter search spaces of the four pipelines. LASSO, $\ell_1$-penalized linear regression on the binary outcome used as a feature selector; RFI, ranking by the gain-based feature importance of a random-forest gradient-boosting model; XGBRF, \texttt{XGBRFClassifier}; LinearSVC, linear support vector classifier with an $\ell_1$ penalty, which acts as selector and classifier simultaneously; NuSVC, $\nu$-support vector classifier with a radial basis function kernel. In the Component column, (S) marks the component acting as feature selector, (C) as classifier, and (S+C) a single object that performs both roles. Scale indicates how Optuna sampled the parameter: \emph{log}, log-uniform over the interval; \emph{int}, uniform over the integers of the interval; \emph{float}, uniform over the interval.}
\label{tab:hyperparameters}
\smallskip
\input{\tablesdir/table2_hyperparameters}
\end{table}

\smallskip
{\footnotesize
\sloppy
\noindent
\textit{Notes to Table~\ref{tab:hyperparameters}.}
$n$ and \texttt{X.shape[1]} denote the number of rows and the number of columns of the training matrix passed to the search.
Rows are patients for the global-based and distribution-based datasets and voxels for the voxel-based datasets, so the same expression yields different intervals across datasets; the number of columns reaches 20660 for the distribution-based datasets.
The minority fraction is the proportion of the smaller outcome class among the training labels of the given fold.
In the RFI--RF pipeline a single parameter dictionary configures both the model whose gain importances rank the features and the classifier fitted on the retained ones, which is why its rows are marked (S+C) rather than (C).
All search spaces were explored with Optuna using the Gaussian-process sampler (\texttt{GPSampler}) with 20 trials per dataset-by-pipeline combination.
The objective maximized was the \emph{mean accuracy over the five inner cross-validation folds}, not AUC; the inner folds were generated by \texttt{StratifiedGroupKFold} grouped by \texttt{patient\_id}.
Settings that were fixed and not searched, one per line so that the long identifiers have somewhere to break:\\
\texttt{Lasso(fit\_intercept=True, max\_iter=5000)}\\
\texttt{LinearSVC(penalty='l1', max\_iter=5000, class\_weight='balanced')}\\
\texttt{NuSVC(kernel='rbf', probability=True, class\_weight='balanced')}\\
\texttt{XGBRFClassifier(random\_state=42)}\\
Class weighting was therefore applied in the two support vector branches only.
The models were seeded with \texttt{random\_state=42}, but \texttt{GPSampler()} was instantiated without a seed, so the sequence of 20 trials is not reproducible.
\par}

\clearpage

\begin{table}[htbp]
\centering
\caption{Cross-validated performance of all 24 dataset-by-pipeline combinations in the MR CLEAN NO-IV cohort ($n = 141$), over five outer folds. Feature set: \emph{global}, eight summary statistics of each radiomic feature over the thrombus, one row per patient; \emph{distribution}, a ten-bin histogram of each radiomic feature over the thrombus, one row per patient; \emph{voxel}, the radiomic feature values themselves, one row per voxel. ROI: \emph{whole}, the dilated thrombus mask; \emph{centerline}, the dilated skeleton of the thrombus mask. Selector (Sel.): LASSO, $\ell_1$-penalized regression; L1, embedded selection by the non-zero coefficients of the classifier itself; RFI, gain-based feature importance of an \texttt{XGBRFClassifier}. Classifier (Class.): XGBRF, \texttt{XGBRFClassifier}; LinSVC, linear support vector classifier with an $\ell_1$ penalty; NuSVC, $\nu$-support vector classifier with a radial basis function kernel. In the body of the table the distribution-based feature set is abbreviated \emph{distrib.}\ and the centerline ROI \emph{cent.} AUC, area under the receiver operating characteristic curve; ACC, accuracy; IQR, interquartile range across the five folds; min--max, the smallest and largest value across the five folds; Med.\ $n$, median number of features retained by the selector; Null AUC, mean AUC of the permutation baseline; $\kappa$, the $\kappa$-type index of the observed AUC relative to that baseline; $p$, permutation $p$-value; $q$, the same $p$-value after Benjamini--Hochberg correction across all 24 combinations.}
\label{tab:cv-performance}
\smallskip
\input{\tablesdir/table3_cv_performance}
\end{table}

\smallskip
{\footnotesize
\sloppy
\noindent
\textit{Notes to Table~\ref{tab:cv-performance}.}
Accuracy must be read against the majority-class accuracy of the development cohort, $82/141 = 0.582$; only 4 of the 24 combinations exceed it by median accuracy, and 86 of the 120 individual fold results (71.7\%) fall below it.
The class ratio itself varies across folds (0.448, 0.714, 0.536, 0.679, 0.536), so in the first fold the majority class is the absence rather than the presence of near-complete or complete reperfusion.
The largest single fold result of the whole study is 0.810, obtained by distribution/whole/RFI/NuSVC in fold 4.
Averaged over the eight whole-thrombus patient-level combinations, AUC per fold is 0.510, 0.661, 0.606, 0.587, 0.623, a spread within a single fold of up to 0.531; between-fold variability therefore exceeds between-method variability.
The permutation null was obtained from 100 label permutations at the patient level per outer fold; with $m = 100$ permutations the smallest attainable $p$-value is $p = (0+1)/(100+1) = 0.0099$, which bounds the resolution of every permutation result in this table.
The Null AUC, $\kappa$, $p$ and $q$ columns were computed from the \emph{mean} AUC across folds, whereas the AUC column reports the \emph{median}; the two summaries are not interchangeable and are shown side by side deliberately.
10 of the 24 combinations reach $p < 0.05$ uncorrected and 6 survive Benjamini--Hochberg correction\NOTETODO{cite Benjamini and Hochberg 1995; the key is missing from BibFile.bib} for AUC, while none survives it for accuracy.
No centerline combination survives correction; across the eight of them the mean AUC ranges from 0.396 to 0.495 against a null of 0.489--0.510, $\kappa$ from $-0.206$ to $-0.014$ and $p$ from 0.5347 to 1.0000.\\
\textsuperscript{a} In these 2 combinations the $\ell_1$ penalty shrank every coefficient to zero in every fold and the code substituted the index of the first column; the retained feature is a fallback, not one the selector identified as informative.
Combinations with a median of one retained feature but no mark reached that median with genuinely non-zero coefficients.\\
\textsuperscript{b} The unit of analysis of the voxel-based datasets is the voxel, not the patient, and their metrics are therefore \emph{not} comparable with the remaining rows.
Across the five folds the voxel/whole datasets contribute 396085 voxel-level predictions (91500, 79327, 75748, 82369, 67141) and the voxel/centerline datasets 8557 (1690, 1571, 1811, 1806, 1679), against 141 patient-level predictions (29, 28, 28, 28, 28) for every other dataset.
The consequence is visible in the permutation null, whose standard deviation is 0.040--0.052 for the voxel-based combinations against 0.063--0.120 for all others, so that the voxel-based $p$-values are computed against a spuriously narrow null and are not comparable with the rest of the table.
3 voxel-based combinations are formally significant for AUC after correction; they must not be read as evidence for a signal.
\NOTETODO{patient-level majority-vote aggregation of the voxel-based predictions, which would make these rows comparable}
\par}

\clearpage

\begin{table}[htbp]
\centering
\caption{External validation in the MR CLEAN Registry cohort. Each pipeline was refitted on all 141 development patients with its tuned hyperparameters and applied once to the Registry cohort. CV, cross-validated median over the five outer folds of the same combination in Table~\ref{tab:cv-performance}; Reg., Registry; ACC, accuracy; $\Delta$AUC, Registry AUC minus cross-validated median AUC, a negative value being a loss of discrimination on transfer; other abbreviations as in Table~\ref{tab:cv-performance}. The voxel-based datasets were not transferred externally and therefore have no row here, which is why the table has 16 rows rather than 24.}
\label{tab:external-validation}
\smallskip
\input{\tablesdir/table4_external_validation}
\end{table}

\smallskip
{\footnotesize
\sloppy
\noindent
\textit{Notes to Table~\ref{tab:external-validation}.}
Across the 16 models the Registry AUC ranged from 0.471 to 0.521 with a median of 0.496, and 5 of them exceeded 0.50.
Registry accuracy ranged from 0.440 to 0.550 with a median of 0.481, and no model reached the majority-class accuracy of the development cohort, 0.582.
The sign of $\Delta$ follows the cross-validated value and nothing else: the 8 combinations whose cross-validated median AUC was above 0.500 all moved downwards on transfer, by $-0.034$ to $-0.190$, and the 8 that were at or below it all moved upwards, by $0.002$ to $0.128$.
The upward movement is therefore a regression to chance of models that had performed below chance internally, not a gain in generalization, and it must not be presented as one.
The only combinations for which the comparison is informative are the 6 that discriminated in cross-validation, that is reached a cross-validated median AUC of at least 0.600; all of them lose between 0.150 and 0.190 of AUC.
\NOTEWARN{$n$ Registry is not stated in the results files and is not the same for every model. Reconstructing the smallest denominator compatible with the decimal expansion of each reported accuracy leaves exactly three cohort sizes: 1649 for eight models, 1760 for six and 1761 for two. Neither the 1906 of the exclusion cascade nor the 1839 identifiers stored with the cohort is compatible with any of the sixteen. Until this is resolved the AUC values in this table are not strictly comparable with one another}
\par}

\clearpage

\begin{table}[htbp]
\centering
\caption{Comparison with the reference pipeline of van Voorst et al. \cite{vanVoorst2022} The first row summarizes the present study; rows two to four were computed on the MR CLEAN NO-IV development cohort with the same five outer folds as in Table~\ref{tab:cv-performance}, so that the comparison is made within one dataset rather than across publications; the last row is the value reported in the original publication, on a different cohort and for a different endpoint. ICC, intraclass correlation coefficient: of the 107 features the reference pipeline extracts, the 52 with ICC $\geq$ 0.75 are retained and the rest discarded. AUC$_{\text{prob}}$ denotes the AUC computed from the predicted class probabilities, and is the quantity reported for the reference pipeline. Values are medians with the interquartile range in brackets, or means $\pm$ standard deviation across folds, as indicated in the note below. An em dash marks a quantity the row does not have: the pooled row spans combinations with different accuracies, and the published row reports neither a range nor an accuracy.}
\label{tab:reference-comparison}
\smallskip
\input{\tablesdir/table5_reference_comparison}
\end{table}

\smallskip
{\footnotesize
\sloppy
\noindent
\textit{Notes to Table~\ref{tab:reference-comparison}.}
The first row pools 25 fold results, 5 combinations $\times$ 5 outer folds, and its bracketed values are the interquartile range; the second row is a mean $\pm$ standard deviation over the 5 splits; rows three and four are medians over the 5 folds, with no further feature selection applied to the 52 ICC-filtered features.
The last row is quoted from the publication: thrombus radiomics for eTICI $\geq$ 2B, a more permissive endpoint than the eTICI 2C--3 outcome of this article, as a mean over 100 bootstrap samples of 100 patients.
The plain \texttt{AUC} column of the reference pipeline degenerates to exactly 0.50 in 4 of the five splits, because the fitted model assigns the majority class to every patient in those splits and its hard predictions therefore carry no ranking information.
AUC$_{\text{prob}}$, computed from the predicted probabilities, is the only quantity of that pipeline that is comparable with the AUC values reported elsewhere in this article, and it is the one used here.
By median accuracy the reference pipeline and both of our classifiers restricted to its 52 ICC-filtered features stay below the majority-class accuracy of 0.582, although the reference pipeline exceeds it in one individual split.
Rows two to four were recomputed against the eTICI 2C--3 outcome of this article --- the stored label column is named \texttt{TICI\_2B-3}, but its per-fold positive counts are exactly the 82/141 of this study --- and are directly comparable with the first row, whereas the published 0.52 of the last row was obtained for eTICI $\geq$ 2B in a different cohort and differs from them in endpoint and cohort as much as in pipeline.
\NOTEQUERY{the pool of the first row omits one whole-thrombus combination that also discriminated in cross-validation, and the reason is not documented in the source material}
\NOTEQUERY{decide which of the two mutually incompatible sets of numbers in the source data applies before this table is finalized: set A, 0.481 $\pm$ 0.145 on 141 patients against set B, 0.467 $\pm$ 0.099 on 170}
\par}

\clearpage

\begin{table}[htbp]
\centering
\caption{Ensemble analysis: the 15 configurations of discrete, that is binary, AdaBoost \cite{Freund1997} applied to the stored prediction vectors of the 16 patient-level weak classifiers (four patient-level datasets $\times$ four pipelines). Each configuration was fitted on the five outer folds, giving 75 ensemble models. ACC $\varnothing$, mean accuracy of the configuration across the five folds; best weak classifier, the single pipeline in the configuration's pool with the highest mean accuracy, and its value; $\Delta$, the mean accuracy of the ensemble minus that of the best weak classifier. A negative $\Delta$ means the ensemble is worse than the best model it combines. The configuration name states which of the four patient-level datasets form its pool, each contributing all four pipelines: \texttt{type1} pools all four datasets (16 weak classifiers), \texttt{type2} a whole-thrombus dataset with its own centerline counterpart (8), \texttt{type3} a single dataset (4), \texttt{type4} two datasets that are not such a pair (8) and \texttt{type5} all datasets but one (12). \texttt{distrib.}\ abbreviates the distribution-based dataset and \emph{cent.} the centerline region of interest, as in Table~\ref{tab:cv-performance}. The two voxel-based datasets were excluded from every pool because their predictions are defined per voxel.}
\label{tab:ensemble}
\smallskip
\input{\tablesdir/table6_ensemble}
\end{table}

\smallskip
{\footnotesize
\sloppy
\noindent
\textit{Notes to Table~\ref{tab:ensemble}.}
$\Delta$ is negative for all 15 configurations, between $-0.022$ and $-0.121$; no ensemble outperformed the best single pipeline in its own pool.
The worst result belongs to \texttt{type1\_\allowbreak all\_\allowbreak featuresets}, the configuration that pools all feature sets.
Across all 75 fitted ensemble models accuracy averages $0.543 \pm 0.083$ with a range of 0.414--0.714.
The mechanism is documented by the fitted weights: in the all-feature-set configuration, 55 of the 80 boosting rounds (16 rounds $\times$ five folds) received a negative weight $\alpha$, meaning that the weak classifier admitted in that round performed worse than chance on the reweighted data.
The individual pipelines therefore do not make complementary errors, and the boosting assumption of weak classifiers reliably better than chance is not met.
Labels were coded as $\{-1,+1\}$, the weight of a round was $\alpha = \tfrac{1}{2}\ln\frac{1-\varepsilon}{\varepsilon}$ with $\varepsilon$ the error of the admitted classifier on the currently reweighted data, and each weak classifier could be admitted at most once, so a pool of $k$ members yields at most $k$ rounds.
The weak classifiers were not retrained: the procedure is applied to models that had already been fitted, which makes it methodologically a form of stacking rather than boosting proper.
The best weak classifier of \texttt{type3\_\allowbreak global} and of \texttt{type1\_\allowbreak all\_\allowbreak featuresets} is the same pipeline, because the pool of the latter contains the pool of the former.
\par}

\clearpage

\begin{table}[htbp]
\centering
\caption{Pairwise comparison of the pipelines by the corrected resampled $t$-test of Nadeau and Bengio \cite{Nadeau2003}. The upper block compares the four pipelines of this study with one another on the whole-thrombus global-based dataset; the lower block compares them, and our two classifiers restricted to the reference pipeline's own 52 ICC-filtered features, against the reference pipeline of van Voorst et al.\ recomputed on the same cohort. Every model in the table was evaluated on the same 5 outer folds of the same 141 patients, which is what makes the paired test applicable. Because the training sets of $k$-fold cross-validation overlap, the per-fold differences are not independent and the ordinary paired $t$-test underestimates their variance; the correction multiplies the variance by $1/k + n_{\text{test}}/n_{\text{train}}$, here $1/5 + 28.2/112.8 = 0.4500$, which inflates the standard error 1.50-fold at $\mathrm{df} = 4$. Fold sizes are taken from the data, not assumed. Mean $\Delta$AUC is the mean of the 5 per-fold differences in AUC between pipeline A and pipeline B; a positive value favours A. $t$ (uncorrected) is the ordinary paired statistic, $t$ (corrected) the same quantity with the inflated variance. The two-sided critical value at $\alpha = 0.05$ and $\mathrm{df} = 4$ is 2.776. Pipeline names combine selector and classifier as in Table~\ref{tab:cv-performance}.}
\label{tab:nadeau-bengio}
\smallskip
\input{\tablesdir/table7_nadeau_bengio}
\end{table}

\smallskip
{\footnotesize
\sloppy
\noindent
\textit{Notes to Table~\ref{tab:nadeau-bengio}.}
Not one of the 12 comparisons reaches significance: the largest corrected statistic in absolute value is 1.54, against a critical value of 2.776.
The four pipelines of this study are therefore statistically indistinguishable from one another, and none of them is distinguishable from the reference pipeline either, in either direction.
For the reference pipeline the compared quantity is AUC$_{\text{prob}}$, because its plain \texttt{AUC} column degenerates to exactly 0.50 in four of the five splits; see Table~\ref{tab:reference-comparison}.
The pooled row of Table~\ref{tab:reference-comparison} is deliberately absent here: those 5 combinations were picked by their results on these same folds, so testing them against the reference pipeline would be biased in favour of the present study. All four pipelines of the dataset are tested instead, whatever their outcome.
With 5 folds the test has $\mathrm{df} = 4$ and very little power, so these results are an absence of evidence for a difference, not evidence that the pipelines perform identically. The comparison that the data would support best, a DeLong test on the external ROC curves, was not possible because per-patient predicted scores on the Registry were not retained.
\par}

\clearpage


%%
%% Tabulky se vejdou i do uzke sazby `elsarticle` (390 pt), ale zdrobni.
%% Aby vypadaly jako v samostatnem supplementu, obal se vlozeni geometrii:
%%
%%     \usepackage{geometry}          % v preambuli, za elsarticle
%%     ...
%%     \clearpage
%%     \newgeometry{margin=15mm}
%%     %% -------------------------------------------------------------------------
%% GENEROVANO: make_all.py -- needitovat rucne.
%%
%% Popisky, tabulky a poznamky bez preambule. Do clanku staci jeden radek:
%%
%%     %% -------------------------------------------------------------------------
%% GENEROVANO: make_all.py -- needitovat rucne.
%%
%% Popisky, tabulky a poznamky bez preambule. Do clanku staci jeden radek:
%%
%%     %% -------------------------------------------------------------------------
%% GENEROVANO: make_all.py -- needitovat rucne.
%%
%% Popisky, tabulky a poznamky bez preambule. Do clanku staci jeden radek:
%%
%%     \input{tables_tom/output/supplement_body}
%%
%% Tabulky se vejdou i do uzke sazby `elsarticle` (390 pt), ale zdrobni.
%% Aby vypadaly jako v samostatnem supplementu, obal se vlozeni geometrii:
%%
%%     \usepackage{geometry}          % v preambuli, za elsarticle
%%     ...
%%     \clearpage
%%     \newgeometry{margin=15mm}
%%     \input{tables_tom/output/supplement_body}
%%     \restoregeometry
%%
%% Cislovani S1--S7 misto 1--7 zaridi pred vlozenim:
%%     \setcounter{table}{0}\renewcommand{\thetable}{S\arabic{table}}
%%
%% Preambule potrebuje: booktabs, array, amsmath, amssymb.
%% Cesta k tabulkam se bere z \tablesdir; vychozi hodnota plati pro preklad
%% z korene projektu.
%% -------------------------------------------------------------------------
\providecommand{\tablesdir}{tables_tom/output}

%% Zaloha, kdyby `manuscriptnotes` nebyl nactveny: vysadi se jen kotva,
%% poznamka se zahodi -- tedy tentyz vysledek jako v rezimu odesilani.
\providecommand{\NOTE}[2][]{#1}
\providecommand{\NOTEWARN}[2][]{#1}
\providecommand{\NOTECHECK}[2][]{#1}
\providecommand{\NOTETODO}[2][]{#1}
\providecommand{\NOTEQUERY}[2][]{#1}

{\large\bfseries Supplementary tables\par}
\bigskip
\NOTEWARN{the captions and notes in this supplement were drafted with AI assistance and are provisional; every one of them has to be rewritten and checked by an author before submission}

\begin{table}[htbp]
\centering
\caption{Patient selection in the two cohorts. MR CLEAN NO-IV was used for model development and cross-validation, the MR CLEAN Registry as the independent external cohort. Entries are numbers of patients; an em dash marks an exclusion step that does not apply to that cohort. eTICI, expanded Thrombolysis in Cerebral Infarction score; the outcome modeled throughout this article is near-complete or complete reperfusion, defined as a final eTICI score of 2C or 3. Slice spacing refers to the through-plane spacing of the non-contrast CT acquisition.}
\label{tab:cohorts}
\smallskip
\input{\tablesdir/table1_cohorts}
\end{table}

\smallskip
{\footnotesize
\sloppy
\noindent
\textit{Notes to Table~\ref{tab:cohorts}.}
The manual review of segmentation quality was carried out on the MR CLEAN NO-IV cohort only.
No segmentation was reviewed or rejected in the MR CLEAN Registry cohort, so the two cohorts differ in segmentation quality control as well as in population.
This is a property of the design and is the single largest asymmetry between the development and the external cohort.
The review was blinded to the eTICI score: it judged the quality and the anatomical placement of the segmentation only, so the 155 exclusions it produced cannot have been informed by the outcome.
Of the 141 patients that remain, 82 (58.2\%) reached eTICI 2C--3; the proportion varies across the five outer folds (0.448, 0.714, 0.536, 0.679, 0.536), so in the first fold the majority class is the absence rather than the presence of near-complete or complete reperfusion.
\NOTETODO{breakdown of how many of the 155 reviewed cases were rejected by which criterion}
\NOTETODO{number and proportion of patients with eTICI 2C--3 reperfusion in the Registry cohort, which the table leaves as an em dash}
\NOTEWARN{size of the Registry cohort: the exclusion cascade gives 1906, the stored identifier list holds 1839, and the denominators reconstructed from the decimal expansion of the reported external accuracies are 1649 (8 models), 1760 (6 models), 1761 (2 models). Neither 1906 nor 1839 is compatible with any of the sixteen transferred models, and the reconstructed values are mutually incompatible, so the external cohort is not the same across feature sets}
\par}

\clearpage

\begin{table}[htbp]
\centering
\caption{Hyperparameter search spaces of the four pipelines. LASSO, $\ell_1$-penalized linear regression on the binary outcome used as a feature selector; RFI, ranking by the gain-based feature importance of a random-forest gradient-boosting model; XGBRF, \texttt{XGBRFClassifier}; LinearSVC, linear support vector classifier with an $\ell_1$ penalty, which acts as selector and classifier simultaneously; NuSVC, $\nu$-support vector classifier with a radial basis function kernel. In the Component column, (S) marks the component acting as feature selector, (C) as classifier, and (S+C) a single object that performs both roles. Scale indicates how Optuna sampled the parameter: \emph{log}, log-uniform over the interval; \emph{int}, uniform over the integers of the interval; \emph{float}, uniform over the interval.}
\label{tab:hyperparameters}
\smallskip
\input{\tablesdir/table2_hyperparameters}
\end{table}

\smallskip
{\footnotesize
\sloppy
\noindent
\textit{Notes to Table~\ref{tab:hyperparameters}.}
$n$ and \texttt{X.shape[1]} denote the number of rows and the number of columns of the training matrix passed to the search.
Rows are patients for the global-based and distribution-based datasets and voxels for the voxel-based datasets, so the same expression yields different intervals across datasets; the number of columns reaches 20660 for the distribution-based datasets.
The minority fraction is the proportion of the smaller outcome class among the training labels of the given fold.
In the RFI--RF pipeline a single parameter dictionary configures both the model whose gain importances rank the features and the classifier fitted on the retained ones, which is why its rows are marked (S+C) rather than (C).
All search spaces were explored with Optuna using the Gaussian-process sampler (\texttt{GPSampler}) with 20 trials per dataset-by-pipeline combination.
The objective maximized was the \emph{mean accuracy over the five inner cross-validation folds}, not AUC; the inner folds were generated by \texttt{StratifiedGroupKFold} grouped by \texttt{patient\_id}.
Settings that were fixed and not searched, one per line so that the long identifiers have somewhere to break:\\
\texttt{Lasso(fit\_intercept=True, max\_iter=5000)}\\
\texttt{LinearSVC(penalty='l1', max\_iter=5000, class\_weight='balanced')}\\
\texttt{NuSVC(kernel='rbf', probability=True, class\_weight='balanced')}\\
\texttt{XGBRFClassifier(random\_state=42)}\\
Class weighting was therefore applied in the two support vector branches only.
The models were seeded with \texttt{random\_state=42}, but \texttt{GPSampler()} was instantiated without a seed, so the sequence of 20 trials is not reproducible.
\par}

\clearpage

\begin{table}[htbp]
\centering
\caption{Cross-validated performance of all 24 dataset-by-pipeline combinations in the MR CLEAN NO-IV cohort ($n = 141$), over five outer folds. Feature set: \emph{global}, eight summary statistics of each radiomic feature over the thrombus, one row per patient; \emph{distribution}, a ten-bin histogram of each radiomic feature over the thrombus, one row per patient; \emph{voxel}, the radiomic feature values themselves, one row per voxel. ROI: \emph{whole}, the dilated thrombus mask; \emph{centerline}, the dilated skeleton of the thrombus mask. Selector (Sel.): LASSO, $\ell_1$-penalized regression; L1, embedded selection by the non-zero coefficients of the classifier itself; RFI, gain-based feature importance of an \texttt{XGBRFClassifier}. Classifier (Class.): XGBRF, \texttt{XGBRFClassifier}; LinSVC, linear support vector classifier with an $\ell_1$ penalty; NuSVC, $\nu$-support vector classifier with a radial basis function kernel. In the body of the table the distribution-based feature set is abbreviated \emph{distrib.}\ and the centerline ROI \emph{cent.} AUC, area under the receiver operating characteristic curve; ACC, accuracy; IQR, interquartile range across the five folds; min--max, the smallest and largest value across the five folds; Med.\ $n$, median number of features retained by the selector; Null AUC, mean AUC of the permutation baseline; $\kappa$, the $\kappa$-type index of the observed AUC relative to that baseline; $p$, permutation $p$-value; $q$, the same $p$-value after Benjamini--Hochberg correction across all 24 combinations.}
\label{tab:cv-performance}
\smallskip
\input{\tablesdir/table3_cv_performance}
\end{table}

\smallskip
{\footnotesize
\sloppy
\noindent
\textit{Notes to Table~\ref{tab:cv-performance}.}
Accuracy must be read against the majority-class accuracy of the development cohort, $82/141 = 0.582$; only 4 of the 24 combinations exceed it by median accuracy, and 86 of the 120 individual fold results (71.7\%) fall below it.
The class ratio itself varies across folds (0.448, 0.714, 0.536, 0.679, 0.536), so in the first fold the majority class is the absence rather than the presence of near-complete or complete reperfusion.
The largest single fold result of the whole study is 0.810, obtained by distribution/whole/RFI/NuSVC in fold 4.
Averaged over the eight whole-thrombus patient-level combinations, AUC per fold is 0.510, 0.661, 0.606, 0.587, 0.623, a spread within a single fold of up to 0.531; between-fold variability therefore exceeds between-method variability.
The permutation null was obtained from 100 label permutations at the patient level per outer fold; with $m = 100$ permutations the smallest attainable $p$-value is $p = (0+1)/(100+1) = 0.0099$, which bounds the resolution of every permutation result in this table.
The Null AUC, $\kappa$, $p$ and $q$ columns were computed from the \emph{mean} AUC across folds, whereas the AUC column reports the \emph{median}; the two summaries are not interchangeable and are shown side by side deliberately.
10 of the 24 combinations reach $p < 0.05$ uncorrected and 6 survive Benjamini--Hochberg correction\NOTETODO{cite Benjamini and Hochberg 1995; the key is missing from BibFile.bib} for AUC, while none survives it for accuracy.
No centerline combination survives correction; across the eight of them the mean AUC ranges from 0.396 to 0.495 against a null of 0.489--0.510, $\kappa$ from $-0.206$ to $-0.014$ and $p$ from 0.5347 to 1.0000.\\
\textsuperscript{a} In these 2 combinations the $\ell_1$ penalty shrank every coefficient to zero in every fold and the code substituted the index of the first column; the retained feature is a fallback, not one the selector identified as informative.
Combinations with a median of one retained feature but no mark reached that median with genuinely non-zero coefficients.\\
\textsuperscript{b} The unit of analysis of the voxel-based datasets is the voxel, not the patient, and their metrics are therefore \emph{not} comparable with the remaining rows.
Across the five folds the voxel/whole datasets contribute 396085 voxel-level predictions (91500, 79327, 75748, 82369, 67141) and the voxel/centerline datasets 8557 (1690, 1571, 1811, 1806, 1679), against 141 patient-level predictions (29, 28, 28, 28, 28) for every other dataset.
The consequence is visible in the permutation null, whose standard deviation is 0.040--0.052 for the voxel-based combinations against 0.063--0.120 for all others, so that the voxel-based $p$-values are computed against a spuriously narrow null and are not comparable with the rest of the table.
3 voxel-based combinations are formally significant for AUC after correction; they must not be read as evidence for a signal.
\NOTETODO{patient-level majority-vote aggregation of the voxel-based predictions, which would make these rows comparable}
\par}

\clearpage

\begin{table}[htbp]
\centering
\caption{External validation in the MR CLEAN Registry cohort. Each pipeline was refitted on all 141 development patients with its tuned hyperparameters and applied once to the Registry cohort. CV, cross-validated median over the five outer folds of the same combination in Table~\ref{tab:cv-performance}; Reg., Registry; ACC, accuracy; $\Delta$AUC, Registry AUC minus cross-validated median AUC, a negative value being a loss of discrimination on transfer; other abbreviations as in Table~\ref{tab:cv-performance}. The voxel-based datasets were not transferred externally and therefore have no row here, which is why the table has 16 rows rather than 24.}
\label{tab:external-validation}
\smallskip
\input{\tablesdir/table4_external_validation}
\end{table}

\smallskip
{\footnotesize
\sloppy
\noindent
\textit{Notes to Table~\ref{tab:external-validation}.}
Across the 16 models the Registry AUC ranged from 0.471 to 0.521 with a median of 0.496, and 5 of them exceeded 0.50.
Registry accuracy ranged from 0.440 to 0.550 with a median of 0.481, and no model reached the majority-class accuracy of the development cohort, 0.582.
The sign of $\Delta$ follows the cross-validated value and nothing else: the 8 combinations whose cross-validated median AUC was above 0.500 all moved downwards on transfer, by $-0.034$ to $-0.190$, and the 8 that were at or below it all moved upwards, by $0.002$ to $0.128$.
The upward movement is therefore a regression to chance of models that had performed below chance internally, not a gain in generalization, and it must not be presented as one.
The only combinations for which the comparison is informative are the 6 that discriminated in cross-validation, that is reached a cross-validated median AUC of at least 0.600; all of them lose between 0.150 and 0.190 of AUC.
\NOTEWARN{$n$ Registry is not stated in the results files and is not the same for every model. Reconstructing the smallest denominator compatible with the decimal expansion of each reported accuracy leaves exactly three cohort sizes: 1649 for eight models, 1760 for six and 1761 for two. Neither the 1906 of the exclusion cascade nor the 1839 identifiers stored with the cohort is compatible with any of the sixteen. Until this is resolved the AUC values in this table are not strictly comparable with one another}
\par}

\clearpage

\begin{table}[htbp]
\centering
\caption{Comparison with the reference pipeline of van Voorst et al. \cite{vanVoorst2022} The first row summarizes the present study; rows two to four were computed on the MR CLEAN NO-IV development cohort with the same five outer folds as in Table~\ref{tab:cv-performance}, so that the comparison is made within one dataset rather than across publications; the last row is the value reported in the original publication, on a different cohort and for a different endpoint. ICC, intraclass correlation coefficient: of the 107 features the reference pipeline extracts, the 52 with ICC $\geq$ 0.75 are retained and the rest discarded. AUC$_{\text{prob}}$ denotes the AUC computed from the predicted class probabilities, and is the quantity reported for the reference pipeline. Values are medians with the interquartile range in brackets, or means $\pm$ standard deviation across folds, as indicated in the note below. An em dash marks a quantity the row does not have: the pooled row spans combinations with different accuracies, and the published row reports neither a range nor an accuracy.}
\label{tab:reference-comparison}
\smallskip
\input{\tablesdir/table5_reference_comparison}
\end{table}

\smallskip
{\footnotesize
\sloppy
\noindent
\textit{Notes to Table~\ref{tab:reference-comparison}.}
The first row pools 25 fold results, 5 combinations $\times$ 5 outer folds, and its bracketed values are the interquartile range; the second row is a mean $\pm$ standard deviation over the 5 splits; rows three and four are medians over the 5 folds, with no further feature selection applied to the 52 ICC-filtered features.
The last row is quoted from the publication: thrombus radiomics for eTICI $\geq$ 2B, a more permissive endpoint than the eTICI 2C--3 outcome of this article, as a mean over 100 bootstrap samples of 100 patients.
The plain \texttt{AUC} column of the reference pipeline degenerates to exactly 0.50 in 4 of the five splits, because the fitted model assigns the majority class to every patient in those splits and its hard predictions therefore carry no ranking information.
AUC$_{\text{prob}}$, computed from the predicted probabilities, is the only quantity of that pipeline that is comparable with the AUC values reported elsewhere in this article, and it is the one used here.
By median accuracy the reference pipeline and both of our classifiers restricted to its 52 ICC-filtered features stay below the majority-class accuracy of 0.582, although the reference pipeline exceeds it in one individual split.
Rows two to four were recomputed against the eTICI 2C--3 outcome of this article --- the stored label column is named \texttt{TICI\_2B-3}, but its per-fold positive counts are exactly the 82/141 of this study --- and are directly comparable with the first row, whereas the published 0.52 of the last row was obtained for eTICI $\geq$ 2B in a different cohort and differs from them in endpoint and cohort as much as in pipeline.
\NOTEQUERY{the pool of the first row omits one whole-thrombus combination that also discriminated in cross-validation, and the reason is not documented in the source material}
\NOTEQUERY{decide which of the two mutually incompatible sets of numbers in the source data applies before this table is finalized: set A, 0.481 $\pm$ 0.145 on 141 patients against set B, 0.467 $\pm$ 0.099 on 170}
\par}

\clearpage

\begin{table}[htbp]
\centering
\caption{Ensemble analysis: the 15 configurations of discrete, that is binary, AdaBoost \cite{Freund1997} applied to the stored prediction vectors of the 16 patient-level weak classifiers (four patient-level datasets $\times$ four pipelines). Each configuration was fitted on the five outer folds, giving 75 ensemble models. ACC $\varnothing$, mean accuracy of the configuration across the five folds; best weak classifier, the single pipeline in the configuration's pool with the highest mean accuracy, and its value; $\Delta$, the mean accuracy of the ensemble minus that of the best weak classifier. A negative $\Delta$ means the ensemble is worse than the best model it combines. The configuration name states which of the four patient-level datasets form its pool, each contributing all four pipelines: \texttt{type1} pools all four datasets (16 weak classifiers), \texttt{type2} a whole-thrombus dataset with its own centerline counterpart (8), \texttt{type3} a single dataset (4), \texttt{type4} two datasets that are not such a pair (8) and \texttt{type5} all datasets but one (12). \texttt{distrib.}\ abbreviates the distribution-based dataset and \emph{cent.} the centerline region of interest, as in Table~\ref{tab:cv-performance}. The two voxel-based datasets were excluded from every pool because their predictions are defined per voxel.}
\label{tab:ensemble}
\smallskip
\input{\tablesdir/table6_ensemble}
\end{table}

\smallskip
{\footnotesize
\sloppy
\noindent
\textit{Notes to Table~\ref{tab:ensemble}.}
$\Delta$ is negative for all 15 configurations, between $-0.022$ and $-0.121$; no ensemble outperformed the best single pipeline in its own pool.
The worst result belongs to \texttt{type1\_\allowbreak all\_\allowbreak featuresets}, the configuration that pools all feature sets.
Across all 75 fitted ensemble models accuracy averages $0.543 \pm 0.083$ with a range of 0.414--0.714.
The mechanism is documented by the fitted weights: in the all-feature-set configuration, 55 of the 80 boosting rounds (16 rounds $\times$ five folds) received a negative weight $\alpha$, meaning that the weak classifier admitted in that round performed worse than chance on the reweighted data.
The individual pipelines therefore do not make complementary errors, and the boosting assumption of weak classifiers reliably better than chance is not met.
Labels were coded as $\{-1,+1\}$, the weight of a round was $\alpha = \tfrac{1}{2}\ln\frac{1-\varepsilon}{\varepsilon}$ with $\varepsilon$ the error of the admitted classifier on the currently reweighted data, and each weak classifier could be admitted at most once, so a pool of $k$ members yields at most $k$ rounds.
The weak classifiers were not retrained: the procedure is applied to models that had already been fitted, which makes it methodologically a form of stacking rather than boosting proper.
The best weak classifier of \texttt{type3\_\allowbreak global} and of \texttt{type1\_\allowbreak all\_\allowbreak featuresets} is the same pipeline, because the pool of the latter contains the pool of the former.
\par}

\clearpage

\begin{table}[htbp]
\centering
\caption{Pairwise comparison of the pipelines by the corrected resampled $t$-test of Nadeau and Bengio \cite{Nadeau2003}. The upper block compares the four pipelines of this study with one another on the whole-thrombus global-based dataset; the lower block compares them, and our two classifiers restricted to the reference pipeline's own 52 ICC-filtered features, against the reference pipeline of van Voorst et al.\ recomputed on the same cohort. Every model in the table was evaluated on the same 5 outer folds of the same 141 patients, which is what makes the paired test applicable. Because the training sets of $k$-fold cross-validation overlap, the per-fold differences are not independent and the ordinary paired $t$-test underestimates their variance; the correction multiplies the variance by $1/k + n_{\text{test}}/n_{\text{train}}$, here $1/5 + 28.2/112.8 = 0.4500$, which inflates the standard error 1.50-fold at $\mathrm{df} = 4$. Fold sizes are taken from the data, not assumed. Mean $\Delta$AUC is the mean of the 5 per-fold differences in AUC between pipeline A and pipeline B; a positive value favours A. $t$ (uncorrected) is the ordinary paired statistic, $t$ (corrected) the same quantity with the inflated variance. The two-sided critical value at $\alpha = 0.05$ and $\mathrm{df} = 4$ is 2.776. Pipeline names combine selector and classifier as in Table~\ref{tab:cv-performance}.}
\label{tab:nadeau-bengio}
\smallskip
\input{\tablesdir/table7_nadeau_bengio}
\end{table}

\smallskip
{\footnotesize
\sloppy
\noindent
\textit{Notes to Table~\ref{tab:nadeau-bengio}.}
Not one of the 12 comparisons reaches significance: the largest corrected statistic in absolute value is 1.54, against a critical value of 2.776.
The four pipelines of this study are therefore statistically indistinguishable from one another, and none of them is distinguishable from the reference pipeline either, in either direction.
For the reference pipeline the compared quantity is AUC$_{\text{prob}}$, because its plain \texttt{AUC} column degenerates to exactly 0.50 in four of the five splits; see Table~\ref{tab:reference-comparison}.
The pooled row of Table~\ref{tab:reference-comparison} is deliberately absent here: those 5 combinations were picked by their results on these same folds, so testing them against the reference pipeline would be biased in favour of the present study. All four pipelines of the dataset are tested instead, whatever their outcome.
With 5 folds the test has $\mathrm{df} = 4$ and very little power, so these results are an absence of evidence for a difference, not evidence that the pipelines perform identically. The comparison that the data would support best, a DeLong test on the external ROC curves, was not possible because per-patient predicted scores on the Registry were not retained.
\par}

\clearpage


%%
%% Tabulky se vejdou i do uzke sazby `elsarticle` (390 pt), ale zdrobni.
%% Aby vypadaly jako v samostatnem supplementu, obal se vlozeni geometrii:
%%
%%     \usepackage{geometry}          % v preambuli, za elsarticle
%%     ...
%%     \clearpage
%%     \newgeometry{margin=15mm}
%%     %% -------------------------------------------------------------------------
%% GENEROVANO: make_all.py -- needitovat rucne.
%%
%% Popisky, tabulky a poznamky bez preambule. Do clanku staci jeden radek:
%%
%%     \input{tables_tom/output/supplement_body}
%%
%% Tabulky se vejdou i do uzke sazby `elsarticle` (390 pt), ale zdrobni.
%% Aby vypadaly jako v samostatnem supplementu, obal se vlozeni geometrii:
%%
%%     \usepackage{geometry}          % v preambuli, za elsarticle
%%     ...
%%     \clearpage
%%     \newgeometry{margin=15mm}
%%     \input{tables_tom/output/supplement_body}
%%     \restoregeometry
%%
%% Cislovani S1--S7 misto 1--7 zaridi pred vlozenim:
%%     \setcounter{table}{0}\renewcommand{\thetable}{S\arabic{table}}
%%
%% Preambule potrebuje: booktabs, array, amsmath, amssymb.
%% Cesta k tabulkam se bere z \tablesdir; vychozi hodnota plati pro preklad
%% z korene projektu.
%% -------------------------------------------------------------------------
\providecommand{\tablesdir}{tables_tom/output}

%% Zaloha, kdyby `manuscriptnotes` nebyl nactveny: vysadi se jen kotva,
%% poznamka se zahodi -- tedy tentyz vysledek jako v rezimu odesilani.
\providecommand{\NOTE}[2][]{#1}
\providecommand{\NOTEWARN}[2][]{#1}
\providecommand{\NOTECHECK}[2][]{#1}
\providecommand{\NOTETODO}[2][]{#1}
\providecommand{\NOTEQUERY}[2][]{#1}

{\large\bfseries Supplementary tables\par}
\bigskip
\NOTEWARN{the captions and notes in this supplement were drafted with AI assistance and are provisional; every one of them has to be rewritten and checked by an author before submission}

\begin{table}[htbp]
\centering
\caption{Patient selection in the two cohorts. MR CLEAN NO-IV was used for model development and cross-validation, the MR CLEAN Registry as the independent external cohort. Entries are numbers of patients; an em dash marks an exclusion step that does not apply to that cohort. eTICI, expanded Thrombolysis in Cerebral Infarction score; the outcome modeled throughout this article is near-complete or complete reperfusion, defined as a final eTICI score of 2C or 3. Slice spacing refers to the through-plane spacing of the non-contrast CT acquisition.}
\label{tab:cohorts}
\smallskip
\input{\tablesdir/table1_cohorts}
\end{table}

\smallskip
{\footnotesize
\sloppy
\noindent
\textit{Notes to Table~\ref{tab:cohorts}.}
The manual review of segmentation quality was carried out on the MR CLEAN NO-IV cohort only.
No segmentation was reviewed or rejected in the MR CLEAN Registry cohort, so the two cohorts differ in segmentation quality control as well as in population.
This is a property of the design and is the single largest asymmetry between the development and the external cohort.
The review was blinded to the eTICI score: it judged the quality and the anatomical placement of the segmentation only, so the 155 exclusions it produced cannot have been informed by the outcome.
Of the 141 patients that remain, 82 (58.2\%) reached eTICI 2C--3; the proportion varies across the five outer folds (0.448, 0.714, 0.536, 0.679, 0.536), so in the first fold the majority class is the absence rather than the presence of near-complete or complete reperfusion.
\NOTETODO{breakdown of how many of the 155 reviewed cases were rejected by which criterion}
\NOTETODO{number and proportion of patients with eTICI 2C--3 reperfusion in the Registry cohort, which the table leaves as an em dash}
\NOTEWARN{size of the Registry cohort: the exclusion cascade gives 1906, the stored identifier list holds 1839, and the denominators reconstructed from the decimal expansion of the reported external accuracies are 1649 (8 models), 1760 (6 models), 1761 (2 models). Neither 1906 nor 1839 is compatible with any of the sixteen transferred models, and the reconstructed values are mutually incompatible, so the external cohort is not the same across feature sets}
\par}

\clearpage

\begin{table}[htbp]
\centering
\caption{Hyperparameter search spaces of the four pipelines. LASSO, $\ell_1$-penalized linear regression on the binary outcome used as a feature selector; RFI, ranking by the gain-based feature importance of a random-forest gradient-boosting model; XGBRF, \texttt{XGBRFClassifier}; LinearSVC, linear support vector classifier with an $\ell_1$ penalty, which acts as selector and classifier simultaneously; NuSVC, $\nu$-support vector classifier with a radial basis function kernel. In the Component column, (S) marks the component acting as feature selector, (C) as classifier, and (S+C) a single object that performs both roles. Scale indicates how Optuna sampled the parameter: \emph{log}, log-uniform over the interval; \emph{int}, uniform over the integers of the interval; \emph{float}, uniform over the interval.}
\label{tab:hyperparameters}
\smallskip
\input{\tablesdir/table2_hyperparameters}
\end{table}

\smallskip
{\footnotesize
\sloppy
\noindent
\textit{Notes to Table~\ref{tab:hyperparameters}.}
$n$ and \texttt{X.shape[1]} denote the number of rows and the number of columns of the training matrix passed to the search.
Rows are patients for the global-based and distribution-based datasets and voxels for the voxel-based datasets, so the same expression yields different intervals across datasets; the number of columns reaches 20660 for the distribution-based datasets.
The minority fraction is the proportion of the smaller outcome class among the training labels of the given fold.
In the RFI--RF pipeline a single parameter dictionary configures both the model whose gain importances rank the features and the classifier fitted on the retained ones, which is why its rows are marked (S+C) rather than (C).
All search spaces were explored with Optuna using the Gaussian-process sampler (\texttt{GPSampler}) with 20 trials per dataset-by-pipeline combination.
The objective maximized was the \emph{mean accuracy over the five inner cross-validation folds}, not AUC; the inner folds were generated by \texttt{StratifiedGroupKFold} grouped by \texttt{patient\_id}.
Settings that were fixed and not searched, one per line so that the long identifiers have somewhere to break:\\
\texttt{Lasso(fit\_intercept=True, max\_iter=5000)}\\
\texttt{LinearSVC(penalty='l1', max\_iter=5000, class\_weight='balanced')}\\
\texttt{NuSVC(kernel='rbf', probability=True, class\_weight='balanced')}\\
\texttt{XGBRFClassifier(random\_state=42)}\\
Class weighting was therefore applied in the two support vector branches only.
The models were seeded with \texttt{random\_state=42}, but \texttt{GPSampler()} was instantiated without a seed, so the sequence of 20 trials is not reproducible.
\par}

\clearpage

\begin{table}[htbp]
\centering
\caption{Cross-validated performance of all 24 dataset-by-pipeline combinations in the MR CLEAN NO-IV cohort ($n = 141$), over five outer folds. Feature set: \emph{global}, eight summary statistics of each radiomic feature over the thrombus, one row per patient; \emph{distribution}, a ten-bin histogram of each radiomic feature over the thrombus, one row per patient; \emph{voxel}, the radiomic feature values themselves, one row per voxel. ROI: \emph{whole}, the dilated thrombus mask; \emph{centerline}, the dilated skeleton of the thrombus mask. Selector (Sel.): LASSO, $\ell_1$-penalized regression; L1, embedded selection by the non-zero coefficients of the classifier itself; RFI, gain-based feature importance of an \texttt{XGBRFClassifier}. Classifier (Class.): XGBRF, \texttt{XGBRFClassifier}; LinSVC, linear support vector classifier with an $\ell_1$ penalty; NuSVC, $\nu$-support vector classifier with a radial basis function kernel. In the body of the table the distribution-based feature set is abbreviated \emph{distrib.}\ and the centerline ROI \emph{cent.} AUC, area under the receiver operating characteristic curve; ACC, accuracy; IQR, interquartile range across the five folds; min--max, the smallest and largest value across the five folds; Med.\ $n$, median number of features retained by the selector; Null AUC, mean AUC of the permutation baseline; $\kappa$, the $\kappa$-type index of the observed AUC relative to that baseline; $p$, permutation $p$-value; $q$, the same $p$-value after Benjamini--Hochberg correction across all 24 combinations.}
\label{tab:cv-performance}
\smallskip
\input{\tablesdir/table3_cv_performance}
\end{table}

\smallskip
{\footnotesize
\sloppy
\noindent
\textit{Notes to Table~\ref{tab:cv-performance}.}
Accuracy must be read against the majority-class accuracy of the development cohort, $82/141 = 0.582$; only 4 of the 24 combinations exceed it by median accuracy, and 86 of the 120 individual fold results (71.7\%) fall below it.
The class ratio itself varies across folds (0.448, 0.714, 0.536, 0.679, 0.536), so in the first fold the majority class is the absence rather than the presence of near-complete or complete reperfusion.
The largest single fold result of the whole study is 0.810, obtained by distribution/whole/RFI/NuSVC in fold 4.
Averaged over the eight whole-thrombus patient-level combinations, AUC per fold is 0.510, 0.661, 0.606, 0.587, 0.623, a spread within a single fold of up to 0.531; between-fold variability therefore exceeds between-method variability.
The permutation null was obtained from 100 label permutations at the patient level per outer fold; with $m = 100$ permutations the smallest attainable $p$-value is $p = (0+1)/(100+1) = 0.0099$, which bounds the resolution of every permutation result in this table.
The Null AUC, $\kappa$, $p$ and $q$ columns were computed from the \emph{mean} AUC across folds, whereas the AUC column reports the \emph{median}; the two summaries are not interchangeable and are shown side by side deliberately.
10 of the 24 combinations reach $p < 0.05$ uncorrected and 6 survive Benjamini--Hochberg correction\NOTETODO{cite Benjamini and Hochberg 1995; the key is missing from BibFile.bib} for AUC, while none survives it for accuracy.
No centerline combination survives correction; across the eight of them the mean AUC ranges from 0.396 to 0.495 against a null of 0.489--0.510, $\kappa$ from $-0.206$ to $-0.014$ and $p$ from 0.5347 to 1.0000.\\
\textsuperscript{a} In these 2 combinations the $\ell_1$ penalty shrank every coefficient to zero in every fold and the code substituted the index of the first column; the retained feature is a fallback, not one the selector identified as informative.
Combinations with a median of one retained feature but no mark reached that median with genuinely non-zero coefficients.\\
\textsuperscript{b} The unit of analysis of the voxel-based datasets is the voxel, not the patient, and their metrics are therefore \emph{not} comparable with the remaining rows.
Across the five folds the voxel/whole datasets contribute 396085 voxel-level predictions (91500, 79327, 75748, 82369, 67141) and the voxel/centerline datasets 8557 (1690, 1571, 1811, 1806, 1679), against 141 patient-level predictions (29, 28, 28, 28, 28) for every other dataset.
The consequence is visible in the permutation null, whose standard deviation is 0.040--0.052 for the voxel-based combinations against 0.063--0.120 for all others, so that the voxel-based $p$-values are computed against a spuriously narrow null and are not comparable with the rest of the table.
3 voxel-based combinations are formally significant for AUC after correction; they must not be read as evidence for a signal.
\NOTETODO{patient-level majority-vote aggregation of the voxel-based predictions, which would make these rows comparable}
\par}

\clearpage

\begin{table}[htbp]
\centering
\caption{External validation in the MR CLEAN Registry cohort. Each pipeline was refitted on all 141 development patients with its tuned hyperparameters and applied once to the Registry cohort. CV, cross-validated median over the five outer folds of the same combination in Table~\ref{tab:cv-performance}; Reg., Registry; ACC, accuracy; $\Delta$AUC, Registry AUC minus cross-validated median AUC, a negative value being a loss of discrimination on transfer; other abbreviations as in Table~\ref{tab:cv-performance}. The voxel-based datasets were not transferred externally and therefore have no row here, which is why the table has 16 rows rather than 24.}
\label{tab:external-validation}
\smallskip
\input{\tablesdir/table4_external_validation}
\end{table}

\smallskip
{\footnotesize
\sloppy
\noindent
\textit{Notes to Table~\ref{tab:external-validation}.}
Across the 16 models the Registry AUC ranged from 0.471 to 0.521 with a median of 0.496, and 5 of them exceeded 0.50.
Registry accuracy ranged from 0.440 to 0.550 with a median of 0.481, and no model reached the majority-class accuracy of the development cohort, 0.582.
The sign of $\Delta$ follows the cross-validated value and nothing else: the 8 combinations whose cross-validated median AUC was above 0.500 all moved downwards on transfer, by $-0.034$ to $-0.190$, and the 8 that were at or below it all moved upwards, by $0.002$ to $0.128$.
The upward movement is therefore a regression to chance of models that had performed below chance internally, not a gain in generalization, and it must not be presented as one.
The only combinations for which the comparison is informative are the 6 that discriminated in cross-validation, that is reached a cross-validated median AUC of at least 0.600; all of them lose between 0.150 and 0.190 of AUC.
\NOTEWARN{$n$ Registry is not stated in the results files and is not the same for every model. Reconstructing the smallest denominator compatible with the decimal expansion of each reported accuracy leaves exactly three cohort sizes: 1649 for eight models, 1760 for six and 1761 for two. Neither the 1906 of the exclusion cascade nor the 1839 identifiers stored with the cohort is compatible with any of the sixteen. Until this is resolved the AUC values in this table are not strictly comparable with one another}
\par}

\clearpage

\begin{table}[htbp]
\centering
\caption{Comparison with the reference pipeline of van Voorst et al. \cite{vanVoorst2022} The first row summarizes the present study; rows two to four were computed on the MR CLEAN NO-IV development cohort with the same five outer folds as in Table~\ref{tab:cv-performance}, so that the comparison is made within one dataset rather than across publications; the last row is the value reported in the original publication, on a different cohort and for a different endpoint. ICC, intraclass correlation coefficient: of the 107 features the reference pipeline extracts, the 52 with ICC $\geq$ 0.75 are retained and the rest discarded. AUC$_{\text{prob}}$ denotes the AUC computed from the predicted class probabilities, and is the quantity reported for the reference pipeline. Values are medians with the interquartile range in brackets, or means $\pm$ standard deviation across folds, as indicated in the note below. An em dash marks a quantity the row does not have: the pooled row spans combinations with different accuracies, and the published row reports neither a range nor an accuracy.}
\label{tab:reference-comparison}
\smallskip
\input{\tablesdir/table5_reference_comparison}
\end{table}

\smallskip
{\footnotesize
\sloppy
\noindent
\textit{Notes to Table~\ref{tab:reference-comparison}.}
The first row pools 25 fold results, 5 combinations $\times$ 5 outer folds, and its bracketed values are the interquartile range; the second row is a mean $\pm$ standard deviation over the 5 splits; rows three and four are medians over the 5 folds, with no further feature selection applied to the 52 ICC-filtered features.
The last row is quoted from the publication: thrombus radiomics for eTICI $\geq$ 2B, a more permissive endpoint than the eTICI 2C--3 outcome of this article, as a mean over 100 bootstrap samples of 100 patients.
The plain \texttt{AUC} column of the reference pipeline degenerates to exactly 0.50 in 4 of the five splits, because the fitted model assigns the majority class to every patient in those splits and its hard predictions therefore carry no ranking information.
AUC$_{\text{prob}}$, computed from the predicted probabilities, is the only quantity of that pipeline that is comparable with the AUC values reported elsewhere in this article, and it is the one used here.
By median accuracy the reference pipeline and both of our classifiers restricted to its 52 ICC-filtered features stay below the majority-class accuracy of 0.582, although the reference pipeline exceeds it in one individual split.
Rows two to four were recomputed against the eTICI 2C--3 outcome of this article --- the stored label column is named \texttt{TICI\_2B-3}, but its per-fold positive counts are exactly the 82/141 of this study --- and are directly comparable with the first row, whereas the published 0.52 of the last row was obtained for eTICI $\geq$ 2B in a different cohort and differs from them in endpoint and cohort as much as in pipeline.
\NOTEQUERY{the pool of the first row omits one whole-thrombus combination that also discriminated in cross-validation, and the reason is not documented in the source material}
\NOTEQUERY{decide which of the two mutually incompatible sets of numbers in the source data applies before this table is finalized: set A, 0.481 $\pm$ 0.145 on 141 patients against set B, 0.467 $\pm$ 0.099 on 170}
\par}

\clearpage

\begin{table}[htbp]
\centering
\caption{Ensemble analysis: the 15 configurations of discrete, that is binary, AdaBoost \cite{Freund1997} applied to the stored prediction vectors of the 16 patient-level weak classifiers (four patient-level datasets $\times$ four pipelines). Each configuration was fitted on the five outer folds, giving 75 ensemble models. ACC $\varnothing$, mean accuracy of the configuration across the five folds; best weak classifier, the single pipeline in the configuration's pool with the highest mean accuracy, and its value; $\Delta$, the mean accuracy of the ensemble minus that of the best weak classifier. A negative $\Delta$ means the ensemble is worse than the best model it combines. The configuration name states which of the four patient-level datasets form its pool, each contributing all four pipelines: \texttt{type1} pools all four datasets (16 weak classifiers), \texttt{type2} a whole-thrombus dataset with its own centerline counterpart (8), \texttt{type3} a single dataset (4), \texttt{type4} two datasets that are not such a pair (8) and \texttt{type5} all datasets but one (12). \texttt{distrib.}\ abbreviates the distribution-based dataset and \emph{cent.} the centerline region of interest, as in Table~\ref{tab:cv-performance}. The two voxel-based datasets were excluded from every pool because their predictions are defined per voxel.}
\label{tab:ensemble}
\smallskip
\input{\tablesdir/table6_ensemble}
\end{table}

\smallskip
{\footnotesize
\sloppy
\noindent
\textit{Notes to Table~\ref{tab:ensemble}.}
$\Delta$ is negative for all 15 configurations, between $-0.022$ and $-0.121$; no ensemble outperformed the best single pipeline in its own pool.
The worst result belongs to \texttt{type1\_\allowbreak all\_\allowbreak featuresets}, the configuration that pools all feature sets.
Across all 75 fitted ensemble models accuracy averages $0.543 \pm 0.083$ with a range of 0.414--0.714.
The mechanism is documented by the fitted weights: in the all-feature-set configuration, 55 of the 80 boosting rounds (16 rounds $\times$ five folds) received a negative weight $\alpha$, meaning that the weak classifier admitted in that round performed worse than chance on the reweighted data.
The individual pipelines therefore do not make complementary errors, and the boosting assumption of weak classifiers reliably better than chance is not met.
Labels were coded as $\{-1,+1\}$, the weight of a round was $\alpha = \tfrac{1}{2}\ln\frac{1-\varepsilon}{\varepsilon}$ with $\varepsilon$ the error of the admitted classifier on the currently reweighted data, and each weak classifier could be admitted at most once, so a pool of $k$ members yields at most $k$ rounds.
The weak classifiers were not retrained: the procedure is applied to models that had already been fitted, which makes it methodologically a form of stacking rather than boosting proper.
The best weak classifier of \texttt{type3\_\allowbreak global} and of \texttt{type1\_\allowbreak all\_\allowbreak featuresets} is the same pipeline, because the pool of the latter contains the pool of the former.
\par}

\clearpage

\begin{table}[htbp]
\centering
\caption{Pairwise comparison of the pipelines by the corrected resampled $t$-test of Nadeau and Bengio \cite{Nadeau2003}. The upper block compares the four pipelines of this study with one another on the whole-thrombus global-based dataset; the lower block compares them, and our two classifiers restricted to the reference pipeline's own 52 ICC-filtered features, against the reference pipeline of van Voorst et al.\ recomputed on the same cohort. Every model in the table was evaluated on the same 5 outer folds of the same 141 patients, which is what makes the paired test applicable. Because the training sets of $k$-fold cross-validation overlap, the per-fold differences are not independent and the ordinary paired $t$-test underestimates their variance; the correction multiplies the variance by $1/k + n_{\text{test}}/n_{\text{train}}$, here $1/5 + 28.2/112.8 = 0.4500$, which inflates the standard error 1.50-fold at $\mathrm{df} = 4$. Fold sizes are taken from the data, not assumed. Mean $\Delta$AUC is the mean of the 5 per-fold differences in AUC between pipeline A and pipeline B; a positive value favours A. $t$ (uncorrected) is the ordinary paired statistic, $t$ (corrected) the same quantity with the inflated variance. The two-sided critical value at $\alpha = 0.05$ and $\mathrm{df} = 4$ is 2.776. Pipeline names combine selector and classifier as in Table~\ref{tab:cv-performance}.}
\label{tab:nadeau-bengio}
\smallskip
\input{\tablesdir/table7_nadeau_bengio}
\end{table}

\smallskip
{\footnotesize
\sloppy
\noindent
\textit{Notes to Table~\ref{tab:nadeau-bengio}.}
Not one of the 12 comparisons reaches significance: the largest corrected statistic in absolute value is 1.54, against a critical value of 2.776.
The four pipelines of this study are therefore statistically indistinguishable from one another, and none of them is distinguishable from the reference pipeline either, in either direction.
For the reference pipeline the compared quantity is AUC$_{\text{prob}}$, because its plain \texttt{AUC} column degenerates to exactly 0.50 in four of the five splits; see Table~\ref{tab:reference-comparison}.
The pooled row of Table~\ref{tab:reference-comparison} is deliberately absent here: those 5 combinations were picked by their results on these same folds, so testing them against the reference pipeline would be biased in favour of the present study. All four pipelines of the dataset are tested instead, whatever their outcome.
With 5 folds the test has $\mathrm{df} = 4$ and very little power, so these results are an absence of evidence for a difference, not evidence that the pipelines perform identically. The comparison that the data would support best, a DeLong test on the external ROC curves, was not possible because per-patient predicted scores on the Registry were not retained.
\par}

\clearpage


%%     \restoregeometry
%%
%% Cislovani S1--S7 misto 1--7 zaridi pred vlozenim:
%%     \setcounter{table}{0}\renewcommand{\thetable}{S\arabic{table}}
%%
%% Preambule potrebuje: booktabs, array, amsmath, amssymb.
%% Cesta k tabulkam se bere z \tablesdir; vychozi hodnota plati pro preklad
%% z korene projektu.
%% -------------------------------------------------------------------------
\providecommand{\tablesdir}{tables_tom/output}

%% Zaloha, kdyby `manuscriptnotes` nebyl nactveny: vysadi se jen kotva,
%% poznamka se zahodi -- tedy tentyz vysledek jako v rezimu odesilani.
\providecommand{\NOTE}[2][]{#1}
\providecommand{\NOTEWARN}[2][]{#1}
\providecommand{\NOTECHECK}[2][]{#1}
\providecommand{\NOTETODO}[2][]{#1}
\providecommand{\NOTEQUERY}[2][]{#1}

{\large\bfseries Supplementary tables\par}
\bigskip
\NOTEWARN{the captions and notes in this supplement were drafted with AI assistance and are provisional; every one of them has to be rewritten and checked by an author before submission}

\begin{table}[htbp]
\centering
\caption{Patient selection in the two cohorts. MR CLEAN NO-IV was used for model development and cross-validation, the MR CLEAN Registry as the independent external cohort. Entries are numbers of patients; an em dash marks an exclusion step that does not apply to that cohort. eTICI, expanded Thrombolysis in Cerebral Infarction score; the outcome modeled throughout this article is near-complete or complete reperfusion, defined as a final eTICI score of 2C or 3. Slice spacing refers to the through-plane spacing of the non-contrast CT acquisition.}
\label{tab:cohorts}
\smallskip
\input{\tablesdir/table1_cohorts}
\end{table}

\smallskip
{\footnotesize
\sloppy
\noindent
\textit{Notes to Table~\ref{tab:cohorts}.}
The manual review of segmentation quality was carried out on the MR CLEAN NO-IV cohort only.
No segmentation was reviewed or rejected in the MR CLEAN Registry cohort, so the two cohorts differ in segmentation quality control as well as in population.
This is a property of the design and is the single largest asymmetry between the development and the external cohort.
The review was blinded to the eTICI score: it judged the quality and the anatomical placement of the segmentation only, so the 155 exclusions it produced cannot have been informed by the outcome.
Of the 141 patients that remain, 82 (58.2\%) reached eTICI 2C--3; the proportion varies across the five outer folds (0.448, 0.714, 0.536, 0.679, 0.536), so in the first fold the majority class is the absence rather than the presence of near-complete or complete reperfusion.
\NOTETODO{breakdown of how many of the 155 reviewed cases were rejected by which criterion}
\NOTETODO{number and proportion of patients with eTICI 2C--3 reperfusion in the Registry cohort, which the table leaves as an em dash}
\NOTEWARN{size of the Registry cohort: the exclusion cascade gives 1906, the stored identifier list holds 1839, and the denominators reconstructed from the decimal expansion of the reported external accuracies are 1649 (8 models), 1760 (6 models), 1761 (2 models). Neither 1906 nor 1839 is compatible with any of the sixteen transferred models, and the reconstructed values are mutually incompatible, so the external cohort is not the same across feature sets}
\par}

\clearpage

\begin{table}[htbp]
\centering
\caption{Hyperparameter search spaces of the four pipelines. LASSO, $\ell_1$-penalized linear regression on the binary outcome used as a feature selector; RFI, ranking by the gain-based feature importance of a random-forest gradient-boosting model; XGBRF, \texttt{XGBRFClassifier}; LinearSVC, linear support vector classifier with an $\ell_1$ penalty, which acts as selector and classifier simultaneously; NuSVC, $\nu$-support vector classifier with a radial basis function kernel. In the Component column, (S) marks the component acting as feature selector, (C) as classifier, and (S+C) a single object that performs both roles. Scale indicates how Optuna sampled the parameter: \emph{log}, log-uniform over the interval; \emph{int}, uniform over the integers of the interval; \emph{float}, uniform over the interval.}
\label{tab:hyperparameters}
\smallskip
\input{\tablesdir/table2_hyperparameters}
\end{table}

\smallskip
{\footnotesize
\sloppy
\noindent
\textit{Notes to Table~\ref{tab:hyperparameters}.}
$n$ and \texttt{X.shape[1]} denote the number of rows and the number of columns of the training matrix passed to the search.
Rows are patients for the global-based and distribution-based datasets and voxels for the voxel-based datasets, so the same expression yields different intervals across datasets; the number of columns reaches 20660 for the distribution-based datasets.
The minority fraction is the proportion of the smaller outcome class among the training labels of the given fold.
In the RFI--RF pipeline a single parameter dictionary configures both the model whose gain importances rank the features and the classifier fitted on the retained ones, which is why its rows are marked (S+C) rather than (C).
All search spaces were explored with Optuna using the Gaussian-process sampler (\texttt{GPSampler}) with 20 trials per dataset-by-pipeline combination.
The objective maximized was the \emph{mean accuracy over the five inner cross-validation folds}, not AUC; the inner folds were generated by \texttt{StratifiedGroupKFold} grouped by \texttt{patient\_id}.
Settings that were fixed and not searched, one per line so that the long identifiers have somewhere to break:\\
\texttt{Lasso(fit\_intercept=True, max\_iter=5000)}\\
\texttt{LinearSVC(penalty='l1', max\_iter=5000, class\_weight='balanced')}\\
\texttt{NuSVC(kernel='rbf', probability=True, class\_weight='balanced')}\\
\texttt{XGBRFClassifier(random\_state=42)}\\
Class weighting was therefore applied in the two support vector branches only.
The models were seeded with \texttt{random\_state=42}, but \texttt{GPSampler()} was instantiated without a seed, so the sequence of 20 trials is not reproducible.
\par}

\clearpage

\begin{table}[htbp]
\centering
\caption{Cross-validated performance of all 24 dataset-by-pipeline combinations in the MR CLEAN NO-IV cohort ($n = 141$), over five outer folds. Feature set: \emph{global}, eight summary statistics of each radiomic feature over the thrombus, one row per patient; \emph{distribution}, a ten-bin histogram of each radiomic feature over the thrombus, one row per patient; \emph{voxel}, the radiomic feature values themselves, one row per voxel. ROI: \emph{whole}, the dilated thrombus mask; \emph{centerline}, the dilated skeleton of the thrombus mask. Selector (Sel.): LASSO, $\ell_1$-penalized regression; L1, embedded selection by the non-zero coefficients of the classifier itself; RFI, gain-based feature importance of an \texttt{XGBRFClassifier}. Classifier (Class.): XGBRF, \texttt{XGBRFClassifier}; LinSVC, linear support vector classifier with an $\ell_1$ penalty; NuSVC, $\nu$-support vector classifier with a radial basis function kernel. In the body of the table the distribution-based feature set is abbreviated \emph{distrib.}\ and the centerline ROI \emph{cent.} AUC, area under the receiver operating characteristic curve; ACC, accuracy; IQR, interquartile range across the five folds; min--max, the smallest and largest value across the five folds; Med.\ $n$, median number of features retained by the selector; Null AUC, mean AUC of the permutation baseline; $\kappa$, the $\kappa$-type index of the observed AUC relative to that baseline; $p$, permutation $p$-value; $q$, the same $p$-value after Benjamini--Hochberg correction across all 24 combinations.}
\label{tab:cv-performance}
\smallskip
\input{\tablesdir/table3_cv_performance}
\end{table}

\smallskip
{\footnotesize
\sloppy
\noindent
\textit{Notes to Table~\ref{tab:cv-performance}.}
Accuracy must be read against the majority-class accuracy of the development cohort, $82/141 = 0.582$; only 4 of the 24 combinations exceed it by median accuracy, and 86 of the 120 individual fold results (71.7\%) fall below it.
The class ratio itself varies across folds (0.448, 0.714, 0.536, 0.679, 0.536), so in the first fold the majority class is the absence rather than the presence of near-complete or complete reperfusion.
The largest single fold result of the whole study is 0.810, obtained by distribution/whole/RFI/NuSVC in fold 4.
Averaged over the eight whole-thrombus patient-level combinations, AUC per fold is 0.510, 0.661, 0.606, 0.587, 0.623, a spread within a single fold of up to 0.531; between-fold variability therefore exceeds between-method variability.
The permutation null was obtained from 100 label permutations at the patient level per outer fold; with $m = 100$ permutations the smallest attainable $p$-value is $p = (0+1)/(100+1) = 0.0099$, which bounds the resolution of every permutation result in this table.
The Null AUC, $\kappa$, $p$ and $q$ columns were computed from the \emph{mean} AUC across folds, whereas the AUC column reports the \emph{median}; the two summaries are not interchangeable and are shown side by side deliberately.
10 of the 24 combinations reach $p < 0.05$ uncorrected and 6 survive Benjamini--Hochberg correction\NOTETODO{cite Benjamini and Hochberg 1995; the key is missing from BibFile.bib} for AUC, while none survives it for accuracy.
No centerline combination survives correction; across the eight of them the mean AUC ranges from 0.396 to 0.495 against a null of 0.489--0.510, $\kappa$ from $-0.206$ to $-0.014$ and $p$ from 0.5347 to 1.0000.\\
\textsuperscript{a} In these 2 combinations the $\ell_1$ penalty shrank every coefficient to zero in every fold and the code substituted the index of the first column; the retained feature is a fallback, not one the selector identified as informative.
Combinations with a median of one retained feature but no mark reached that median with genuinely non-zero coefficients.\\
\textsuperscript{b} The unit of analysis of the voxel-based datasets is the voxel, not the patient, and their metrics are therefore \emph{not} comparable with the remaining rows.
Across the five folds the voxel/whole datasets contribute 396085 voxel-level predictions (91500, 79327, 75748, 82369, 67141) and the voxel/centerline datasets 8557 (1690, 1571, 1811, 1806, 1679), against 141 patient-level predictions (29, 28, 28, 28, 28) for every other dataset.
The consequence is visible in the permutation null, whose standard deviation is 0.040--0.052 for the voxel-based combinations against 0.063--0.120 for all others, so that the voxel-based $p$-values are computed against a spuriously narrow null and are not comparable with the rest of the table.
3 voxel-based combinations are formally significant for AUC after correction; they must not be read as evidence for a signal.
\NOTETODO{patient-level majority-vote aggregation of the voxel-based predictions, which would make these rows comparable}
\par}

\clearpage

\begin{table}[htbp]
\centering
\caption{External validation in the MR CLEAN Registry cohort. Each pipeline was refitted on all 141 development patients with its tuned hyperparameters and applied once to the Registry cohort. CV, cross-validated median over the five outer folds of the same combination in Table~\ref{tab:cv-performance}; Reg., Registry; ACC, accuracy; $\Delta$AUC, Registry AUC minus cross-validated median AUC, a negative value being a loss of discrimination on transfer; other abbreviations as in Table~\ref{tab:cv-performance}. The voxel-based datasets were not transferred externally and therefore have no row here, which is why the table has 16 rows rather than 24.}
\label{tab:external-validation}
\smallskip
\input{\tablesdir/table4_external_validation}
\end{table}

\smallskip
{\footnotesize
\sloppy
\noindent
\textit{Notes to Table~\ref{tab:external-validation}.}
Across the 16 models the Registry AUC ranged from 0.471 to 0.521 with a median of 0.496, and 5 of them exceeded 0.50.
Registry accuracy ranged from 0.440 to 0.550 with a median of 0.481, and no model reached the majority-class accuracy of the development cohort, 0.582.
The sign of $\Delta$ follows the cross-validated value and nothing else: the 8 combinations whose cross-validated median AUC was above 0.500 all moved downwards on transfer, by $-0.034$ to $-0.190$, and the 8 that were at or below it all moved upwards, by $0.002$ to $0.128$.
The upward movement is therefore a regression to chance of models that had performed below chance internally, not a gain in generalization, and it must not be presented as one.
The only combinations for which the comparison is informative are the 6 that discriminated in cross-validation, that is reached a cross-validated median AUC of at least 0.600; all of them lose between 0.150 and 0.190 of AUC.
\NOTEWARN{$n$ Registry is not stated in the results files and is not the same for every model. Reconstructing the smallest denominator compatible with the decimal expansion of each reported accuracy leaves exactly three cohort sizes: 1649 for eight models, 1760 for six and 1761 for two. Neither the 1906 of the exclusion cascade nor the 1839 identifiers stored with the cohort is compatible with any of the sixteen. Until this is resolved the AUC values in this table are not strictly comparable with one another}
\par}

\clearpage

\begin{table}[htbp]
\centering
\caption{Comparison with the reference pipeline of van Voorst et al. \cite{vanVoorst2022} The first row summarizes the present study; rows two to four were computed on the MR CLEAN NO-IV development cohort with the same five outer folds as in Table~\ref{tab:cv-performance}, so that the comparison is made within one dataset rather than across publications; the last row is the value reported in the original publication, on a different cohort and for a different endpoint. ICC, intraclass correlation coefficient: of the 107 features the reference pipeline extracts, the 52 with ICC $\geq$ 0.75 are retained and the rest discarded. AUC$_{\text{prob}}$ denotes the AUC computed from the predicted class probabilities, and is the quantity reported for the reference pipeline. Values are medians with the interquartile range in brackets, or means $\pm$ standard deviation across folds, as indicated in the note below. An em dash marks a quantity the row does not have: the pooled row spans combinations with different accuracies, and the published row reports neither a range nor an accuracy.}
\label{tab:reference-comparison}
\smallskip
\input{\tablesdir/table5_reference_comparison}
\end{table}

\smallskip
{\footnotesize
\sloppy
\noindent
\textit{Notes to Table~\ref{tab:reference-comparison}.}
The first row pools 25 fold results, 5 combinations $\times$ 5 outer folds, and its bracketed values are the interquartile range; the second row is a mean $\pm$ standard deviation over the 5 splits; rows three and four are medians over the 5 folds, with no further feature selection applied to the 52 ICC-filtered features.
The last row is quoted from the publication: thrombus radiomics for eTICI $\geq$ 2B, a more permissive endpoint than the eTICI 2C--3 outcome of this article, as a mean over 100 bootstrap samples of 100 patients.
The plain \texttt{AUC} column of the reference pipeline degenerates to exactly 0.50 in 4 of the five splits, because the fitted model assigns the majority class to every patient in those splits and its hard predictions therefore carry no ranking information.
AUC$_{\text{prob}}$, computed from the predicted probabilities, is the only quantity of that pipeline that is comparable with the AUC values reported elsewhere in this article, and it is the one used here.
By median accuracy the reference pipeline and both of our classifiers restricted to its 52 ICC-filtered features stay below the majority-class accuracy of 0.582, although the reference pipeline exceeds it in one individual split.
Rows two to four were recomputed against the eTICI 2C--3 outcome of this article --- the stored label column is named \texttt{TICI\_2B-3}, but its per-fold positive counts are exactly the 82/141 of this study --- and are directly comparable with the first row, whereas the published 0.52 of the last row was obtained for eTICI $\geq$ 2B in a different cohort and differs from them in endpoint and cohort as much as in pipeline.
\NOTEQUERY{the pool of the first row omits one whole-thrombus combination that also discriminated in cross-validation, and the reason is not documented in the source material}
\NOTEQUERY{decide which of the two mutually incompatible sets of numbers in the source data applies before this table is finalized: set A, 0.481 $\pm$ 0.145 on 141 patients against set B, 0.467 $\pm$ 0.099 on 170}
\par}

\clearpage

\begin{table}[htbp]
\centering
\caption{Ensemble analysis: the 15 configurations of discrete, that is binary, AdaBoost \cite{Freund1997} applied to the stored prediction vectors of the 16 patient-level weak classifiers (four patient-level datasets $\times$ four pipelines). Each configuration was fitted on the five outer folds, giving 75 ensemble models. ACC $\varnothing$, mean accuracy of the configuration across the five folds; best weak classifier, the single pipeline in the configuration's pool with the highest mean accuracy, and its value; $\Delta$, the mean accuracy of the ensemble minus that of the best weak classifier. A negative $\Delta$ means the ensemble is worse than the best model it combines. The configuration name states which of the four patient-level datasets form its pool, each contributing all four pipelines: \texttt{type1} pools all four datasets (16 weak classifiers), \texttt{type2} a whole-thrombus dataset with its own centerline counterpart (8), \texttt{type3} a single dataset (4), \texttt{type4} two datasets that are not such a pair (8) and \texttt{type5} all datasets but one (12). \texttt{distrib.}\ abbreviates the distribution-based dataset and \emph{cent.} the centerline region of interest, as in Table~\ref{tab:cv-performance}. The two voxel-based datasets were excluded from every pool because their predictions are defined per voxel.}
\label{tab:ensemble}
\smallskip
\input{\tablesdir/table6_ensemble}
\end{table}

\smallskip
{\footnotesize
\sloppy
\noindent
\textit{Notes to Table~\ref{tab:ensemble}.}
$\Delta$ is negative for all 15 configurations, between $-0.022$ and $-0.121$; no ensemble outperformed the best single pipeline in its own pool.
The worst result belongs to \texttt{type1\_\allowbreak all\_\allowbreak featuresets}, the configuration that pools all feature sets.
Across all 75 fitted ensemble models accuracy averages $0.543 \pm 0.083$ with a range of 0.414--0.714.
The mechanism is documented by the fitted weights: in the all-feature-set configuration, 55 of the 80 boosting rounds (16 rounds $\times$ five folds) received a negative weight $\alpha$, meaning that the weak classifier admitted in that round performed worse than chance on the reweighted data.
The individual pipelines therefore do not make complementary errors, and the boosting assumption of weak classifiers reliably better than chance is not met.
Labels were coded as $\{-1,+1\}$, the weight of a round was $\alpha = \tfrac{1}{2}\ln\frac{1-\varepsilon}{\varepsilon}$ with $\varepsilon$ the error of the admitted classifier on the currently reweighted data, and each weak classifier could be admitted at most once, so a pool of $k$ members yields at most $k$ rounds.
The weak classifiers were not retrained: the procedure is applied to models that had already been fitted, which makes it methodologically a form of stacking rather than boosting proper.
The best weak classifier of \texttt{type3\_\allowbreak global} and of \texttt{type1\_\allowbreak all\_\allowbreak featuresets} is the same pipeline, because the pool of the latter contains the pool of the former.
\par}

\clearpage

\begin{table}[htbp]
\centering
\caption{Pairwise comparison of the pipelines by the corrected resampled $t$-test of Nadeau and Bengio \cite{Nadeau2003}. The upper block compares the four pipelines of this study with one another on the whole-thrombus global-based dataset; the lower block compares them, and our two classifiers restricted to the reference pipeline's own 52 ICC-filtered features, against the reference pipeline of van Voorst et al.\ recomputed on the same cohort. Every model in the table was evaluated on the same 5 outer folds of the same 141 patients, which is what makes the paired test applicable. Because the training sets of $k$-fold cross-validation overlap, the per-fold differences are not independent and the ordinary paired $t$-test underestimates their variance; the correction multiplies the variance by $1/k + n_{\text{test}}/n_{\text{train}}$, here $1/5 + 28.2/112.8 = 0.4500$, which inflates the standard error 1.50-fold at $\mathrm{df} = 4$. Fold sizes are taken from the data, not assumed. Mean $\Delta$AUC is the mean of the 5 per-fold differences in AUC between pipeline A and pipeline B; a positive value favours A. $t$ (uncorrected) is the ordinary paired statistic, $t$ (corrected) the same quantity with the inflated variance. The two-sided critical value at $\alpha = 0.05$ and $\mathrm{df} = 4$ is 2.776. Pipeline names combine selector and classifier as in Table~\ref{tab:cv-performance}.}
\label{tab:nadeau-bengio}
\smallskip
\input{\tablesdir/table7_nadeau_bengio}
\end{table}

\smallskip
{\footnotesize
\sloppy
\noindent
\textit{Notes to Table~\ref{tab:nadeau-bengio}.}
Not one of the 12 comparisons reaches significance: the largest corrected statistic in absolute value is 1.54, against a critical value of 2.776.
The four pipelines of this study are therefore statistically indistinguishable from one another, and none of them is distinguishable from the reference pipeline either, in either direction.
For the reference pipeline the compared quantity is AUC$_{\text{prob}}$, because its plain \texttt{AUC} column degenerates to exactly 0.50 in four of the five splits; see Table~\ref{tab:reference-comparison}.
The pooled row of Table~\ref{tab:reference-comparison} is deliberately absent here: those 5 combinations were picked by their results on these same folds, so testing them against the reference pipeline would be biased in favour of the present study. All four pipelines of the dataset are tested instead, whatever their outcome.
With 5 folds the test has $\mathrm{df} = 4$ and very little power, so these results are an absence of evidence for a difference, not evidence that the pipelines perform identically. The comparison that the data would support best, a DeLong test on the external ROC curves, was not possible because per-patient predicted scores on the Registry were not retained.
\par}

\clearpage


%%
%% Tabulky se vejdou i do uzke sazby `elsarticle` (390 pt), ale zdrobni.
%% Aby vypadaly jako v samostatnem supplementu, obal se vlozeni geometrii:
%%
%%     \usepackage{geometry}          % v preambuli, za elsarticle
%%     ...
%%     \clearpage
%%     \newgeometry{margin=15mm}
%%     %% -------------------------------------------------------------------------
%% GENEROVANO: make_all.py -- needitovat rucne.
%%
%% Popisky, tabulky a poznamky bez preambule. Do clanku staci jeden radek:
%%
%%     %% -------------------------------------------------------------------------
%% GENEROVANO: make_all.py -- needitovat rucne.
%%
%% Popisky, tabulky a poznamky bez preambule. Do clanku staci jeden radek:
%%
%%     \input{tables_tom/output/supplement_body}
%%
%% Tabulky se vejdou i do uzke sazby `elsarticle` (390 pt), ale zdrobni.
%% Aby vypadaly jako v samostatnem supplementu, obal se vlozeni geometrii:
%%
%%     \usepackage{geometry}          % v preambuli, za elsarticle
%%     ...
%%     \clearpage
%%     \newgeometry{margin=15mm}
%%     \input{tables_tom/output/supplement_body}
%%     \restoregeometry
%%
%% Cislovani S1--S7 misto 1--7 zaridi pred vlozenim:
%%     \setcounter{table}{0}\renewcommand{\thetable}{S\arabic{table}}
%%
%% Preambule potrebuje: booktabs, array, amsmath, amssymb.
%% Cesta k tabulkam se bere z \tablesdir; vychozi hodnota plati pro preklad
%% z korene projektu.
%% -------------------------------------------------------------------------
\providecommand{\tablesdir}{tables_tom/output}

%% Zaloha, kdyby `manuscriptnotes` nebyl nactveny: vysadi se jen kotva,
%% poznamka se zahodi -- tedy tentyz vysledek jako v rezimu odesilani.
\providecommand{\NOTE}[2][]{#1}
\providecommand{\NOTEWARN}[2][]{#1}
\providecommand{\NOTECHECK}[2][]{#1}
\providecommand{\NOTETODO}[2][]{#1}
\providecommand{\NOTEQUERY}[2][]{#1}

{\large\bfseries Supplementary tables\par}
\bigskip
\NOTEWARN{the captions and notes in this supplement were drafted with AI assistance and are provisional; every one of them has to be rewritten and checked by an author before submission}

\begin{table}[htbp]
\centering
\caption{Patient selection in the two cohorts. MR CLEAN NO-IV was used for model development and cross-validation, the MR CLEAN Registry as the independent external cohort. Entries are numbers of patients; an em dash marks an exclusion step that does not apply to that cohort. eTICI, expanded Thrombolysis in Cerebral Infarction score; the outcome modeled throughout this article is near-complete or complete reperfusion, defined as a final eTICI score of 2C or 3. Slice spacing refers to the through-plane spacing of the non-contrast CT acquisition.}
\label{tab:cohorts}
\smallskip
\input{\tablesdir/table1_cohorts}
\end{table}

\smallskip
{\footnotesize
\sloppy
\noindent
\textit{Notes to Table~\ref{tab:cohorts}.}
The manual review of segmentation quality was carried out on the MR CLEAN NO-IV cohort only.
No segmentation was reviewed or rejected in the MR CLEAN Registry cohort, so the two cohorts differ in segmentation quality control as well as in population.
This is a property of the design and is the single largest asymmetry between the development and the external cohort.
The review was blinded to the eTICI score: it judged the quality and the anatomical placement of the segmentation only, so the 155 exclusions it produced cannot have been informed by the outcome.
Of the 141 patients that remain, 82 (58.2\%) reached eTICI 2C--3; the proportion varies across the five outer folds (0.448, 0.714, 0.536, 0.679, 0.536), so in the first fold the majority class is the absence rather than the presence of near-complete or complete reperfusion.
\NOTETODO{breakdown of how many of the 155 reviewed cases were rejected by which criterion}
\NOTETODO{number and proportion of patients with eTICI 2C--3 reperfusion in the Registry cohort, which the table leaves as an em dash}
\NOTEWARN{size of the Registry cohort: the exclusion cascade gives 1906, the stored identifier list holds 1839, and the denominators reconstructed from the decimal expansion of the reported external accuracies are 1649 (8 models), 1760 (6 models), 1761 (2 models). Neither 1906 nor 1839 is compatible with any of the sixteen transferred models, and the reconstructed values are mutually incompatible, so the external cohort is not the same across feature sets}
\par}

\clearpage

\begin{table}[htbp]
\centering
\caption{Hyperparameter search spaces of the four pipelines. LASSO, $\ell_1$-penalized linear regression on the binary outcome used as a feature selector; RFI, ranking by the gain-based feature importance of a random-forest gradient-boosting model; XGBRF, \texttt{XGBRFClassifier}; LinearSVC, linear support vector classifier with an $\ell_1$ penalty, which acts as selector and classifier simultaneously; NuSVC, $\nu$-support vector classifier with a radial basis function kernel. In the Component column, (S) marks the component acting as feature selector, (C) as classifier, and (S+C) a single object that performs both roles. Scale indicates how Optuna sampled the parameter: \emph{log}, log-uniform over the interval; \emph{int}, uniform over the integers of the interval; \emph{float}, uniform over the interval.}
\label{tab:hyperparameters}
\smallskip
\input{\tablesdir/table2_hyperparameters}
\end{table}

\smallskip
{\footnotesize
\sloppy
\noindent
\textit{Notes to Table~\ref{tab:hyperparameters}.}
$n$ and \texttt{X.shape[1]} denote the number of rows and the number of columns of the training matrix passed to the search.
Rows are patients for the global-based and distribution-based datasets and voxels for the voxel-based datasets, so the same expression yields different intervals across datasets; the number of columns reaches 20660 for the distribution-based datasets.
The minority fraction is the proportion of the smaller outcome class among the training labels of the given fold.
In the RFI--RF pipeline a single parameter dictionary configures both the model whose gain importances rank the features and the classifier fitted on the retained ones, which is why its rows are marked (S+C) rather than (C).
All search spaces were explored with Optuna using the Gaussian-process sampler (\texttt{GPSampler}) with 20 trials per dataset-by-pipeline combination.
The objective maximized was the \emph{mean accuracy over the five inner cross-validation folds}, not AUC; the inner folds were generated by \texttt{StratifiedGroupKFold} grouped by \texttt{patient\_id}.
Settings that were fixed and not searched, one per line so that the long identifiers have somewhere to break:\\
\texttt{Lasso(fit\_intercept=True, max\_iter=5000)}\\
\texttt{LinearSVC(penalty='l1', max\_iter=5000, class\_weight='balanced')}\\
\texttt{NuSVC(kernel='rbf', probability=True, class\_weight='balanced')}\\
\texttt{XGBRFClassifier(random\_state=42)}\\
Class weighting was therefore applied in the two support vector branches only.
The models were seeded with \texttt{random\_state=42}, but \texttt{GPSampler()} was instantiated without a seed, so the sequence of 20 trials is not reproducible.
\par}

\clearpage

\begin{table}[htbp]
\centering
\caption{Cross-validated performance of all 24 dataset-by-pipeline combinations in the MR CLEAN NO-IV cohort ($n = 141$), over five outer folds. Feature set: \emph{global}, eight summary statistics of each radiomic feature over the thrombus, one row per patient; \emph{distribution}, a ten-bin histogram of each radiomic feature over the thrombus, one row per patient; \emph{voxel}, the radiomic feature values themselves, one row per voxel. ROI: \emph{whole}, the dilated thrombus mask; \emph{centerline}, the dilated skeleton of the thrombus mask. Selector (Sel.): LASSO, $\ell_1$-penalized regression; L1, embedded selection by the non-zero coefficients of the classifier itself; RFI, gain-based feature importance of an \texttt{XGBRFClassifier}. Classifier (Class.): XGBRF, \texttt{XGBRFClassifier}; LinSVC, linear support vector classifier with an $\ell_1$ penalty; NuSVC, $\nu$-support vector classifier with a radial basis function kernel. In the body of the table the distribution-based feature set is abbreviated \emph{distrib.}\ and the centerline ROI \emph{cent.} AUC, area under the receiver operating characteristic curve; ACC, accuracy; IQR, interquartile range across the five folds; min--max, the smallest and largest value across the five folds; Med.\ $n$, median number of features retained by the selector; Null AUC, mean AUC of the permutation baseline; $\kappa$, the $\kappa$-type index of the observed AUC relative to that baseline; $p$, permutation $p$-value; $q$, the same $p$-value after Benjamini--Hochberg correction across all 24 combinations.}
\label{tab:cv-performance}
\smallskip
\input{\tablesdir/table3_cv_performance}
\end{table}

\smallskip
{\footnotesize
\sloppy
\noindent
\textit{Notes to Table~\ref{tab:cv-performance}.}
Accuracy must be read against the majority-class accuracy of the development cohort, $82/141 = 0.582$; only 4 of the 24 combinations exceed it by median accuracy, and 86 of the 120 individual fold results (71.7\%) fall below it.
The class ratio itself varies across folds (0.448, 0.714, 0.536, 0.679, 0.536), so in the first fold the majority class is the absence rather than the presence of near-complete or complete reperfusion.
The largest single fold result of the whole study is 0.810, obtained by distribution/whole/RFI/NuSVC in fold 4.
Averaged over the eight whole-thrombus patient-level combinations, AUC per fold is 0.510, 0.661, 0.606, 0.587, 0.623, a spread within a single fold of up to 0.531; between-fold variability therefore exceeds between-method variability.
The permutation null was obtained from 100 label permutations at the patient level per outer fold; with $m = 100$ permutations the smallest attainable $p$-value is $p = (0+1)/(100+1) = 0.0099$, which bounds the resolution of every permutation result in this table.
The Null AUC, $\kappa$, $p$ and $q$ columns were computed from the \emph{mean} AUC across folds, whereas the AUC column reports the \emph{median}; the two summaries are not interchangeable and are shown side by side deliberately.
10 of the 24 combinations reach $p < 0.05$ uncorrected and 6 survive Benjamini--Hochberg correction\NOTETODO{cite Benjamini and Hochberg 1995; the key is missing from BibFile.bib} for AUC, while none survives it for accuracy.
No centerline combination survives correction; across the eight of them the mean AUC ranges from 0.396 to 0.495 against a null of 0.489--0.510, $\kappa$ from $-0.206$ to $-0.014$ and $p$ from 0.5347 to 1.0000.\\
\textsuperscript{a} In these 2 combinations the $\ell_1$ penalty shrank every coefficient to zero in every fold and the code substituted the index of the first column; the retained feature is a fallback, not one the selector identified as informative.
Combinations with a median of one retained feature but no mark reached that median with genuinely non-zero coefficients.\\
\textsuperscript{b} The unit of analysis of the voxel-based datasets is the voxel, not the patient, and their metrics are therefore \emph{not} comparable with the remaining rows.
Across the five folds the voxel/whole datasets contribute 396085 voxel-level predictions (91500, 79327, 75748, 82369, 67141) and the voxel/centerline datasets 8557 (1690, 1571, 1811, 1806, 1679), against 141 patient-level predictions (29, 28, 28, 28, 28) for every other dataset.
The consequence is visible in the permutation null, whose standard deviation is 0.040--0.052 for the voxel-based combinations against 0.063--0.120 for all others, so that the voxel-based $p$-values are computed against a spuriously narrow null and are not comparable with the rest of the table.
3 voxel-based combinations are formally significant for AUC after correction; they must not be read as evidence for a signal.
\NOTETODO{patient-level majority-vote aggregation of the voxel-based predictions, which would make these rows comparable}
\par}

\clearpage

\begin{table}[htbp]
\centering
\caption{External validation in the MR CLEAN Registry cohort. Each pipeline was refitted on all 141 development patients with its tuned hyperparameters and applied once to the Registry cohort. CV, cross-validated median over the five outer folds of the same combination in Table~\ref{tab:cv-performance}; Reg., Registry; ACC, accuracy; $\Delta$AUC, Registry AUC minus cross-validated median AUC, a negative value being a loss of discrimination on transfer; other abbreviations as in Table~\ref{tab:cv-performance}. The voxel-based datasets were not transferred externally and therefore have no row here, which is why the table has 16 rows rather than 24.}
\label{tab:external-validation}
\smallskip
\input{\tablesdir/table4_external_validation}
\end{table}

\smallskip
{\footnotesize
\sloppy
\noindent
\textit{Notes to Table~\ref{tab:external-validation}.}
Across the 16 models the Registry AUC ranged from 0.471 to 0.521 with a median of 0.496, and 5 of them exceeded 0.50.
Registry accuracy ranged from 0.440 to 0.550 with a median of 0.481, and no model reached the majority-class accuracy of the development cohort, 0.582.
The sign of $\Delta$ follows the cross-validated value and nothing else: the 8 combinations whose cross-validated median AUC was above 0.500 all moved downwards on transfer, by $-0.034$ to $-0.190$, and the 8 that were at or below it all moved upwards, by $0.002$ to $0.128$.
The upward movement is therefore a regression to chance of models that had performed below chance internally, not a gain in generalization, and it must not be presented as one.
The only combinations for which the comparison is informative are the 6 that discriminated in cross-validation, that is reached a cross-validated median AUC of at least 0.600; all of them lose between 0.150 and 0.190 of AUC.
\NOTEWARN{$n$ Registry is not stated in the results files and is not the same for every model. Reconstructing the smallest denominator compatible with the decimal expansion of each reported accuracy leaves exactly three cohort sizes: 1649 for eight models, 1760 for six and 1761 for two. Neither the 1906 of the exclusion cascade nor the 1839 identifiers stored with the cohort is compatible with any of the sixteen. Until this is resolved the AUC values in this table are not strictly comparable with one another}
\par}

\clearpage

\begin{table}[htbp]
\centering
\caption{Comparison with the reference pipeline of van Voorst et al. \cite{vanVoorst2022} The first row summarizes the present study; rows two to four were computed on the MR CLEAN NO-IV development cohort with the same five outer folds as in Table~\ref{tab:cv-performance}, so that the comparison is made within one dataset rather than across publications; the last row is the value reported in the original publication, on a different cohort and for a different endpoint. ICC, intraclass correlation coefficient: of the 107 features the reference pipeline extracts, the 52 with ICC $\geq$ 0.75 are retained and the rest discarded. AUC$_{\text{prob}}$ denotes the AUC computed from the predicted class probabilities, and is the quantity reported for the reference pipeline. Values are medians with the interquartile range in brackets, or means $\pm$ standard deviation across folds, as indicated in the note below. An em dash marks a quantity the row does not have: the pooled row spans combinations with different accuracies, and the published row reports neither a range nor an accuracy.}
\label{tab:reference-comparison}
\smallskip
\input{\tablesdir/table5_reference_comparison}
\end{table}

\smallskip
{\footnotesize
\sloppy
\noindent
\textit{Notes to Table~\ref{tab:reference-comparison}.}
The first row pools 25 fold results, 5 combinations $\times$ 5 outer folds, and its bracketed values are the interquartile range; the second row is a mean $\pm$ standard deviation over the 5 splits; rows three and four are medians over the 5 folds, with no further feature selection applied to the 52 ICC-filtered features.
The last row is quoted from the publication: thrombus radiomics for eTICI $\geq$ 2B, a more permissive endpoint than the eTICI 2C--3 outcome of this article, as a mean over 100 bootstrap samples of 100 patients.
The plain \texttt{AUC} column of the reference pipeline degenerates to exactly 0.50 in 4 of the five splits, because the fitted model assigns the majority class to every patient in those splits and its hard predictions therefore carry no ranking information.
AUC$_{\text{prob}}$, computed from the predicted probabilities, is the only quantity of that pipeline that is comparable with the AUC values reported elsewhere in this article, and it is the one used here.
By median accuracy the reference pipeline and both of our classifiers restricted to its 52 ICC-filtered features stay below the majority-class accuracy of 0.582, although the reference pipeline exceeds it in one individual split.
Rows two to four were recomputed against the eTICI 2C--3 outcome of this article --- the stored label column is named \texttt{TICI\_2B-3}, but its per-fold positive counts are exactly the 82/141 of this study --- and are directly comparable with the first row, whereas the published 0.52 of the last row was obtained for eTICI $\geq$ 2B in a different cohort and differs from them in endpoint and cohort as much as in pipeline.
\NOTEQUERY{the pool of the first row omits one whole-thrombus combination that also discriminated in cross-validation, and the reason is not documented in the source material}
\NOTEQUERY{decide which of the two mutually incompatible sets of numbers in the source data applies before this table is finalized: set A, 0.481 $\pm$ 0.145 on 141 patients against set B, 0.467 $\pm$ 0.099 on 170}
\par}

\clearpage

\begin{table}[htbp]
\centering
\caption{Ensemble analysis: the 15 configurations of discrete, that is binary, AdaBoost \cite{Freund1997} applied to the stored prediction vectors of the 16 patient-level weak classifiers (four patient-level datasets $\times$ four pipelines). Each configuration was fitted on the five outer folds, giving 75 ensemble models. ACC $\varnothing$, mean accuracy of the configuration across the five folds; best weak classifier, the single pipeline in the configuration's pool with the highest mean accuracy, and its value; $\Delta$, the mean accuracy of the ensemble minus that of the best weak classifier. A negative $\Delta$ means the ensemble is worse than the best model it combines. The configuration name states which of the four patient-level datasets form its pool, each contributing all four pipelines: \texttt{type1} pools all four datasets (16 weak classifiers), \texttt{type2} a whole-thrombus dataset with its own centerline counterpart (8), \texttt{type3} a single dataset (4), \texttt{type4} two datasets that are not such a pair (8) and \texttt{type5} all datasets but one (12). \texttt{distrib.}\ abbreviates the distribution-based dataset and \emph{cent.} the centerline region of interest, as in Table~\ref{tab:cv-performance}. The two voxel-based datasets were excluded from every pool because their predictions are defined per voxel.}
\label{tab:ensemble}
\smallskip
\input{\tablesdir/table6_ensemble}
\end{table}

\smallskip
{\footnotesize
\sloppy
\noindent
\textit{Notes to Table~\ref{tab:ensemble}.}
$\Delta$ is negative for all 15 configurations, between $-0.022$ and $-0.121$; no ensemble outperformed the best single pipeline in its own pool.
The worst result belongs to \texttt{type1\_\allowbreak all\_\allowbreak featuresets}, the configuration that pools all feature sets.
Across all 75 fitted ensemble models accuracy averages $0.543 \pm 0.083$ with a range of 0.414--0.714.
The mechanism is documented by the fitted weights: in the all-feature-set configuration, 55 of the 80 boosting rounds (16 rounds $\times$ five folds) received a negative weight $\alpha$, meaning that the weak classifier admitted in that round performed worse than chance on the reweighted data.
The individual pipelines therefore do not make complementary errors, and the boosting assumption of weak classifiers reliably better than chance is not met.
Labels were coded as $\{-1,+1\}$, the weight of a round was $\alpha = \tfrac{1}{2}\ln\frac{1-\varepsilon}{\varepsilon}$ with $\varepsilon$ the error of the admitted classifier on the currently reweighted data, and each weak classifier could be admitted at most once, so a pool of $k$ members yields at most $k$ rounds.
The weak classifiers were not retrained: the procedure is applied to models that had already been fitted, which makes it methodologically a form of stacking rather than boosting proper.
The best weak classifier of \texttt{type3\_\allowbreak global} and of \texttt{type1\_\allowbreak all\_\allowbreak featuresets} is the same pipeline, because the pool of the latter contains the pool of the former.
\par}

\clearpage

\begin{table}[htbp]
\centering
\caption{Pairwise comparison of the pipelines by the corrected resampled $t$-test of Nadeau and Bengio \cite{Nadeau2003}. The upper block compares the four pipelines of this study with one another on the whole-thrombus global-based dataset; the lower block compares them, and our two classifiers restricted to the reference pipeline's own 52 ICC-filtered features, against the reference pipeline of van Voorst et al.\ recomputed on the same cohort. Every model in the table was evaluated on the same 5 outer folds of the same 141 patients, which is what makes the paired test applicable. Because the training sets of $k$-fold cross-validation overlap, the per-fold differences are not independent and the ordinary paired $t$-test underestimates their variance; the correction multiplies the variance by $1/k + n_{\text{test}}/n_{\text{train}}$, here $1/5 + 28.2/112.8 = 0.4500$, which inflates the standard error 1.50-fold at $\mathrm{df} = 4$. Fold sizes are taken from the data, not assumed. Mean $\Delta$AUC is the mean of the 5 per-fold differences in AUC between pipeline A and pipeline B; a positive value favours A. $t$ (uncorrected) is the ordinary paired statistic, $t$ (corrected) the same quantity with the inflated variance. The two-sided critical value at $\alpha = 0.05$ and $\mathrm{df} = 4$ is 2.776. Pipeline names combine selector and classifier as in Table~\ref{tab:cv-performance}.}
\label{tab:nadeau-bengio}
\smallskip
\input{\tablesdir/table7_nadeau_bengio}
\end{table}

\smallskip
{\footnotesize
\sloppy
\noindent
\textit{Notes to Table~\ref{tab:nadeau-bengio}.}
Not one of the 12 comparisons reaches significance: the largest corrected statistic in absolute value is 1.54, against a critical value of 2.776.
The four pipelines of this study are therefore statistically indistinguishable from one another, and none of them is distinguishable from the reference pipeline either, in either direction.
For the reference pipeline the compared quantity is AUC$_{\text{prob}}$, because its plain \texttt{AUC} column degenerates to exactly 0.50 in four of the five splits; see Table~\ref{tab:reference-comparison}.
The pooled row of Table~\ref{tab:reference-comparison} is deliberately absent here: those 5 combinations were picked by their results on these same folds, so testing them against the reference pipeline would be biased in favour of the present study. All four pipelines of the dataset are tested instead, whatever their outcome.
With 5 folds the test has $\mathrm{df} = 4$ and very little power, so these results are an absence of evidence for a difference, not evidence that the pipelines perform identically. The comparison that the data would support best, a DeLong test on the external ROC curves, was not possible because per-patient predicted scores on the Registry were not retained.
\par}

\clearpage


%%
%% Tabulky se vejdou i do uzke sazby `elsarticle` (390 pt), ale zdrobni.
%% Aby vypadaly jako v samostatnem supplementu, obal se vlozeni geometrii:
%%
%%     \usepackage{geometry}          % v preambuli, za elsarticle
%%     ...
%%     \clearpage
%%     \newgeometry{margin=15mm}
%%     %% -------------------------------------------------------------------------
%% GENEROVANO: make_all.py -- needitovat rucne.
%%
%% Popisky, tabulky a poznamky bez preambule. Do clanku staci jeden radek:
%%
%%     \input{tables_tom/output/supplement_body}
%%
%% Tabulky se vejdou i do uzke sazby `elsarticle` (390 pt), ale zdrobni.
%% Aby vypadaly jako v samostatnem supplementu, obal se vlozeni geometrii:
%%
%%     \usepackage{geometry}          % v preambuli, za elsarticle
%%     ...
%%     \clearpage
%%     \newgeometry{margin=15mm}
%%     \input{tables_tom/output/supplement_body}
%%     \restoregeometry
%%
%% Cislovani S1--S7 misto 1--7 zaridi pred vlozenim:
%%     \setcounter{table}{0}\renewcommand{\thetable}{S\arabic{table}}
%%
%% Preambule potrebuje: booktabs, array, amsmath, amssymb.
%% Cesta k tabulkam se bere z \tablesdir; vychozi hodnota plati pro preklad
%% z korene projektu.
%% -------------------------------------------------------------------------
\providecommand{\tablesdir}{tables_tom/output}

%% Zaloha, kdyby `manuscriptnotes` nebyl nactveny: vysadi se jen kotva,
%% poznamka se zahodi -- tedy tentyz vysledek jako v rezimu odesilani.
\providecommand{\NOTE}[2][]{#1}
\providecommand{\NOTEWARN}[2][]{#1}
\providecommand{\NOTECHECK}[2][]{#1}
\providecommand{\NOTETODO}[2][]{#1}
\providecommand{\NOTEQUERY}[2][]{#1}

{\large\bfseries Supplementary tables\par}
\bigskip
\NOTEWARN{the captions and notes in this supplement were drafted with AI assistance and are provisional; every one of them has to be rewritten and checked by an author before submission}

\begin{table}[htbp]
\centering
\caption{Patient selection in the two cohorts. MR CLEAN NO-IV was used for model development and cross-validation, the MR CLEAN Registry as the independent external cohort. Entries are numbers of patients; an em dash marks an exclusion step that does not apply to that cohort. eTICI, expanded Thrombolysis in Cerebral Infarction score; the outcome modeled throughout this article is near-complete or complete reperfusion, defined as a final eTICI score of 2C or 3. Slice spacing refers to the through-plane spacing of the non-contrast CT acquisition.}
\label{tab:cohorts}
\smallskip
\input{\tablesdir/table1_cohorts}
\end{table}

\smallskip
{\footnotesize
\sloppy
\noindent
\textit{Notes to Table~\ref{tab:cohorts}.}
The manual review of segmentation quality was carried out on the MR CLEAN NO-IV cohort only.
No segmentation was reviewed or rejected in the MR CLEAN Registry cohort, so the two cohorts differ in segmentation quality control as well as in population.
This is a property of the design and is the single largest asymmetry between the development and the external cohort.
The review was blinded to the eTICI score: it judged the quality and the anatomical placement of the segmentation only, so the 155 exclusions it produced cannot have been informed by the outcome.
Of the 141 patients that remain, 82 (58.2\%) reached eTICI 2C--3; the proportion varies across the five outer folds (0.448, 0.714, 0.536, 0.679, 0.536), so in the first fold the majority class is the absence rather than the presence of near-complete or complete reperfusion.
\NOTETODO{breakdown of how many of the 155 reviewed cases were rejected by which criterion}
\NOTETODO{number and proportion of patients with eTICI 2C--3 reperfusion in the Registry cohort, which the table leaves as an em dash}
\NOTEWARN{size of the Registry cohort: the exclusion cascade gives 1906, the stored identifier list holds 1839, and the denominators reconstructed from the decimal expansion of the reported external accuracies are 1649 (8 models), 1760 (6 models), 1761 (2 models). Neither 1906 nor 1839 is compatible with any of the sixteen transferred models, and the reconstructed values are mutually incompatible, so the external cohort is not the same across feature sets}
\par}

\clearpage

\begin{table}[htbp]
\centering
\caption{Hyperparameter search spaces of the four pipelines. LASSO, $\ell_1$-penalized linear regression on the binary outcome used as a feature selector; RFI, ranking by the gain-based feature importance of a random-forest gradient-boosting model; XGBRF, \texttt{XGBRFClassifier}; LinearSVC, linear support vector classifier with an $\ell_1$ penalty, which acts as selector and classifier simultaneously; NuSVC, $\nu$-support vector classifier with a radial basis function kernel. In the Component column, (S) marks the component acting as feature selector, (C) as classifier, and (S+C) a single object that performs both roles. Scale indicates how Optuna sampled the parameter: \emph{log}, log-uniform over the interval; \emph{int}, uniform over the integers of the interval; \emph{float}, uniform over the interval.}
\label{tab:hyperparameters}
\smallskip
\input{\tablesdir/table2_hyperparameters}
\end{table}

\smallskip
{\footnotesize
\sloppy
\noindent
\textit{Notes to Table~\ref{tab:hyperparameters}.}
$n$ and \texttt{X.shape[1]} denote the number of rows and the number of columns of the training matrix passed to the search.
Rows are patients for the global-based and distribution-based datasets and voxels for the voxel-based datasets, so the same expression yields different intervals across datasets; the number of columns reaches 20660 for the distribution-based datasets.
The minority fraction is the proportion of the smaller outcome class among the training labels of the given fold.
In the RFI--RF pipeline a single parameter dictionary configures both the model whose gain importances rank the features and the classifier fitted on the retained ones, which is why its rows are marked (S+C) rather than (C).
All search spaces were explored with Optuna using the Gaussian-process sampler (\texttt{GPSampler}) with 20 trials per dataset-by-pipeline combination.
The objective maximized was the \emph{mean accuracy over the five inner cross-validation folds}, not AUC; the inner folds were generated by \texttt{StratifiedGroupKFold} grouped by \texttt{patient\_id}.
Settings that were fixed and not searched, one per line so that the long identifiers have somewhere to break:\\
\texttt{Lasso(fit\_intercept=True, max\_iter=5000)}\\
\texttt{LinearSVC(penalty='l1', max\_iter=5000, class\_weight='balanced')}\\
\texttt{NuSVC(kernel='rbf', probability=True, class\_weight='balanced')}\\
\texttt{XGBRFClassifier(random\_state=42)}\\
Class weighting was therefore applied in the two support vector branches only.
The models were seeded with \texttt{random\_state=42}, but \texttt{GPSampler()} was instantiated without a seed, so the sequence of 20 trials is not reproducible.
\par}

\clearpage

\begin{table}[htbp]
\centering
\caption{Cross-validated performance of all 24 dataset-by-pipeline combinations in the MR CLEAN NO-IV cohort ($n = 141$), over five outer folds. Feature set: \emph{global}, eight summary statistics of each radiomic feature over the thrombus, one row per patient; \emph{distribution}, a ten-bin histogram of each radiomic feature over the thrombus, one row per patient; \emph{voxel}, the radiomic feature values themselves, one row per voxel. ROI: \emph{whole}, the dilated thrombus mask; \emph{centerline}, the dilated skeleton of the thrombus mask. Selector (Sel.): LASSO, $\ell_1$-penalized regression; L1, embedded selection by the non-zero coefficients of the classifier itself; RFI, gain-based feature importance of an \texttt{XGBRFClassifier}. Classifier (Class.): XGBRF, \texttt{XGBRFClassifier}; LinSVC, linear support vector classifier with an $\ell_1$ penalty; NuSVC, $\nu$-support vector classifier with a radial basis function kernel. In the body of the table the distribution-based feature set is abbreviated \emph{distrib.}\ and the centerline ROI \emph{cent.} AUC, area under the receiver operating characteristic curve; ACC, accuracy; IQR, interquartile range across the five folds; min--max, the smallest and largest value across the five folds; Med.\ $n$, median number of features retained by the selector; Null AUC, mean AUC of the permutation baseline; $\kappa$, the $\kappa$-type index of the observed AUC relative to that baseline; $p$, permutation $p$-value; $q$, the same $p$-value after Benjamini--Hochberg correction across all 24 combinations.}
\label{tab:cv-performance}
\smallskip
\input{\tablesdir/table3_cv_performance}
\end{table}

\smallskip
{\footnotesize
\sloppy
\noindent
\textit{Notes to Table~\ref{tab:cv-performance}.}
Accuracy must be read against the majority-class accuracy of the development cohort, $82/141 = 0.582$; only 4 of the 24 combinations exceed it by median accuracy, and 86 of the 120 individual fold results (71.7\%) fall below it.
The class ratio itself varies across folds (0.448, 0.714, 0.536, 0.679, 0.536), so in the first fold the majority class is the absence rather than the presence of near-complete or complete reperfusion.
The largest single fold result of the whole study is 0.810, obtained by distribution/whole/RFI/NuSVC in fold 4.
Averaged over the eight whole-thrombus patient-level combinations, AUC per fold is 0.510, 0.661, 0.606, 0.587, 0.623, a spread within a single fold of up to 0.531; between-fold variability therefore exceeds between-method variability.
The permutation null was obtained from 100 label permutations at the patient level per outer fold; with $m = 100$ permutations the smallest attainable $p$-value is $p = (0+1)/(100+1) = 0.0099$, which bounds the resolution of every permutation result in this table.
The Null AUC, $\kappa$, $p$ and $q$ columns were computed from the \emph{mean} AUC across folds, whereas the AUC column reports the \emph{median}; the two summaries are not interchangeable and are shown side by side deliberately.
10 of the 24 combinations reach $p < 0.05$ uncorrected and 6 survive Benjamini--Hochberg correction\NOTETODO{cite Benjamini and Hochberg 1995; the key is missing from BibFile.bib} for AUC, while none survives it for accuracy.
No centerline combination survives correction; across the eight of them the mean AUC ranges from 0.396 to 0.495 against a null of 0.489--0.510, $\kappa$ from $-0.206$ to $-0.014$ and $p$ from 0.5347 to 1.0000.\\
\textsuperscript{a} In these 2 combinations the $\ell_1$ penalty shrank every coefficient to zero in every fold and the code substituted the index of the first column; the retained feature is a fallback, not one the selector identified as informative.
Combinations with a median of one retained feature but no mark reached that median with genuinely non-zero coefficients.\\
\textsuperscript{b} The unit of analysis of the voxel-based datasets is the voxel, not the patient, and their metrics are therefore \emph{not} comparable with the remaining rows.
Across the five folds the voxel/whole datasets contribute 396085 voxel-level predictions (91500, 79327, 75748, 82369, 67141) and the voxel/centerline datasets 8557 (1690, 1571, 1811, 1806, 1679), against 141 patient-level predictions (29, 28, 28, 28, 28) for every other dataset.
The consequence is visible in the permutation null, whose standard deviation is 0.040--0.052 for the voxel-based combinations against 0.063--0.120 for all others, so that the voxel-based $p$-values are computed against a spuriously narrow null and are not comparable with the rest of the table.
3 voxel-based combinations are formally significant for AUC after correction; they must not be read as evidence for a signal.
\NOTETODO{patient-level majority-vote aggregation of the voxel-based predictions, which would make these rows comparable}
\par}

\clearpage

\begin{table}[htbp]
\centering
\caption{External validation in the MR CLEAN Registry cohort. Each pipeline was refitted on all 141 development patients with its tuned hyperparameters and applied once to the Registry cohort. CV, cross-validated median over the five outer folds of the same combination in Table~\ref{tab:cv-performance}; Reg., Registry; ACC, accuracy; $\Delta$AUC, Registry AUC minus cross-validated median AUC, a negative value being a loss of discrimination on transfer; other abbreviations as in Table~\ref{tab:cv-performance}. The voxel-based datasets were not transferred externally and therefore have no row here, which is why the table has 16 rows rather than 24.}
\label{tab:external-validation}
\smallskip
\input{\tablesdir/table4_external_validation}
\end{table}

\smallskip
{\footnotesize
\sloppy
\noindent
\textit{Notes to Table~\ref{tab:external-validation}.}
Across the 16 models the Registry AUC ranged from 0.471 to 0.521 with a median of 0.496, and 5 of them exceeded 0.50.
Registry accuracy ranged from 0.440 to 0.550 with a median of 0.481, and no model reached the majority-class accuracy of the development cohort, 0.582.
The sign of $\Delta$ follows the cross-validated value and nothing else: the 8 combinations whose cross-validated median AUC was above 0.500 all moved downwards on transfer, by $-0.034$ to $-0.190$, and the 8 that were at or below it all moved upwards, by $0.002$ to $0.128$.
The upward movement is therefore a regression to chance of models that had performed below chance internally, not a gain in generalization, and it must not be presented as one.
The only combinations for which the comparison is informative are the 6 that discriminated in cross-validation, that is reached a cross-validated median AUC of at least 0.600; all of them lose between 0.150 and 0.190 of AUC.
\NOTEWARN{$n$ Registry is not stated in the results files and is not the same for every model. Reconstructing the smallest denominator compatible with the decimal expansion of each reported accuracy leaves exactly three cohort sizes: 1649 for eight models, 1760 for six and 1761 for two. Neither the 1906 of the exclusion cascade nor the 1839 identifiers stored with the cohort is compatible with any of the sixteen. Until this is resolved the AUC values in this table are not strictly comparable with one another}
\par}

\clearpage

\begin{table}[htbp]
\centering
\caption{Comparison with the reference pipeline of van Voorst et al. \cite{vanVoorst2022} The first row summarizes the present study; rows two to four were computed on the MR CLEAN NO-IV development cohort with the same five outer folds as in Table~\ref{tab:cv-performance}, so that the comparison is made within one dataset rather than across publications; the last row is the value reported in the original publication, on a different cohort and for a different endpoint. ICC, intraclass correlation coefficient: of the 107 features the reference pipeline extracts, the 52 with ICC $\geq$ 0.75 are retained and the rest discarded. AUC$_{\text{prob}}$ denotes the AUC computed from the predicted class probabilities, and is the quantity reported for the reference pipeline. Values are medians with the interquartile range in brackets, or means $\pm$ standard deviation across folds, as indicated in the note below. An em dash marks a quantity the row does not have: the pooled row spans combinations with different accuracies, and the published row reports neither a range nor an accuracy.}
\label{tab:reference-comparison}
\smallskip
\input{\tablesdir/table5_reference_comparison}
\end{table}

\smallskip
{\footnotesize
\sloppy
\noindent
\textit{Notes to Table~\ref{tab:reference-comparison}.}
The first row pools 25 fold results, 5 combinations $\times$ 5 outer folds, and its bracketed values are the interquartile range; the second row is a mean $\pm$ standard deviation over the 5 splits; rows three and four are medians over the 5 folds, with no further feature selection applied to the 52 ICC-filtered features.
The last row is quoted from the publication: thrombus radiomics for eTICI $\geq$ 2B, a more permissive endpoint than the eTICI 2C--3 outcome of this article, as a mean over 100 bootstrap samples of 100 patients.
The plain \texttt{AUC} column of the reference pipeline degenerates to exactly 0.50 in 4 of the five splits, because the fitted model assigns the majority class to every patient in those splits and its hard predictions therefore carry no ranking information.
AUC$_{\text{prob}}$, computed from the predicted probabilities, is the only quantity of that pipeline that is comparable with the AUC values reported elsewhere in this article, and it is the one used here.
By median accuracy the reference pipeline and both of our classifiers restricted to its 52 ICC-filtered features stay below the majority-class accuracy of 0.582, although the reference pipeline exceeds it in one individual split.
Rows two to four were recomputed against the eTICI 2C--3 outcome of this article --- the stored label column is named \texttt{TICI\_2B-3}, but its per-fold positive counts are exactly the 82/141 of this study --- and are directly comparable with the first row, whereas the published 0.52 of the last row was obtained for eTICI $\geq$ 2B in a different cohort and differs from them in endpoint and cohort as much as in pipeline.
\NOTEQUERY{the pool of the first row omits one whole-thrombus combination that also discriminated in cross-validation, and the reason is not documented in the source material}
\NOTEQUERY{decide which of the two mutually incompatible sets of numbers in the source data applies before this table is finalized: set A, 0.481 $\pm$ 0.145 on 141 patients against set B, 0.467 $\pm$ 0.099 on 170}
\par}

\clearpage

\begin{table}[htbp]
\centering
\caption{Ensemble analysis: the 15 configurations of discrete, that is binary, AdaBoost \cite{Freund1997} applied to the stored prediction vectors of the 16 patient-level weak classifiers (four patient-level datasets $\times$ four pipelines). Each configuration was fitted on the five outer folds, giving 75 ensemble models. ACC $\varnothing$, mean accuracy of the configuration across the five folds; best weak classifier, the single pipeline in the configuration's pool with the highest mean accuracy, and its value; $\Delta$, the mean accuracy of the ensemble minus that of the best weak classifier. A negative $\Delta$ means the ensemble is worse than the best model it combines. The configuration name states which of the four patient-level datasets form its pool, each contributing all four pipelines: \texttt{type1} pools all four datasets (16 weak classifiers), \texttt{type2} a whole-thrombus dataset with its own centerline counterpart (8), \texttt{type3} a single dataset (4), \texttt{type4} two datasets that are not such a pair (8) and \texttt{type5} all datasets but one (12). \texttt{distrib.}\ abbreviates the distribution-based dataset and \emph{cent.} the centerline region of interest, as in Table~\ref{tab:cv-performance}. The two voxel-based datasets were excluded from every pool because their predictions are defined per voxel.}
\label{tab:ensemble}
\smallskip
\input{\tablesdir/table6_ensemble}
\end{table}

\smallskip
{\footnotesize
\sloppy
\noindent
\textit{Notes to Table~\ref{tab:ensemble}.}
$\Delta$ is negative for all 15 configurations, between $-0.022$ and $-0.121$; no ensemble outperformed the best single pipeline in its own pool.
The worst result belongs to \texttt{type1\_\allowbreak all\_\allowbreak featuresets}, the configuration that pools all feature sets.
Across all 75 fitted ensemble models accuracy averages $0.543 \pm 0.083$ with a range of 0.414--0.714.
The mechanism is documented by the fitted weights: in the all-feature-set configuration, 55 of the 80 boosting rounds (16 rounds $\times$ five folds) received a negative weight $\alpha$, meaning that the weak classifier admitted in that round performed worse than chance on the reweighted data.
The individual pipelines therefore do not make complementary errors, and the boosting assumption of weak classifiers reliably better than chance is not met.
Labels were coded as $\{-1,+1\}$, the weight of a round was $\alpha = \tfrac{1}{2}\ln\frac{1-\varepsilon}{\varepsilon}$ with $\varepsilon$ the error of the admitted classifier on the currently reweighted data, and each weak classifier could be admitted at most once, so a pool of $k$ members yields at most $k$ rounds.
The weak classifiers were not retrained: the procedure is applied to models that had already been fitted, which makes it methodologically a form of stacking rather than boosting proper.
The best weak classifier of \texttt{type3\_\allowbreak global} and of \texttt{type1\_\allowbreak all\_\allowbreak featuresets} is the same pipeline, because the pool of the latter contains the pool of the former.
\par}

\clearpage

\begin{table}[htbp]
\centering
\caption{Pairwise comparison of the pipelines by the corrected resampled $t$-test of Nadeau and Bengio \cite{Nadeau2003}. The upper block compares the four pipelines of this study with one another on the whole-thrombus global-based dataset; the lower block compares them, and our two classifiers restricted to the reference pipeline's own 52 ICC-filtered features, against the reference pipeline of van Voorst et al.\ recomputed on the same cohort. Every model in the table was evaluated on the same 5 outer folds of the same 141 patients, which is what makes the paired test applicable. Because the training sets of $k$-fold cross-validation overlap, the per-fold differences are not independent and the ordinary paired $t$-test underestimates their variance; the correction multiplies the variance by $1/k + n_{\text{test}}/n_{\text{train}}$, here $1/5 + 28.2/112.8 = 0.4500$, which inflates the standard error 1.50-fold at $\mathrm{df} = 4$. Fold sizes are taken from the data, not assumed. Mean $\Delta$AUC is the mean of the 5 per-fold differences in AUC between pipeline A and pipeline B; a positive value favours A. $t$ (uncorrected) is the ordinary paired statistic, $t$ (corrected) the same quantity with the inflated variance. The two-sided critical value at $\alpha = 0.05$ and $\mathrm{df} = 4$ is 2.776. Pipeline names combine selector and classifier as in Table~\ref{tab:cv-performance}.}
\label{tab:nadeau-bengio}
\smallskip
\input{\tablesdir/table7_nadeau_bengio}
\end{table}

\smallskip
{\footnotesize
\sloppy
\noindent
\textit{Notes to Table~\ref{tab:nadeau-bengio}.}
Not one of the 12 comparisons reaches significance: the largest corrected statistic in absolute value is 1.54, against a critical value of 2.776.
The four pipelines of this study are therefore statistically indistinguishable from one another, and none of them is distinguishable from the reference pipeline either, in either direction.
For the reference pipeline the compared quantity is AUC$_{\text{prob}}$, because its plain \texttt{AUC} column degenerates to exactly 0.50 in four of the five splits; see Table~\ref{tab:reference-comparison}.
The pooled row of Table~\ref{tab:reference-comparison} is deliberately absent here: those 5 combinations were picked by their results on these same folds, so testing them against the reference pipeline would be biased in favour of the present study. All four pipelines of the dataset are tested instead, whatever their outcome.
With 5 folds the test has $\mathrm{df} = 4$ and very little power, so these results are an absence of evidence for a difference, not evidence that the pipelines perform identically. The comparison that the data would support best, a DeLong test on the external ROC curves, was not possible because per-patient predicted scores on the Registry were not retained.
\par}

\clearpage


%%     \restoregeometry
%%
%% Cislovani S1--S7 misto 1--7 zaridi pred vlozenim:
%%     \setcounter{table}{0}\renewcommand{\thetable}{S\arabic{table}}
%%
%% Preambule potrebuje: booktabs, array, amsmath, amssymb.
%% Cesta k tabulkam se bere z \tablesdir; vychozi hodnota plati pro preklad
%% z korene projektu.
%% -------------------------------------------------------------------------
\providecommand{\tablesdir}{tables_tom/output}

%% Zaloha, kdyby `manuscriptnotes` nebyl nactveny: vysadi se jen kotva,
%% poznamka se zahodi -- tedy tentyz vysledek jako v rezimu odesilani.
\providecommand{\NOTE}[2][]{#1}
\providecommand{\NOTEWARN}[2][]{#1}
\providecommand{\NOTECHECK}[2][]{#1}
\providecommand{\NOTETODO}[2][]{#1}
\providecommand{\NOTEQUERY}[2][]{#1}

{\large\bfseries Supplementary tables\par}
\bigskip
\NOTEWARN{the captions and notes in this supplement were drafted with AI assistance and are provisional; every one of them has to be rewritten and checked by an author before submission}

\begin{table}[htbp]
\centering
\caption{Patient selection in the two cohorts. MR CLEAN NO-IV was used for model development and cross-validation, the MR CLEAN Registry as the independent external cohort. Entries are numbers of patients; an em dash marks an exclusion step that does not apply to that cohort. eTICI, expanded Thrombolysis in Cerebral Infarction score; the outcome modeled throughout this article is near-complete or complete reperfusion, defined as a final eTICI score of 2C or 3. Slice spacing refers to the through-plane spacing of the non-contrast CT acquisition.}
\label{tab:cohorts}
\smallskip
\input{\tablesdir/table1_cohorts}
\end{table}

\smallskip
{\footnotesize
\sloppy
\noindent
\textit{Notes to Table~\ref{tab:cohorts}.}
The manual review of segmentation quality was carried out on the MR CLEAN NO-IV cohort only.
No segmentation was reviewed or rejected in the MR CLEAN Registry cohort, so the two cohorts differ in segmentation quality control as well as in population.
This is a property of the design and is the single largest asymmetry between the development and the external cohort.
The review was blinded to the eTICI score: it judged the quality and the anatomical placement of the segmentation only, so the 155 exclusions it produced cannot have been informed by the outcome.
Of the 141 patients that remain, 82 (58.2\%) reached eTICI 2C--3; the proportion varies across the five outer folds (0.448, 0.714, 0.536, 0.679, 0.536), so in the first fold the majority class is the absence rather than the presence of near-complete or complete reperfusion.
\NOTETODO{breakdown of how many of the 155 reviewed cases were rejected by which criterion}
\NOTETODO{number and proportion of patients with eTICI 2C--3 reperfusion in the Registry cohort, which the table leaves as an em dash}
\NOTEWARN{size of the Registry cohort: the exclusion cascade gives 1906, the stored identifier list holds 1839, and the denominators reconstructed from the decimal expansion of the reported external accuracies are 1649 (8 models), 1760 (6 models), 1761 (2 models). Neither 1906 nor 1839 is compatible with any of the sixteen transferred models, and the reconstructed values are mutually incompatible, so the external cohort is not the same across feature sets}
\par}

\clearpage

\begin{table}[htbp]
\centering
\caption{Hyperparameter search spaces of the four pipelines. LASSO, $\ell_1$-penalized linear regression on the binary outcome used as a feature selector; RFI, ranking by the gain-based feature importance of a random-forest gradient-boosting model; XGBRF, \texttt{XGBRFClassifier}; LinearSVC, linear support vector classifier with an $\ell_1$ penalty, which acts as selector and classifier simultaneously; NuSVC, $\nu$-support vector classifier with a radial basis function kernel. In the Component column, (S) marks the component acting as feature selector, (C) as classifier, and (S+C) a single object that performs both roles. Scale indicates how Optuna sampled the parameter: \emph{log}, log-uniform over the interval; \emph{int}, uniform over the integers of the interval; \emph{float}, uniform over the interval.}
\label{tab:hyperparameters}
\smallskip
\input{\tablesdir/table2_hyperparameters}
\end{table}

\smallskip
{\footnotesize
\sloppy
\noindent
\textit{Notes to Table~\ref{tab:hyperparameters}.}
$n$ and \texttt{X.shape[1]} denote the number of rows and the number of columns of the training matrix passed to the search.
Rows are patients for the global-based and distribution-based datasets and voxels for the voxel-based datasets, so the same expression yields different intervals across datasets; the number of columns reaches 20660 for the distribution-based datasets.
The minority fraction is the proportion of the smaller outcome class among the training labels of the given fold.
In the RFI--RF pipeline a single parameter dictionary configures both the model whose gain importances rank the features and the classifier fitted on the retained ones, which is why its rows are marked (S+C) rather than (C).
All search spaces were explored with Optuna using the Gaussian-process sampler (\texttt{GPSampler}) with 20 trials per dataset-by-pipeline combination.
The objective maximized was the \emph{mean accuracy over the five inner cross-validation folds}, not AUC; the inner folds were generated by \texttt{StratifiedGroupKFold} grouped by \texttt{patient\_id}.
Settings that were fixed and not searched, one per line so that the long identifiers have somewhere to break:\\
\texttt{Lasso(fit\_intercept=True, max\_iter=5000)}\\
\texttt{LinearSVC(penalty='l1', max\_iter=5000, class\_weight='balanced')}\\
\texttt{NuSVC(kernel='rbf', probability=True, class\_weight='balanced')}\\
\texttt{XGBRFClassifier(random\_state=42)}\\
Class weighting was therefore applied in the two support vector branches only.
The models were seeded with \texttt{random\_state=42}, but \texttt{GPSampler()} was instantiated without a seed, so the sequence of 20 trials is not reproducible.
\par}

\clearpage

\begin{table}[htbp]
\centering
\caption{Cross-validated performance of all 24 dataset-by-pipeline combinations in the MR CLEAN NO-IV cohort ($n = 141$), over five outer folds. Feature set: \emph{global}, eight summary statistics of each radiomic feature over the thrombus, one row per patient; \emph{distribution}, a ten-bin histogram of each radiomic feature over the thrombus, one row per patient; \emph{voxel}, the radiomic feature values themselves, one row per voxel. ROI: \emph{whole}, the dilated thrombus mask; \emph{centerline}, the dilated skeleton of the thrombus mask. Selector (Sel.): LASSO, $\ell_1$-penalized regression; L1, embedded selection by the non-zero coefficients of the classifier itself; RFI, gain-based feature importance of an \texttt{XGBRFClassifier}. Classifier (Class.): XGBRF, \texttt{XGBRFClassifier}; LinSVC, linear support vector classifier with an $\ell_1$ penalty; NuSVC, $\nu$-support vector classifier with a radial basis function kernel. In the body of the table the distribution-based feature set is abbreviated \emph{distrib.}\ and the centerline ROI \emph{cent.} AUC, area under the receiver operating characteristic curve; ACC, accuracy; IQR, interquartile range across the five folds; min--max, the smallest and largest value across the five folds; Med.\ $n$, median number of features retained by the selector; Null AUC, mean AUC of the permutation baseline; $\kappa$, the $\kappa$-type index of the observed AUC relative to that baseline; $p$, permutation $p$-value; $q$, the same $p$-value after Benjamini--Hochberg correction across all 24 combinations.}
\label{tab:cv-performance}
\smallskip
\input{\tablesdir/table3_cv_performance}
\end{table}

\smallskip
{\footnotesize
\sloppy
\noindent
\textit{Notes to Table~\ref{tab:cv-performance}.}
Accuracy must be read against the majority-class accuracy of the development cohort, $82/141 = 0.582$; only 4 of the 24 combinations exceed it by median accuracy, and 86 of the 120 individual fold results (71.7\%) fall below it.
The class ratio itself varies across folds (0.448, 0.714, 0.536, 0.679, 0.536), so in the first fold the majority class is the absence rather than the presence of near-complete or complete reperfusion.
The largest single fold result of the whole study is 0.810, obtained by distribution/whole/RFI/NuSVC in fold 4.
Averaged over the eight whole-thrombus patient-level combinations, AUC per fold is 0.510, 0.661, 0.606, 0.587, 0.623, a spread within a single fold of up to 0.531; between-fold variability therefore exceeds between-method variability.
The permutation null was obtained from 100 label permutations at the patient level per outer fold; with $m = 100$ permutations the smallest attainable $p$-value is $p = (0+1)/(100+1) = 0.0099$, which bounds the resolution of every permutation result in this table.
The Null AUC, $\kappa$, $p$ and $q$ columns were computed from the \emph{mean} AUC across folds, whereas the AUC column reports the \emph{median}; the two summaries are not interchangeable and are shown side by side deliberately.
10 of the 24 combinations reach $p < 0.05$ uncorrected and 6 survive Benjamini--Hochberg correction\NOTETODO{cite Benjamini and Hochberg 1995; the key is missing from BibFile.bib} for AUC, while none survives it for accuracy.
No centerline combination survives correction; across the eight of them the mean AUC ranges from 0.396 to 0.495 against a null of 0.489--0.510, $\kappa$ from $-0.206$ to $-0.014$ and $p$ from 0.5347 to 1.0000.\\
\textsuperscript{a} In these 2 combinations the $\ell_1$ penalty shrank every coefficient to zero in every fold and the code substituted the index of the first column; the retained feature is a fallback, not one the selector identified as informative.
Combinations with a median of one retained feature but no mark reached that median with genuinely non-zero coefficients.\\
\textsuperscript{b} The unit of analysis of the voxel-based datasets is the voxel, not the patient, and their metrics are therefore \emph{not} comparable with the remaining rows.
Across the five folds the voxel/whole datasets contribute 396085 voxel-level predictions (91500, 79327, 75748, 82369, 67141) and the voxel/centerline datasets 8557 (1690, 1571, 1811, 1806, 1679), against 141 patient-level predictions (29, 28, 28, 28, 28) for every other dataset.
The consequence is visible in the permutation null, whose standard deviation is 0.040--0.052 for the voxel-based combinations against 0.063--0.120 for all others, so that the voxel-based $p$-values are computed against a spuriously narrow null and are not comparable with the rest of the table.
3 voxel-based combinations are formally significant for AUC after correction; they must not be read as evidence for a signal.
\NOTETODO{patient-level majority-vote aggregation of the voxel-based predictions, which would make these rows comparable}
\par}

\clearpage

\begin{table}[htbp]
\centering
\caption{External validation in the MR CLEAN Registry cohort. Each pipeline was refitted on all 141 development patients with its tuned hyperparameters and applied once to the Registry cohort. CV, cross-validated median over the five outer folds of the same combination in Table~\ref{tab:cv-performance}; Reg., Registry; ACC, accuracy; $\Delta$AUC, Registry AUC minus cross-validated median AUC, a negative value being a loss of discrimination on transfer; other abbreviations as in Table~\ref{tab:cv-performance}. The voxel-based datasets were not transferred externally and therefore have no row here, which is why the table has 16 rows rather than 24.}
\label{tab:external-validation}
\smallskip
\input{\tablesdir/table4_external_validation}
\end{table}

\smallskip
{\footnotesize
\sloppy
\noindent
\textit{Notes to Table~\ref{tab:external-validation}.}
Across the 16 models the Registry AUC ranged from 0.471 to 0.521 with a median of 0.496, and 5 of them exceeded 0.50.
Registry accuracy ranged from 0.440 to 0.550 with a median of 0.481, and no model reached the majority-class accuracy of the development cohort, 0.582.
The sign of $\Delta$ follows the cross-validated value and nothing else: the 8 combinations whose cross-validated median AUC was above 0.500 all moved downwards on transfer, by $-0.034$ to $-0.190$, and the 8 that were at or below it all moved upwards, by $0.002$ to $0.128$.
The upward movement is therefore a regression to chance of models that had performed below chance internally, not a gain in generalization, and it must not be presented as one.
The only combinations for which the comparison is informative are the 6 that discriminated in cross-validation, that is reached a cross-validated median AUC of at least 0.600; all of them lose between 0.150 and 0.190 of AUC.
\NOTEWARN{$n$ Registry is not stated in the results files and is not the same for every model. Reconstructing the smallest denominator compatible with the decimal expansion of each reported accuracy leaves exactly three cohort sizes: 1649 for eight models, 1760 for six and 1761 for two. Neither the 1906 of the exclusion cascade nor the 1839 identifiers stored with the cohort is compatible with any of the sixteen. Until this is resolved the AUC values in this table are not strictly comparable with one another}
\par}

\clearpage

\begin{table}[htbp]
\centering
\caption{Comparison with the reference pipeline of van Voorst et al. \cite{vanVoorst2022} The first row summarizes the present study; rows two to four were computed on the MR CLEAN NO-IV development cohort with the same five outer folds as in Table~\ref{tab:cv-performance}, so that the comparison is made within one dataset rather than across publications; the last row is the value reported in the original publication, on a different cohort and for a different endpoint. ICC, intraclass correlation coefficient: of the 107 features the reference pipeline extracts, the 52 with ICC $\geq$ 0.75 are retained and the rest discarded. AUC$_{\text{prob}}$ denotes the AUC computed from the predicted class probabilities, and is the quantity reported for the reference pipeline. Values are medians with the interquartile range in brackets, or means $\pm$ standard deviation across folds, as indicated in the note below. An em dash marks a quantity the row does not have: the pooled row spans combinations with different accuracies, and the published row reports neither a range nor an accuracy.}
\label{tab:reference-comparison}
\smallskip
\input{\tablesdir/table5_reference_comparison}
\end{table}

\smallskip
{\footnotesize
\sloppy
\noindent
\textit{Notes to Table~\ref{tab:reference-comparison}.}
The first row pools 25 fold results, 5 combinations $\times$ 5 outer folds, and its bracketed values are the interquartile range; the second row is a mean $\pm$ standard deviation over the 5 splits; rows three and four are medians over the 5 folds, with no further feature selection applied to the 52 ICC-filtered features.
The last row is quoted from the publication: thrombus radiomics for eTICI $\geq$ 2B, a more permissive endpoint than the eTICI 2C--3 outcome of this article, as a mean over 100 bootstrap samples of 100 patients.
The plain \texttt{AUC} column of the reference pipeline degenerates to exactly 0.50 in 4 of the five splits, because the fitted model assigns the majority class to every patient in those splits and its hard predictions therefore carry no ranking information.
AUC$_{\text{prob}}$, computed from the predicted probabilities, is the only quantity of that pipeline that is comparable with the AUC values reported elsewhere in this article, and it is the one used here.
By median accuracy the reference pipeline and both of our classifiers restricted to its 52 ICC-filtered features stay below the majority-class accuracy of 0.582, although the reference pipeline exceeds it in one individual split.
Rows two to four were recomputed against the eTICI 2C--3 outcome of this article --- the stored label column is named \texttt{TICI\_2B-3}, but its per-fold positive counts are exactly the 82/141 of this study --- and are directly comparable with the first row, whereas the published 0.52 of the last row was obtained for eTICI $\geq$ 2B in a different cohort and differs from them in endpoint and cohort as much as in pipeline.
\NOTEQUERY{the pool of the first row omits one whole-thrombus combination that also discriminated in cross-validation, and the reason is not documented in the source material}
\NOTEQUERY{decide which of the two mutually incompatible sets of numbers in the source data applies before this table is finalized: set A, 0.481 $\pm$ 0.145 on 141 patients against set B, 0.467 $\pm$ 0.099 on 170}
\par}

\clearpage

\begin{table}[htbp]
\centering
\caption{Ensemble analysis: the 15 configurations of discrete, that is binary, AdaBoost \cite{Freund1997} applied to the stored prediction vectors of the 16 patient-level weak classifiers (four patient-level datasets $\times$ four pipelines). Each configuration was fitted on the five outer folds, giving 75 ensemble models. ACC $\varnothing$, mean accuracy of the configuration across the five folds; best weak classifier, the single pipeline in the configuration's pool with the highest mean accuracy, and its value; $\Delta$, the mean accuracy of the ensemble minus that of the best weak classifier. A negative $\Delta$ means the ensemble is worse than the best model it combines. The configuration name states which of the four patient-level datasets form its pool, each contributing all four pipelines: \texttt{type1} pools all four datasets (16 weak classifiers), \texttt{type2} a whole-thrombus dataset with its own centerline counterpart (8), \texttt{type3} a single dataset (4), \texttt{type4} two datasets that are not such a pair (8) and \texttt{type5} all datasets but one (12). \texttt{distrib.}\ abbreviates the distribution-based dataset and \emph{cent.} the centerline region of interest, as in Table~\ref{tab:cv-performance}. The two voxel-based datasets were excluded from every pool because their predictions are defined per voxel.}
\label{tab:ensemble}
\smallskip
\input{\tablesdir/table6_ensemble}
\end{table}

\smallskip
{\footnotesize
\sloppy
\noindent
\textit{Notes to Table~\ref{tab:ensemble}.}
$\Delta$ is negative for all 15 configurations, between $-0.022$ and $-0.121$; no ensemble outperformed the best single pipeline in its own pool.
The worst result belongs to \texttt{type1\_\allowbreak all\_\allowbreak featuresets}, the configuration that pools all feature sets.
Across all 75 fitted ensemble models accuracy averages $0.543 \pm 0.083$ with a range of 0.414--0.714.
The mechanism is documented by the fitted weights: in the all-feature-set configuration, 55 of the 80 boosting rounds (16 rounds $\times$ five folds) received a negative weight $\alpha$, meaning that the weak classifier admitted in that round performed worse than chance on the reweighted data.
The individual pipelines therefore do not make complementary errors, and the boosting assumption of weak classifiers reliably better than chance is not met.
Labels were coded as $\{-1,+1\}$, the weight of a round was $\alpha = \tfrac{1}{2}\ln\frac{1-\varepsilon}{\varepsilon}$ with $\varepsilon$ the error of the admitted classifier on the currently reweighted data, and each weak classifier could be admitted at most once, so a pool of $k$ members yields at most $k$ rounds.
The weak classifiers were not retrained: the procedure is applied to models that had already been fitted, which makes it methodologically a form of stacking rather than boosting proper.
The best weak classifier of \texttt{type3\_\allowbreak global} and of \texttt{type1\_\allowbreak all\_\allowbreak featuresets} is the same pipeline, because the pool of the latter contains the pool of the former.
\par}

\clearpage

\begin{table}[htbp]
\centering
\caption{Pairwise comparison of the pipelines by the corrected resampled $t$-test of Nadeau and Bengio \cite{Nadeau2003}. The upper block compares the four pipelines of this study with one another on the whole-thrombus global-based dataset; the lower block compares them, and our two classifiers restricted to the reference pipeline's own 52 ICC-filtered features, against the reference pipeline of van Voorst et al.\ recomputed on the same cohort. Every model in the table was evaluated on the same 5 outer folds of the same 141 patients, which is what makes the paired test applicable. Because the training sets of $k$-fold cross-validation overlap, the per-fold differences are not independent and the ordinary paired $t$-test underestimates their variance; the correction multiplies the variance by $1/k + n_{\text{test}}/n_{\text{train}}$, here $1/5 + 28.2/112.8 = 0.4500$, which inflates the standard error 1.50-fold at $\mathrm{df} = 4$. Fold sizes are taken from the data, not assumed. Mean $\Delta$AUC is the mean of the 5 per-fold differences in AUC between pipeline A and pipeline B; a positive value favours A. $t$ (uncorrected) is the ordinary paired statistic, $t$ (corrected) the same quantity with the inflated variance. The two-sided critical value at $\alpha = 0.05$ and $\mathrm{df} = 4$ is 2.776. Pipeline names combine selector and classifier as in Table~\ref{tab:cv-performance}.}
\label{tab:nadeau-bengio}
\smallskip
\input{\tablesdir/table7_nadeau_bengio}
\end{table}

\smallskip
{\footnotesize
\sloppy
\noindent
\textit{Notes to Table~\ref{tab:nadeau-bengio}.}
Not one of the 12 comparisons reaches significance: the largest corrected statistic in absolute value is 1.54, against a critical value of 2.776.
The four pipelines of this study are therefore statistically indistinguishable from one another, and none of them is distinguishable from the reference pipeline either, in either direction.
For the reference pipeline the compared quantity is AUC$_{\text{prob}}$, because its plain \texttt{AUC} column degenerates to exactly 0.50 in four of the five splits; see Table~\ref{tab:reference-comparison}.
The pooled row of Table~\ref{tab:reference-comparison} is deliberately absent here: those 5 combinations were picked by their results on these same folds, so testing them against the reference pipeline would be biased in favour of the present study. All four pipelines of the dataset are tested instead, whatever their outcome.
With 5 folds the test has $\mathrm{df} = 4$ and very little power, so these results are an absence of evidence for a difference, not evidence that the pipelines perform identically. The comparison that the data would support best, a DeLong test on the external ROC curves, was not possible because per-patient predicted scores on the Registry were not retained.
\par}

\clearpage


%%     \restoregeometry
%%
%% Cislovani S1--S7 misto 1--7 zaridi pred vlozenim:
%%     \setcounter{table}{0}\renewcommand{\thetable}{S\arabic{table}}
%%
%% Preambule potrebuje: booktabs, array, amsmath, amssymb.
%% Cesta k tabulkam se bere z \tablesdir; vychozi hodnota plati pro preklad
%% z korene projektu.
%% -------------------------------------------------------------------------
\providecommand{\tablesdir}{tables_tom/output}

%% Zaloha, kdyby `manuscriptnotes` nebyl nactveny: vysadi se jen kotva,
%% poznamka se zahodi -- tedy tentyz vysledek jako v rezimu odesilani.
\providecommand{\NOTE}[2][]{#1}
\providecommand{\NOTEWARN}[2][]{#1}
\providecommand{\NOTECHECK}[2][]{#1}
\providecommand{\NOTETODO}[2][]{#1}
\providecommand{\NOTEQUERY}[2][]{#1}

{\large\bfseries Supplementary tables\par}
\bigskip
\NOTEWARN{the captions and notes in this supplement were drafted with AI assistance and are provisional; every one of them has to be rewritten and checked by an author before submission}

\begin{table}[htbp]
\centering
\caption{Patient selection in the two cohorts. MR CLEAN NO-IV was used for model development and cross-validation, the MR CLEAN Registry as the independent external cohort. Entries are numbers of patients; an em dash marks an exclusion step that does not apply to that cohort. eTICI, expanded Thrombolysis in Cerebral Infarction score; the outcome modeled throughout this article is near-complete or complete reperfusion, defined as a final eTICI score of 2C or 3. Slice spacing refers to the through-plane spacing of the non-contrast CT acquisition.}
\label{tab:cohorts}
\smallskip
\input{\tablesdir/table1_cohorts}
\end{table}

\smallskip
{\footnotesize
\sloppy
\noindent
\textit{Notes to Table~\ref{tab:cohorts}.}
The manual review of segmentation quality was carried out on the MR CLEAN NO-IV cohort only.
No segmentation was reviewed or rejected in the MR CLEAN Registry cohort, so the two cohorts differ in segmentation quality control as well as in population.
This is a property of the design and is the single largest asymmetry between the development and the external cohort.
The review was blinded to the eTICI score: it judged the quality and the anatomical placement of the segmentation only, so the 155 exclusions it produced cannot have been informed by the outcome.
Of the 141 patients that remain, 82 (58.2\%) reached eTICI 2C--3; the proportion varies across the five outer folds (0.448, 0.714, 0.536, 0.679, 0.536), so in the first fold the majority class is the absence rather than the presence of near-complete or complete reperfusion.
\NOTETODO{breakdown of how many of the 155 reviewed cases were rejected by which criterion}
\NOTETODO{number and proportion of patients with eTICI 2C--3 reperfusion in the Registry cohort, which the table leaves as an em dash}
\NOTEWARN{size of the Registry cohort: the exclusion cascade gives 1906, the stored identifier list holds 1839, and the denominators reconstructed from the decimal expansion of the reported external accuracies are 1649 (8 models), 1760 (6 models), 1761 (2 models). Neither 1906 nor 1839 is compatible with any of the sixteen transferred models, and the reconstructed values are mutually incompatible, so the external cohort is not the same across feature sets}
\par}

\clearpage

\begin{table}[htbp]
\centering
\caption{Hyperparameter search spaces of the four pipelines. LASSO, $\ell_1$-penalized linear regression on the binary outcome used as a feature selector; RFI, ranking by the gain-based feature importance of a random-forest gradient-boosting model; XGBRF, \texttt{XGBRFClassifier}; LinearSVC, linear support vector classifier with an $\ell_1$ penalty, which acts as selector and classifier simultaneously; NuSVC, $\nu$-support vector classifier with a radial basis function kernel. In the Component column, (S) marks the component acting as feature selector, (C) as classifier, and (S+C) a single object that performs both roles. Scale indicates how Optuna sampled the parameter: \emph{log}, log-uniform over the interval; \emph{int}, uniform over the integers of the interval; \emph{float}, uniform over the interval.}
\label{tab:hyperparameters}
\smallskip
\input{\tablesdir/table2_hyperparameters}
\end{table}

\smallskip
{\footnotesize
\sloppy
\noindent
\textit{Notes to Table~\ref{tab:hyperparameters}.}
$n$ and \texttt{X.shape[1]} denote the number of rows and the number of columns of the training matrix passed to the search.
Rows are patients for the global-based and distribution-based datasets and voxels for the voxel-based datasets, so the same expression yields different intervals across datasets; the number of columns reaches 20660 for the distribution-based datasets.
The minority fraction is the proportion of the smaller outcome class among the training labels of the given fold.
In the RFI--RF pipeline a single parameter dictionary configures both the model whose gain importances rank the features and the classifier fitted on the retained ones, which is why its rows are marked (S+C) rather than (C).
All search spaces were explored with Optuna using the Gaussian-process sampler (\texttt{GPSampler}) with 20 trials per dataset-by-pipeline combination.
The objective maximized was the \emph{mean accuracy over the five inner cross-validation folds}, not AUC; the inner folds were generated by \texttt{StratifiedGroupKFold} grouped by \texttt{patient\_id}.
Settings that were fixed and not searched, one per line so that the long identifiers have somewhere to break:\\
\texttt{Lasso(fit\_intercept=True, max\_iter=5000)}\\
\texttt{LinearSVC(penalty='l1', max\_iter=5000, class\_weight='balanced')}\\
\texttt{NuSVC(kernel='rbf', probability=True, class\_weight='balanced')}\\
\texttt{XGBRFClassifier(random\_state=42)}\\
Class weighting was therefore applied in the two support vector branches only.
The models were seeded with \texttt{random\_state=42}, but \texttt{GPSampler()} was instantiated without a seed, so the sequence of 20 trials is not reproducible.
\par}

\clearpage

\begin{table}[htbp]
\centering
\caption{Cross-validated performance of all 24 dataset-by-pipeline combinations in the MR CLEAN NO-IV cohort ($n = 141$), over five outer folds. Feature set: \emph{global}, eight summary statistics of each radiomic feature over the thrombus, one row per patient; \emph{distribution}, a ten-bin histogram of each radiomic feature over the thrombus, one row per patient; \emph{voxel}, the radiomic feature values themselves, one row per voxel. ROI: \emph{whole}, the dilated thrombus mask; \emph{centerline}, the dilated skeleton of the thrombus mask. Selector (Sel.): LASSO, $\ell_1$-penalized regression; L1, embedded selection by the non-zero coefficients of the classifier itself; RFI, gain-based feature importance of an \texttt{XGBRFClassifier}. Classifier (Class.): XGBRF, \texttt{XGBRFClassifier}; LinSVC, linear support vector classifier with an $\ell_1$ penalty; NuSVC, $\nu$-support vector classifier with a radial basis function kernel. In the body of the table the distribution-based feature set is abbreviated \emph{distrib.}\ and the centerline ROI \emph{cent.} AUC, area under the receiver operating characteristic curve; ACC, accuracy; IQR, interquartile range across the five folds; min--max, the smallest and largest value across the five folds; Med.\ $n$, median number of features retained by the selector; Null AUC, mean AUC of the permutation baseline; $\kappa$, the $\kappa$-type index of the observed AUC relative to that baseline; $p$, permutation $p$-value; $q$, the same $p$-value after Benjamini--Hochberg correction across all 24 combinations.}
\label{tab:cv-performance}
\smallskip
\input{\tablesdir/table3_cv_performance}
\end{table}

\smallskip
{\footnotesize
\sloppy
\noindent
\textit{Notes to Table~\ref{tab:cv-performance}.}
Accuracy must be read against the majority-class accuracy of the development cohort, $82/141 = 0.582$; only 4 of the 24 combinations exceed it by median accuracy, and 86 of the 120 individual fold results (71.7\%) fall below it.
The class ratio itself varies across folds (0.448, 0.714, 0.536, 0.679, 0.536), so in the first fold the majority class is the absence rather than the presence of near-complete or complete reperfusion.
The largest single fold result of the whole study is 0.810, obtained by distribution/whole/RFI/NuSVC in fold 4.
Averaged over the eight whole-thrombus patient-level combinations, AUC per fold is 0.510, 0.661, 0.606, 0.587, 0.623, a spread within a single fold of up to 0.531; between-fold variability therefore exceeds between-method variability.
The permutation null was obtained from 100 label permutations at the patient level per outer fold; with $m = 100$ permutations the smallest attainable $p$-value is $p = (0+1)/(100+1) = 0.0099$, which bounds the resolution of every permutation result in this table.
The Null AUC, $\kappa$, $p$ and $q$ columns were computed from the \emph{mean} AUC across folds, whereas the AUC column reports the \emph{median}; the two summaries are not interchangeable and are shown side by side deliberately.
10 of the 24 combinations reach $p < 0.05$ uncorrected and 6 survive Benjamini--Hochberg correction\NOTETODO{cite Benjamini and Hochberg 1995; the key is missing from BibFile.bib} for AUC, while none survives it for accuracy.
No centerline combination survives correction; across the eight of them the mean AUC ranges from 0.396 to 0.495 against a null of 0.489--0.510, $\kappa$ from $-0.206$ to $-0.014$ and $p$ from 0.5347 to 1.0000.\\
\textsuperscript{a} In these 2 combinations the $\ell_1$ penalty shrank every coefficient to zero in every fold and the code substituted the index of the first column; the retained feature is a fallback, not one the selector identified as informative.
Combinations with a median of one retained feature but no mark reached that median with genuinely non-zero coefficients.\\
\textsuperscript{b} The unit of analysis of the voxel-based datasets is the voxel, not the patient, and their metrics are therefore \emph{not} comparable with the remaining rows.
Across the five folds the voxel/whole datasets contribute 396085 voxel-level predictions (91500, 79327, 75748, 82369, 67141) and the voxel/centerline datasets 8557 (1690, 1571, 1811, 1806, 1679), against 141 patient-level predictions (29, 28, 28, 28, 28) for every other dataset.
The consequence is visible in the permutation null, whose standard deviation is 0.040--0.052 for the voxel-based combinations against 0.063--0.120 for all others, so that the voxel-based $p$-values are computed against a spuriously narrow null and are not comparable with the rest of the table.
3 voxel-based combinations are formally significant for AUC after correction; they must not be read as evidence for a signal.
\NOTETODO{patient-level majority-vote aggregation of the voxel-based predictions, which would make these rows comparable}
\par}

\clearpage

\begin{table}[htbp]
\centering
\caption{External validation in the MR CLEAN Registry cohort. Each pipeline was refitted on all 141 development patients with its tuned hyperparameters and applied once to the Registry cohort. CV, cross-validated median over the five outer folds of the same combination in Table~\ref{tab:cv-performance}; Reg., Registry; ACC, accuracy; $\Delta$AUC, Registry AUC minus cross-validated median AUC, a negative value being a loss of discrimination on transfer; other abbreviations as in Table~\ref{tab:cv-performance}. The voxel-based datasets were not transferred externally and therefore have no row here, which is why the table has 16 rows rather than 24.}
\label{tab:external-validation}
\smallskip
\input{\tablesdir/table4_external_validation}
\end{table}

\smallskip
{\footnotesize
\sloppy
\noindent
\textit{Notes to Table~\ref{tab:external-validation}.}
Across the 16 models the Registry AUC ranged from 0.471 to 0.521 with a median of 0.496, and 5 of them exceeded 0.50.
Registry accuracy ranged from 0.440 to 0.550 with a median of 0.481, and no model reached the majority-class accuracy of the development cohort, 0.582.
The sign of $\Delta$ follows the cross-validated value and nothing else: the 8 combinations whose cross-validated median AUC was above 0.500 all moved downwards on transfer, by $-0.034$ to $-0.190$, and the 8 that were at or below it all moved upwards, by $0.002$ to $0.128$.
The upward movement is therefore a regression to chance of models that had performed below chance internally, not a gain in generalization, and it must not be presented as one.
The only combinations for which the comparison is informative are the 6 that discriminated in cross-validation, that is reached a cross-validated median AUC of at least 0.600; all of them lose between 0.150 and 0.190 of AUC.
\NOTEWARN{$n$ Registry is not stated in the results files and is not the same for every model. Reconstructing the smallest denominator compatible with the decimal expansion of each reported accuracy leaves exactly three cohort sizes: 1649 for eight models, 1760 for six and 1761 for two. Neither the 1906 of the exclusion cascade nor the 1839 identifiers stored with the cohort is compatible with any of the sixteen. Until this is resolved the AUC values in this table are not strictly comparable with one another}
\par}

\clearpage

\begin{table}[htbp]
\centering
\caption{Comparison with the reference pipeline of van Voorst et al. \cite{vanVoorst2022} The first row summarizes the present study; rows two to four were computed on the MR CLEAN NO-IV development cohort with the same five outer folds as in Table~\ref{tab:cv-performance}, so that the comparison is made within one dataset rather than across publications; the last row is the value reported in the original publication, on a different cohort and for a different endpoint. ICC, intraclass correlation coefficient: of the 107 features the reference pipeline extracts, the 52 with ICC $\geq$ 0.75 are retained and the rest discarded. AUC$_{\text{prob}}$ denotes the AUC computed from the predicted class probabilities, and is the quantity reported for the reference pipeline. Values are medians with the interquartile range in brackets, or means $\pm$ standard deviation across folds, as indicated in the note below. An em dash marks a quantity the row does not have: the pooled row spans combinations with different accuracies, and the published row reports neither a range nor an accuracy.}
\label{tab:reference-comparison}
\smallskip
\input{\tablesdir/table5_reference_comparison}
\end{table}

\smallskip
{\footnotesize
\sloppy
\noindent
\textit{Notes to Table~\ref{tab:reference-comparison}.}
The first row pools 25 fold results, 5 combinations $\times$ 5 outer folds, and its bracketed values are the interquartile range; the second row is a mean $\pm$ standard deviation over the 5 splits; rows three and four are medians over the 5 folds, with no further feature selection applied to the 52 ICC-filtered features.
The last row is quoted from the publication: thrombus radiomics for eTICI $\geq$ 2B, a more permissive endpoint than the eTICI 2C--3 outcome of this article, as a mean over 100 bootstrap samples of 100 patients.
The plain \texttt{AUC} column of the reference pipeline degenerates to exactly 0.50 in 4 of the five splits, because the fitted model assigns the majority class to every patient in those splits and its hard predictions therefore carry no ranking information.
AUC$_{\text{prob}}$, computed from the predicted probabilities, is the only quantity of that pipeline that is comparable with the AUC values reported elsewhere in this article, and it is the one used here.
By median accuracy the reference pipeline and both of our classifiers restricted to its 52 ICC-filtered features stay below the majority-class accuracy of 0.582, although the reference pipeline exceeds it in one individual split.
Rows two to four were recomputed against the eTICI 2C--3 outcome of this article --- the stored label column is named \texttt{TICI\_2B-3}, but its per-fold positive counts are exactly the 82/141 of this study --- and are directly comparable with the first row, whereas the published 0.52 of the last row was obtained for eTICI $\geq$ 2B in a different cohort and differs from them in endpoint and cohort as much as in pipeline.
\NOTEQUERY{the pool of the first row omits one whole-thrombus combination that also discriminated in cross-validation, and the reason is not documented in the source material}
\NOTEQUERY{decide which of the two mutually incompatible sets of numbers in the source data applies before this table is finalized: set A, 0.481 $\pm$ 0.145 on 141 patients against set B, 0.467 $\pm$ 0.099 on 170}
\par}

\clearpage

\begin{table}[htbp]
\centering
\caption{Ensemble analysis: the 15 configurations of discrete, that is binary, AdaBoost \cite{Freund1997} applied to the stored prediction vectors of the 16 patient-level weak classifiers (four patient-level datasets $\times$ four pipelines). Each configuration was fitted on the five outer folds, giving 75 ensemble models. ACC $\varnothing$, mean accuracy of the configuration across the five folds; best weak classifier, the single pipeline in the configuration's pool with the highest mean accuracy, and its value; $\Delta$, the mean accuracy of the ensemble minus that of the best weak classifier. A negative $\Delta$ means the ensemble is worse than the best model it combines. The configuration name states which of the four patient-level datasets form its pool, each contributing all four pipelines: \texttt{type1} pools all four datasets (16 weak classifiers), \texttt{type2} a whole-thrombus dataset with its own centerline counterpart (8), \texttt{type3} a single dataset (4), \texttt{type4} two datasets that are not such a pair (8) and \texttt{type5} all datasets but one (12). \texttt{distrib.}\ abbreviates the distribution-based dataset and \emph{cent.} the centerline region of interest, as in Table~\ref{tab:cv-performance}. The two voxel-based datasets were excluded from every pool because their predictions are defined per voxel.}
\label{tab:ensemble}
\smallskip
\input{\tablesdir/table6_ensemble}
\end{table}

\smallskip
{\footnotesize
\sloppy
\noindent
\textit{Notes to Table~\ref{tab:ensemble}.}
$\Delta$ is negative for all 15 configurations, between $-0.022$ and $-0.121$; no ensemble outperformed the best single pipeline in its own pool.
The worst result belongs to \texttt{type1\_\allowbreak all\_\allowbreak featuresets}, the configuration that pools all feature sets.
Across all 75 fitted ensemble models accuracy averages $0.543 \pm 0.083$ with a range of 0.414--0.714.
The mechanism is documented by the fitted weights: in the all-feature-set configuration, 55 of the 80 boosting rounds (16 rounds $\times$ five folds) received a negative weight $\alpha$, meaning that the weak classifier admitted in that round performed worse than chance on the reweighted data.
The individual pipelines therefore do not make complementary errors, and the boosting assumption of weak classifiers reliably better than chance is not met.
Labels were coded as $\{-1,+1\}$, the weight of a round was $\alpha = \tfrac{1}{2}\ln\frac{1-\varepsilon}{\varepsilon}$ with $\varepsilon$ the error of the admitted classifier on the currently reweighted data, and each weak classifier could be admitted at most once, so a pool of $k$ members yields at most $k$ rounds.
The weak classifiers were not retrained: the procedure is applied to models that had already been fitted, which makes it methodologically a form of stacking rather than boosting proper.
The best weak classifier of \texttt{type3\_\allowbreak global} and of \texttt{type1\_\allowbreak all\_\allowbreak featuresets} is the same pipeline, because the pool of the latter contains the pool of the former.
\par}

\clearpage

\begin{table}[htbp]
\centering
\caption{Pairwise comparison of the pipelines by the corrected resampled $t$-test of Nadeau and Bengio \cite{Nadeau2003}. The upper block compares the four pipelines of this study with one another on the whole-thrombus global-based dataset; the lower block compares them, and our two classifiers restricted to the reference pipeline's own 52 ICC-filtered features, against the reference pipeline of van Voorst et al.\ recomputed on the same cohort. Every model in the table was evaluated on the same 5 outer folds of the same 141 patients, which is what makes the paired test applicable. Because the training sets of $k$-fold cross-validation overlap, the per-fold differences are not independent and the ordinary paired $t$-test underestimates their variance; the correction multiplies the variance by $1/k + n_{\text{test}}/n_{\text{train}}$, here $1/5 + 28.2/112.8 = 0.4500$, which inflates the standard error 1.50-fold at $\mathrm{df} = 4$. Fold sizes are taken from the data, not assumed. Mean $\Delta$AUC is the mean of the 5 per-fold differences in AUC between pipeline A and pipeline B; a positive value favours A. $t$ (uncorrected) is the ordinary paired statistic, $t$ (corrected) the same quantity with the inflated variance. The two-sided critical value at $\alpha = 0.05$ and $\mathrm{df} = 4$ is 2.776. Pipeline names combine selector and classifier as in Table~\ref{tab:cv-performance}.}
\label{tab:nadeau-bengio}
\smallskip
\input{\tablesdir/table7_nadeau_bengio}
\end{table}

\smallskip
{\footnotesize
\sloppy
\noindent
\textit{Notes to Table~\ref{tab:nadeau-bengio}.}
Not one of the 12 comparisons reaches significance: the largest corrected statistic in absolute value is 1.54, against a critical value of 2.776.
The four pipelines of this study are therefore statistically indistinguishable from one another, and none of them is distinguishable from the reference pipeline either, in either direction.
For the reference pipeline the compared quantity is AUC$_{\text{prob}}$, because its plain \texttt{AUC} column degenerates to exactly 0.50 in four of the five splits; see Table~\ref{tab:reference-comparison}.
The pooled row of Table~\ref{tab:reference-comparison} is deliberately absent here: those 5 combinations were picked by their results on these same folds, so testing them against the reference pipeline would be biased in favour of the present study. All four pipelines of the dataset are tested instead, whatever their outcome.
With 5 folds the test has $\mathrm{df} = 4$ and very little power, so these results are an absence of evidence for a difference, not evidence that the pipelines perform identically. The comparison that the data would support best, a DeLong test on the external ROC curves, was not possible because per-patient predicted scores on the Registry were not retained.
\par}

\clearpage


%%     \restoregeometry
%%
%% Cislovani S1--S7 misto 1--7 zaridi pred vlozenim:
%%     \setcounter{table}{0}\renewcommand{\thetable}{S\arabic{table}}
%%
%% Preambule potrebuje: booktabs, array, amsmath, amssymb.
%% Cesta k tabulkam se bere z \tablesdir; vychozi hodnota plati pro preklad
%% z korene projektu.
%% -------------------------------------------------------------------------
\providecommand{\tablesdir}{tables_tom/output}

%% Zaloha, kdyby `manuscriptnotes` nebyl nactveny: vysadi se jen kotva,
%% poznamka se zahodi -- tedy tentyz vysledek jako v rezimu odesilani.
\providecommand{\NOTE}[2][]{#1}
\providecommand{\NOTEWARN}[2][]{#1}
\providecommand{\NOTECHECK}[2][]{#1}
\providecommand{\NOTETODO}[2][]{#1}
\providecommand{\NOTEQUERY}[2][]{#1}

{\large\bfseries Supplementary tables\par}
\bigskip
\NOTEWARN{the captions and notes in this supplement were drafted with AI assistance and are provisional; every one of them has to be rewritten and checked by an author before submission}

\begin{table}[htbp]
\centering
\caption{Patient selection in the two cohorts. MR CLEAN NO-IV was used for model development and cross-validation, the MR CLEAN Registry as the independent external cohort. Entries are numbers of patients; an em dash marks an exclusion step that does not apply to that cohort. eTICI, expanded Thrombolysis in Cerebral Infarction score; the outcome modeled throughout this article is near-complete or complete reperfusion, defined as a final eTICI score of 2C or 3. Slice spacing refers to the through-plane spacing of the non-contrast CT acquisition.}
\label{tab:cohorts}
\smallskip
\input{\tablesdir/table1_cohorts}
\end{table}

\smallskip
{\footnotesize
\sloppy
\noindent
\textit{Notes to Table~\ref{tab:cohorts}.}
The manual review of segmentation quality was carried out on the MR CLEAN NO-IV cohort only.
No segmentation was reviewed or rejected in the MR CLEAN Registry cohort, so the two cohorts differ in segmentation quality control as well as in population.
This is a property of the design and is the single largest asymmetry between the development and the external cohort.
The review was blinded to the eTICI score: it judged the quality and the anatomical placement of the segmentation only, so the 155 exclusions it produced cannot have been informed by the outcome.
Of the 141 patients that remain, 82 (58.2\%) reached eTICI 2C--3; the proportion varies across the five outer folds (0.448, 0.714, 0.536, 0.679, 0.536), so in the first fold the majority class is the absence rather than the presence of near-complete or complete reperfusion.
\NOTETODO{breakdown of how many of the 155 reviewed cases were rejected by which criterion}
\NOTETODO{number and proportion of patients with eTICI 2C--3 reperfusion in the Registry cohort, which the table leaves as an em dash}
\NOTEWARN{size of the Registry cohort: the exclusion cascade gives 1906, the stored identifier list holds 1839, and the denominators reconstructed from the decimal expansion of the reported external accuracies are 1649 (8 models), 1760 (6 models), 1761 (2 models). Neither 1906 nor 1839 is compatible with any of the sixteen transferred models, and the reconstructed values are mutually incompatible, so the external cohort is not the same across feature sets}
\par}

\clearpage

\begin{table}[htbp]
\centering
\caption{Hyperparameter search spaces of the four pipelines. LASSO, $\ell_1$-penalized linear regression on the binary outcome used as a feature selector; RFI, ranking by the gain-based feature importance of a random-forest gradient-boosting model; XGBRF, \texttt{XGBRFClassifier}; LinearSVC, linear support vector classifier with an $\ell_1$ penalty, which acts as selector and classifier simultaneously; NuSVC, $\nu$-support vector classifier with a radial basis function kernel. In the Component column, (S) marks the component acting as feature selector, (C) as classifier, and (S+C) a single object that performs both roles. Scale indicates how Optuna sampled the parameter: \emph{log}, log-uniform over the interval; \emph{int}, uniform over the integers of the interval; \emph{float}, uniform over the interval.}
\label{tab:hyperparameters}
\smallskip
\input{\tablesdir/table2_hyperparameters}
\end{table}

\smallskip
{\footnotesize
\sloppy
\noindent
\textit{Notes to Table~\ref{tab:hyperparameters}.}
$n$ and \texttt{X.shape[1]} denote the number of rows and the number of columns of the training matrix passed to the search.
Rows are patients for the global-based and distribution-based datasets and voxels for the voxel-based datasets, so the same expression yields different intervals across datasets; the number of columns reaches 20660 for the distribution-based datasets.
The minority fraction is the proportion of the smaller outcome class among the training labels of the given fold.
In the RFI--RF pipeline a single parameter dictionary configures both the model whose gain importances rank the features and the classifier fitted on the retained ones, which is why its rows are marked (S+C) rather than (C).
All search spaces were explored with Optuna using the Gaussian-process sampler (\texttt{GPSampler}) with 20 trials per dataset-by-pipeline combination.
The objective maximized was the \emph{mean accuracy over the five inner cross-validation folds}, not AUC; the inner folds were generated by \texttt{StratifiedGroupKFold} grouped by \texttt{patient\_id}.
Settings that were fixed and not searched, one per line so that the long identifiers have somewhere to break:\\
\texttt{Lasso(fit\_intercept=True, max\_iter=5000)}\\
\texttt{LinearSVC(penalty='l1', max\_iter=5000, class\_weight='balanced')}\\
\texttt{NuSVC(kernel='rbf', probability=True, class\_weight='balanced')}\\
\texttt{XGBRFClassifier(random\_state=42)}\\
Class weighting was therefore applied in the two support vector branches only.
The models were seeded with \texttt{random\_state=42}, but \texttt{GPSampler()} was instantiated without a seed, so the sequence of 20 trials is not reproducible.
\par}

\clearpage

\begin{table}[htbp]
\centering
\caption{Cross-validated performance of all 24 dataset-by-pipeline combinations in the MR CLEAN NO-IV cohort ($n = 141$), over five outer folds. Feature set: \emph{global}, eight summary statistics of each radiomic feature over the thrombus, one row per patient; \emph{distribution}, a ten-bin histogram of each radiomic feature over the thrombus, one row per patient; \emph{voxel}, the radiomic feature values themselves, one row per voxel. ROI: \emph{whole}, the dilated thrombus mask; \emph{centerline}, the dilated skeleton of the thrombus mask. Selector (Sel.): LASSO, $\ell_1$-penalized regression; L1, embedded selection by the non-zero coefficients of the classifier itself; RFI, gain-based feature importance of an \texttt{XGBRFClassifier}. Classifier (Class.): XGBRF, \texttt{XGBRFClassifier}; LinSVC, linear support vector classifier with an $\ell_1$ penalty; NuSVC, $\nu$-support vector classifier with a radial basis function kernel. In the body of the table the distribution-based feature set is abbreviated \emph{distrib.}\ and the centerline ROI \emph{cent.} AUC, area under the receiver operating characteristic curve; ACC, accuracy; IQR, interquartile range across the five folds; min--max, the smallest and largest value across the five folds; Med.\ $n$, median number of features retained by the selector; Null AUC, mean AUC of the permutation baseline; $\kappa$, the $\kappa$-type index of the observed AUC relative to that baseline; $p$, permutation $p$-value; $q$, the same $p$-value after Benjamini--Hochberg correction across all 24 combinations.}
\label{tab:cv-performance}
\smallskip
\input{\tablesdir/table3_cv_performance}
\end{table}

\smallskip
{\footnotesize
\sloppy
\noindent
\textit{Notes to Table~\ref{tab:cv-performance}.}
Accuracy must be read against the majority-class accuracy of the development cohort, $82/141 = 0.582$; only 4 of the 24 combinations exceed it by median accuracy, and 86 of the 120 individual fold results (71.7\%) fall below it.
The class ratio itself varies across folds (0.448, 0.714, 0.536, 0.679, 0.536), so in the first fold the majority class is the absence rather than the presence of near-complete or complete reperfusion.
The largest single fold result of the whole study is 0.810, obtained by distribution/whole/RFI/NuSVC in fold 4.
Averaged over the eight whole-thrombus patient-level combinations, AUC per fold is 0.510, 0.661, 0.606, 0.587, 0.623, a spread within a single fold of up to 0.531; between-fold variability therefore exceeds between-method variability.
The permutation null was obtained from 100 label permutations at the patient level per outer fold; with $m = 100$ permutations the smallest attainable $p$-value is $p = (0+1)/(100+1) = 0.0099$, which bounds the resolution of every permutation result in this table.
The Null AUC, $\kappa$, $p$ and $q$ columns were computed from the \emph{mean} AUC across folds, whereas the AUC column reports the \emph{median}; the two summaries are not interchangeable and are shown side by side deliberately.
10 of the 24 combinations reach $p < 0.05$ uncorrected and 6 survive Benjamini--Hochberg correction\NOTETODO{cite Benjamini and Hochberg 1995; the key is missing from BibFile.bib} for AUC, while none survives it for accuracy.
No centerline combination survives correction; across the eight of them the mean AUC ranges from 0.396 to 0.495 against a null of 0.489--0.510, $\kappa$ from $-0.206$ to $-0.014$ and $p$ from 0.5347 to 1.0000.\\
\textsuperscript{a} In these 2 combinations the $\ell_1$ penalty shrank every coefficient to zero in every fold and the code substituted the index of the first column; the retained feature is a fallback, not one the selector identified as informative.
Combinations with a median of one retained feature but no mark reached that median with genuinely non-zero coefficients.\\
\textsuperscript{b} The unit of analysis of the voxel-based datasets is the voxel, not the patient, and their metrics are therefore \emph{not} comparable with the remaining rows.
Across the five folds the voxel/whole datasets contribute 396085 voxel-level predictions (91500, 79327, 75748, 82369, 67141) and the voxel/centerline datasets 8557 (1690, 1571, 1811, 1806, 1679), against 141 patient-level predictions (29, 28, 28, 28, 28) for every other dataset.
The consequence is visible in the permutation null, whose standard deviation is 0.040--0.052 for the voxel-based combinations against 0.063--0.120 for all others, so that the voxel-based $p$-values are computed against a spuriously narrow null and are not comparable with the rest of the table.
3 voxel-based combinations are formally significant for AUC after correction; they must not be read as evidence for a signal.
\NOTETODO{patient-level majority-vote aggregation of the voxel-based predictions, which would make these rows comparable}
\par}

\clearpage

\begin{table}[htbp]
\centering
\caption{External validation in the MR CLEAN Registry cohort. Each pipeline was refitted on all 141 development patients with its tuned hyperparameters and applied once to the Registry cohort. CV, cross-validated median over the five outer folds of the same combination in Table~\ref{tab:cv-performance}; Reg., Registry; ACC, accuracy; $\Delta$AUC, Registry AUC minus cross-validated median AUC, a negative value being a loss of discrimination on transfer; other abbreviations as in Table~\ref{tab:cv-performance}. The voxel-based datasets were not transferred externally and therefore have no row here, which is why the table has 16 rows rather than 24.}
\label{tab:external-validation}
\smallskip
\input{\tablesdir/table4_external_validation}
\end{table}

\smallskip
{\footnotesize
\sloppy
\noindent
\textit{Notes to Table~\ref{tab:external-validation}.}
Across the 16 models the Registry AUC ranged from 0.471 to 0.521 with a median of 0.496, and 5 of them exceeded 0.50.
Registry accuracy ranged from 0.440 to 0.550 with a median of 0.481, and no model reached the majority-class accuracy of the development cohort, 0.582.
The sign of $\Delta$ follows the cross-validated value and nothing else: the 8 combinations whose cross-validated median AUC was above 0.500 all moved downwards on transfer, by $-0.034$ to $-0.190$, and the 8 that were at or below it all moved upwards, by $0.002$ to $0.128$.
The upward movement is therefore a regression to chance of models that had performed below chance internally, not a gain in generalization, and it must not be presented as one.
The only combinations for which the comparison is informative are the 6 that discriminated in cross-validation, that is reached a cross-validated median AUC of at least 0.600; all of them lose between 0.150 and 0.190 of AUC.
\NOTEWARN{$n$ Registry is not stated in the results files and is not the same for every model. Reconstructing the smallest denominator compatible with the decimal expansion of each reported accuracy leaves exactly three cohort sizes: 1649 for eight models, 1760 for six and 1761 for two. Neither the 1906 of the exclusion cascade nor the 1839 identifiers stored with the cohort is compatible with any of the sixteen. Until this is resolved the AUC values in this table are not strictly comparable with one another}
\par}

\clearpage

\begin{table}[htbp]
\centering
\caption{Comparison with the reference pipeline of van Voorst et al. \cite{vanVoorst2022} The first row summarizes the present study; rows two to four were computed on the MR CLEAN NO-IV development cohort with the same five outer folds as in Table~\ref{tab:cv-performance}, so that the comparison is made within one dataset rather than across publications; the last row is the value reported in the original publication, on a different cohort and for a different endpoint. ICC, intraclass correlation coefficient: of the 107 features the reference pipeline extracts, the 52 with ICC $\geq$ 0.75 are retained and the rest discarded. AUC$_{\text{prob}}$ denotes the AUC computed from the predicted class probabilities, and is the quantity reported for the reference pipeline. Values are medians with the interquartile range in brackets, or means $\pm$ standard deviation across folds, as indicated in the note below. An em dash marks a quantity the row does not have: the pooled row spans combinations with different accuracies, and the published row reports neither a range nor an accuracy.}
\label{tab:reference-comparison}
\smallskip
\input{\tablesdir/table5_reference_comparison}
\end{table}

\smallskip
{\footnotesize
\sloppy
\noindent
\textit{Notes to Table~\ref{tab:reference-comparison}.}
The first row pools 25 fold results, 5 combinations $\times$ 5 outer folds, and its bracketed values are the interquartile range; the second row is a mean $\pm$ standard deviation over the 5 splits; rows three and four are medians over the 5 folds, with no further feature selection applied to the 52 ICC-filtered features.
The last row is quoted from the publication: thrombus radiomics for eTICI $\geq$ 2B, a more permissive endpoint than the eTICI 2C--3 outcome of this article, as a mean over 100 bootstrap samples of 100 patients.
The plain \texttt{AUC} column of the reference pipeline degenerates to exactly 0.50 in 4 of the five splits, because the fitted model assigns the majority class to every patient in those splits and its hard predictions therefore carry no ranking information.
AUC$_{\text{prob}}$, computed from the predicted probabilities, is the only quantity of that pipeline that is comparable with the AUC values reported elsewhere in this article, and it is the one used here.
By median accuracy the reference pipeline and both of our classifiers restricted to its 52 ICC-filtered features stay below the majority-class accuracy of 0.582, although the reference pipeline exceeds it in one individual split.
Rows two to four were recomputed against the eTICI 2C--3 outcome of this article --- the stored label column is named \texttt{TICI\_2B-3}, but its per-fold positive counts are exactly the 82/141 of this study --- and are directly comparable with the first row, whereas the published 0.52 of the last row was obtained for eTICI $\geq$ 2B in a different cohort and differs from them in endpoint and cohort as much as in pipeline.
\NOTEQUERY{the pool of the first row omits one whole-thrombus combination that also discriminated in cross-validation, and the reason is not documented in the source material}
\NOTEQUERY{decide which of the two mutually incompatible sets of numbers in the source data applies before this table is finalized: set A, 0.481 $\pm$ 0.145 on 141 patients against set B, 0.467 $\pm$ 0.099 on 170}
\par}

\clearpage

\begin{table}[htbp]
\centering
\caption{Ensemble analysis: the 15 configurations of discrete, that is binary, AdaBoost \cite{Freund1997} applied to the stored prediction vectors of the 16 patient-level weak classifiers (four patient-level datasets $\times$ four pipelines). Each configuration was fitted on the five outer folds, giving 75 ensemble models. ACC $\varnothing$, mean accuracy of the configuration across the five folds; best weak classifier, the single pipeline in the configuration's pool with the highest mean accuracy, and its value; $\Delta$, the mean accuracy of the ensemble minus that of the best weak classifier. A negative $\Delta$ means the ensemble is worse than the best model it combines. The configuration name states which of the four patient-level datasets form its pool, each contributing all four pipelines: \texttt{type1} pools all four datasets (16 weak classifiers), \texttt{type2} a whole-thrombus dataset with its own centerline counterpart (8), \texttt{type3} a single dataset (4), \texttt{type4} two datasets that are not such a pair (8) and \texttt{type5} all datasets but one (12). \texttt{distrib.}\ abbreviates the distribution-based dataset and \emph{cent.} the centerline region of interest, as in Table~\ref{tab:cv-performance}. The two voxel-based datasets were excluded from every pool because their predictions are defined per voxel.}
\label{tab:ensemble}
\smallskip
\input{\tablesdir/table6_ensemble}
\end{table}

\smallskip
{\footnotesize
\sloppy
\noindent
\textit{Notes to Table~\ref{tab:ensemble}.}
$\Delta$ is negative for all 15 configurations, between $-0.022$ and $-0.121$; no ensemble outperformed the best single pipeline in its own pool.
The worst result belongs to \texttt{type1\_\allowbreak all\_\allowbreak featuresets}, the configuration that pools all feature sets.
Across all 75 fitted ensemble models accuracy averages $0.543 \pm 0.083$ with a range of 0.414--0.714.
The mechanism is documented by the fitted weights: in the all-feature-set configuration, 55 of the 80 boosting rounds (16 rounds $\times$ five folds) received a negative weight $\alpha$, meaning that the weak classifier admitted in that round performed worse than chance on the reweighted data.
The individual pipelines therefore do not make complementary errors, and the boosting assumption of weak classifiers reliably better than chance is not met.
Labels were coded as $\{-1,+1\}$, the weight of a round was $\alpha = \tfrac{1}{2}\ln\frac{1-\varepsilon}{\varepsilon}$ with $\varepsilon$ the error of the admitted classifier on the currently reweighted data, and each weak classifier could be admitted at most once, so a pool of $k$ members yields at most $k$ rounds.
The weak classifiers were not retrained: the procedure is applied to models that had already been fitted, which makes it methodologically a form of stacking rather than boosting proper.
The best weak classifier of \texttt{type3\_\allowbreak global} and of \texttt{type1\_\allowbreak all\_\allowbreak featuresets} is the same pipeline, because the pool of the latter contains the pool of the former.
\par}

\clearpage

\begin{table}[htbp]
\centering
\caption{Pairwise comparison of the pipelines by the corrected resampled $t$-test of Nadeau and Bengio \cite{Nadeau2003}. The upper block compares the four pipelines of this study with one another on the whole-thrombus global-based dataset; the lower block compares them, and our two classifiers restricted to the reference pipeline's own 52 ICC-filtered features, against the reference pipeline of van Voorst et al.\ recomputed on the same cohort. Every model in the table was evaluated on the same 5 outer folds of the same 141 patients, which is what makes the paired test applicable. Because the training sets of $k$-fold cross-validation overlap, the per-fold differences are not independent and the ordinary paired $t$-test underestimates their variance; the correction multiplies the variance by $1/k + n_{\text{test}}/n_{\text{train}}$, here $1/5 + 28.2/112.8 = 0.4500$, which inflates the standard error 1.50-fold at $\mathrm{df} = 4$. Fold sizes are taken from the data, not assumed. Mean $\Delta$AUC is the mean of the 5 per-fold differences in AUC between pipeline A and pipeline B; a positive value favours A. $t$ (uncorrected) is the ordinary paired statistic, $t$ (corrected) the same quantity with the inflated variance. The two-sided critical value at $\alpha = 0.05$ and $\mathrm{df} = 4$ is 2.776. Pipeline names combine selector and classifier as in Table~\ref{tab:cv-performance}.}
\label{tab:nadeau-bengio}
\smallskip
\input{\tablesdir/table7_nadeau_bengio}
\end{table}

\smallskip
{\footnotesize
\sloppy
\noindent
\textit{Notes to Table~\ref{tab:nadeau-bengio}.}
Not one of the 12 comparisons reaches significance: the largest corrected statistic in absolute value is 1.54, against a critical value of 2.776.
The four pipelines of this study are therefore statistically indistinguishable from one another, and none of them is distinguishable from the reference pipeline either, in either direction.
For the reference pipeline the compared quantity is AUC$_{\text{prob}}$, because its plain \texttt{AUC} column degenerates to exactly 0.50 in four of the five splits; see Table~\ref{tab:reference-comparison}.
The pooled row of Table~\ref{tab:reference-comparison} is deliberately absent here: those 5 combinations were picked by their results on these same folds, so testing them against the reference pipeline would be biased in favour of the present study. All four pipelines of the dataset are tested instead, whatever their outcome.
With 5 folds the test has $\mathrm{df} = 4$ and very little power, so these results are an absence of evidence for a difference, not evidence that the pipelines perform identically. The comparison that the data would support best, a DeLong test on the external ROC curves, was not possible because per-patient predicted scores on the Registry were not retained.
\par}

\clearpage

